% Section 2 completed with formal algorithm description and breakable pseudocode, 2026-10-11.
\documentclass[11pt,reqno]{amsart}
\usepackage[tt=false]{libertine}
\usepackage{amssymb}
\usepackage{mathtools}
\usepackage[varbb]{newpxmath}
\usepackage[foot]{amsaddr}


\let\savedbigtimes\bigtimes
\let\bigtimes\relax
\usepackage{mathabx}
\let\bigtimes\savedbigtimes

\usepackage[margin=1in, bottom=1in]{geometry}
\usepackage{graphicx}
\usepackage{caption}
\usepackage{subcaption}
\usepackage{enumitem}
\usepackage{algorithm}
\usepackage{algpseudocode}
\usepackage{comment}
\usepackage{bbm}

\usepackage[dvipsnames]{xcolor}
\DeclareRobustCommand{\rev}[1]{{\color{blue}\hypersetup{linkcolor=blue,citecolor=blue,urlcolor=blue}#1}}
\usepackage[colorlinks=true,
  linkcolor=teal!60!black,
  citecolor=PineGreen,
  urlcolor=RedViolet]{hyperref}
\hypersetup{bookmarksopen=true}
\pdfstringdefDisableCommands{\def\rev#1{#1}}

\usepackage[yyyymmdd,hhmmss]{datetime}

\usepackage{tikz}
\usetikzlibrary{decorations.pathreplacing,calligraphy}
\usetikzlibrary{calc}

\linespread{1.1}

\raggedbottom
\setlength{\emergencystretch}{2em}

% define theorem environments, commands, etc here
% theorem environments
\newtheorem{thm}{Theorem}[section]
\newtheorem{lem}[thm]{Lemma}
\newtheorem{ppn}[thm]{Proposition}
\newtheorem{prop}[thm]{Proposition}
\newtheorem{cor}[thm]{Corollary}
\newtheorem{fac}[thm]{Fact}
\newtheorem{conj}[thm]{Conjecture}

\theoremstyle{definition}
\newtheorem{dfn}[thm]{Definition}
\newtheorem{ex}[thm]{Example}
\newtheorem{asn}[thm]{Assumption}

\theoremstyle{remark}
\newtheorem{rmk}[thm]{Remark}

% equation environments
\numberwithin{equation}{section}
\def\beq#1\eeq{%
    \begin{equation}%
    #1%
    \end{equation}%
}
\def\baln#1\ealn{%
    \begin{align}%
    #1%
    \end{align}%
}
\def\balnn#1\ealnn{%
    \begin{align*}%
    #1%
    \end{align*}%
}

% macros
\def\lt{\left}
\def\rt{\right}
\def\fr{\frac}
\def\la{\langle}
\def\ra{\rangle}
\def\eps{\varepsilon}

\def\bA{{\boldsymbol{A}}}
\def\bB{{\boldsymbol{B}}}
\def\bC{{\boldsymbol{C}}}
\def\bD{{\boldsymbol{D}}}
\def\bE{{\boldsymbol{E}}}
\def\bF{{\boldsymbol{F}}}
\def\bG{{\boldsymbol{G}}}
\def\bH{{\boldsymbol{H}}}
\def\bI{{\boldsymbol{I}}}
\def\bJ{{\boldsymbol{J}}}
\def\bK{{\boldsymbol{K}}}
\def\bL{{\boldsymbol{L}}}
\def\bM{{\boldsymbol{M}}}
\def\bN{{\boldsymbol{N}}}
\def\bO{{\boldsymbol{O}}}
\def\bP{{\boldsymbol{P}}}
\def\bQ{{\boldsymbol{Q}}}
\def\bR{{\boldsymbol{R}}}
\def\bS{{\boldsymbol{S}}}
\def\bT{{\boldsymbol{T}}}
\def\bU{{\boldsymbol{U}}}
\def\bV{{\boldsymbol{V}}}
\def\bW{{\boldsymbol{W}}}
\def\bX{{\boldsymbol{X}}}
\def\bY{{\boldsymbol{Y}}}
\def\bZ{{\boldsymbol{Z}}}
\def\bDelta{{\boldsymbol{\Delta}}}
\def\bLambda{{\boldsymbol{\Lambda}}}
\def\bxi{{\boldsymbol{\xi}}}
\def\bomega{{\boldsymbol{\omega}}}
\def\btheta{{\boldsymbol{\theta}}}

\def\ba{{\boldsymbol{a}}}
\def\bb{{\boldsymbol{b}}}
\def\bc{{\boldsymbol{c}}}
\def\bd{{\boldsymbol{d}}}
\def\be{{\boldsymbol{e}}}
\def\bf{{\boldsymbol{f}}}
\def\bg{{\boldsymbol{g}}}
\def\bh{{\boldsymbol{h}}}
\def\bi{{\boldsymbol{i}}}
\def\bj{{\boldsymbol{j}}}
\def\bk{{\boldsymbol{k}}}
\def\bl{{\boldsymbol{\ell}}}
\def\bm{{\boldsymbol{m}}}
\def\bn{{\boldsymbol{n}}}
\def\bo{{\boldsymbol{o}}}
\def\bp{{\boldsymbol{p}}}
\def\bq{{\boldsymbol{q}}}
\def\br{{\boldsymbol{r}}}
\def\bs{{\boldsymbol{s}}}
\def\bt{{\boldsymbol{t}}}
\def\bu{{\boldsymbol{u}}}
\def\bv{{\boldsymbol{v}}}
\def\bw{{\boldsymbol{w}}}
\def\bx{{\boldsymbol{x}}}
\def\by{{\boldsymbol{y}}}
\def\bz{{\boldsymbol{z}}}
\def\bone{{\boldsymbol{1}}}

\def\bbA{{\mathbb{A}}}
\def\bbB{{\mathbb{B}}}
\def\bbC{{\mathbb{C}}}
\def\bbD{{\mathbb{D}}}
\def\bbE{{\mathbb{E}}}
\def\bbF{{\mathbb{F}}}
\def\bbG{{\mathbb{G}}}
\def\bbH{{\mathbb{H}}}
\def\bbI{{\mathbb{I}}}
\def\bbJ{{\mathbb{J}}}
\def\bbK{{\mathbb{K}}}
\def\bbL{{\mathbb{L}}}
\def\bbM{{\mathbb{M}}}
\def\bbN{{\mathbb{N}}}
\def\bbO{{\mathbb{O}}}
\def\bbP{{\mathbb{P}}}
\def\bbQ{{\mathbb{Q}}}
\def\bbR{{\mathbb{R}}}
\def\bbS{{\mathbb{S}}}
\def\bbT{{\mathbb{T}}}
\def\bbU{{\mathbb{U}}}
\def\bbV{{\mathbb{V}}}
\def\bbW{{\mathbb{W}}}
\def\bbX{{\mathbb{X}}}
\def\bbY{{\mathbb{Y}}}
\def\bbZ{{\mathbb{Z}}}
\def\bbone{{\mathbb{1}}}

\def\cA{{\mathcal{A}}}
\def\cB{{\mathcal{B}}}
\def\cC{{\mathcal{C}}}
\def\cD{{\mathcal{D}}}
\def\cE{{\mathcal{E}}}
\def\cF{{\mathcal{F}}}
\def\cG{{\mathcal{G}}}
\def\cH{{\mathcal{H}}}
\def\cI{{\mathcal{I}}}
\def\cJ{{\mathcal{J}}}
\def\cK{{\mathcal{K}}}
\def\cL{{\mathcal{L}}}
\def\cM{{\mathcal{M}}}
\def\cN{{\mathcal{N}}}
\def\cO{{\mathcal{O}}}
\def\cP{{\mathcal{P}}}
\def\cQ{{\mathcal{Q}}}
\def\cR{{\mathcal{R}}}
\def\cS{{\mathcal{S}}}
\def\cT{{\mathcal{T}}}
\def\cU{{\mathcal{U}}}
\def\cV{{\mathcal{V}}}
\def\cW{{\mathcal{W}}}
\def\cX{{\mathcal{X}}}
\def\cY{{\mathcal{Y}}}
\def\cZ{{\mathcal{Z}}}

\def\sA{{\mathscr{A}}}
\def\sB{{\mathscr{B}}}
\def\sC{{\mathscr{C}}}
\def\sD{{\mathscr{D}}}
\def\sE{{\mathscr{E}}}
\def\sF{{\mathscr{F}}}
\def\sG{{\mathscr{G}}}
\def\sH{{\mathscr{H}}}
\def\sI{{\mathscr{I}}}
\def\sJ{{\mathscr{J}}}
\def\sK{{\mathscr{K}}}
\def\sL{{\mathscr{L}}}
\def\sM{{\mathscr{M}}}
\def\sN{{\mathscr{N}}}
\def\sO{{\mathscr{O}}}
\def\sP{{\mathscr{P}}}
\def\sQ{{\mathscr{Q}}}
\def\sR{{\mathscr{R}}}
\def\sS{{\mathscr{S}}}
\def\sT{{\mathscr{T}}}
\def\sU{{\mathscr{U}}}
\def\sV{{\mathscr{V}}}
\def\sW{{\mathscr{W}}}
\def\sX{{\mathscr{X}}}
\def\sY{{\mathscr{Y}}}
\def\sZ{{\mathscr{Z}}}

\DeclareMathOperator{\EE}{{\bbE}}
\DeclareMathOperator{\PP}{{\bbP}}


% for comments
\newcommand{\TODO}[1]{\noindent\textup{\textsf{\color{red}[Todo: #1]}}}

% macros
\def\de{{\mathsf{d}}}
\def\tap{{\mathsf{TAP}}}
\def\pro{{\mathsf{pro}}}
\newcommand{\TV}{\operatorname{TV}}
\newcommand{\KL}{\operatorname{KL}}
\newcommand{\Cov}{\operatorname{Cov}}
\newcommand{\Var}{\operatorname{Var}}
\newcommand{\Tr}{\operatorname{Tr}}
\newcommand{\op}{\operatorname{op}}
\newcommand{\Fop}{\operatorname{F}}
\newcommand{\sech}{\operatorname{sech}}
\newcommand{\atanh}{\operatorname{arctanh}}
\newcommand{\sign}{\operatorname{sign}}
\newcommand{\bzero}{\mathbf 0}

\makeatletter
\newenvironment{breakablealgorithm}
{% \begin{breakablealgorithm}
		\begin{center}
			\refstepcounter{algorithm}% New algorithm
			\hrule height.8pt depth0pt \kern2pt% \@fs@pre for \@fs@ruled
			\renewcommand{\caption}[2][\relax]{% Make a new \caption
				{\raggedright\textbf{\ALG@name~\thealgorithm} ##2\par}%
				\ifx\relax##1\relax % #1 is \relax
				\addcontentsline{loa}{algorithm}{\protect\numberline{\thealgorithm}##2}%
				\else % #1 is not \relax
				\addcontentsline{loa}{algorithm}{\protect\numberline{\thealgorithm}##1}%
				\fi
				\kern2pt\hrule\kern2pt
			}
		}{% \end{breakablealgorithm}
		\kern2pt\hrule\relax% \@fs@post for \@fs@ruled
	\end{center}
}
\makeatother

\title[Sampling Low-Temperature SK model above the Almeida--Thouless Line]{Sampling Low-Temperature Sherrington-Kirkpatrick Model above the Almeida--Thouless Line}

\author{}

\address[]{}
\email[]{}

\date{\today}

\begin{document}
\begin{abstract}
We give a polynomial-time sampler for the Sherrington--Kirkpatrick model at each fixed inverse temperature $\beta>1$ and uniform field $h>0$ satisfying the strict Almeida--Thouless inequality. The conditional output law converges in total variation to the Gibbs law, both in expectation and in probability over the disorder. The construction combines a dimension-free integrated TAP-residual estimate, a cavity analysis of the actual quenched Gaussian observation experiment, fixed-depth AMP initialization, and local convexity of the full TAP functional along that experiment. Finitely many fixed convex neighborhoods cover the bounded-time part of the path, and one convex cap covers the terminal part. Conditioning on correct sign rounding gives an evaluable path target with the exact Gibbs spin marginal. An independence Metropolis correction with a polynomial likelihood cutoff yields the sampler. We specify capped numerical routines and a finite-precision implementation of the final spin output.

\textbf{Note.} This draft is completely AI-generated using GPT-6, prompted by Zhangsong Li. The only input is the strategic decision to follow the ideas developed in \cite{li2026efficient}. 
\end{abstract}
\maketitle
\enlargethispage{3pt}
\setcounter{tocdepth}{1}
\tableofcontents


\section{Introduction and main results}\label{sec:intro}

Disordered systems provide a setting for understanding how random interactions produce collective behavior. A central example is the Sherrington--Kirkpatrick (SK) model, which consists of binary spins with random interactions between every pair of spins \cite{sherrington1975solvable}. The interactions create competing preferences across the system. Despite its simple definition, the model has a rich equilibrium structure and has played a central role in the development of replica symmetry breaking, the cavity method, and iterative methods for solving the Thouless--Anderson--Palmer (TAP) equations \cite{talagrand2010mean,panchenko2013sherrington,thouless1977solution,bolthausen2014iterative}.

Let $\bW$ be a symmetric $N\times N$ random matrix with zero diagonal and independent $\mathcal N(0,N^{-1})$ entries above the diagonal. For inverse temperature $\beta\ge0$ and external field $h\ge0$, write $J=\beta\bW$ and define
\begin{equation*}
 H_J(x)=\frac12x^{\mathsf T}Jx+h\bone^{\mathsf T}x,
 \qquad \mu_J(x)=\frac{e^{H_J(x)}}{Z_J},
 \qquad Z_J=\sum_{x\in\{-1,1\}^N}e^{H_J(x)}.
\end{equation*}
Thus $\mu_J=\mu_{\beta,h}$ in the usual notation. Our normalization places no additional factor of $\beta$ on the field $h$.

At zero external field, $\beta=1$ separates the replica-symmetric and replica-symmetry-breaking phases \cite{talagrand2010mean,panchenko2013sherrington}. A positive field changes this phase diagram. Let $q=q(\beta,h)\in(0,1)$ be the unique solution of
\begin{equation}\label{eq:original-fixed-point}
 q=\bbE\tanh^2(h+\beta\sqrt q\,Z),\qquad Z\sim\mathcal N(0,1).
\end{equation}
Existence and uniqueness for $h>0$ are given in \cite[Proposition~1.3.8]{talagrand2010mean}. The de Almeida--Thouless condition is
\begin{equation}\label{eq:strict-at}
 \alpha:=\beta^2\bbE\sech^4(h+\beta\sqrt q\,Z)<1.
\end{equation}
It originated as a stability criterion for the replica-symmetric solution \cite{dealmeida1978stability}. Replica symmetry throughout this region follows from \cite[Theorem~1.1]{lopatto2026replica}; the quantitative overlap estimate in \cite[Theorem~1.1]{kusuoka2026quantitative} will be a principal input to our analysis.

We study sampling at each fixed low-temperature pair $(\beta,h)$ satisfying \eqref{eq:strict-at}. Our main result is the following.

\begin{thm}\label{thm:main}
Fix $\beta>1$ and $h>0$ satisfying \eqref{eq:strict-at}. There is a randomized algorithm, described in \S\ref{sec:completion}, whose input is $\bW$ and whose running time is polynomial in $N$, such that its conditional output law $\widehat\mu_J$ satisfies
\begin{equation}\label{eq:main-tv}
 \bbE_J\TV(\widehat\mu_J,\mu_J)\longrightarrow0.
\end{equation}
In particular, $\TV(\widehat\mu_J,\mu_J)=o_{\bbP_J}(1)$. The algorithm has a deterministic polynomial running-time bound on every input. Its constants and polynomial exponent may depend on the fixed pair $(\beta,h)$.
\end{thm}

The theorem gives asymptotically vanishing total-variation error. It does not prescribe an inverse-polynomial error rate or a bound uniform as $(\beta,h)$ approaches the AT line. The algorithm uses independent proposals and an independence Metropolis correction. In particular, the conclusion concerns sampling from the Gibbs measure; it is not a mixing-time assertion for single-site Glauber dynamics. We specify the ideal proposal using real Gaussian variables and give a finite-precision implementation of its final spin output in \S\ref{sec:numerical}.

The construction builds on algorithmic stochastic localization \cite{el2022sampling,montanari2023posterior}, the convergence theory for TAP iterations \cite{bolthausen2014iterative,chen2021convergence}, and the local-curvature analysis of \cite{brennecke2023operator,celentano2024sudakov}. The estimates needed here concern a Gaussian observation generated from a Gibbs sample under the original disorder law. We prove those estimates explicitly. The resulting proposal has a target-to-proposal likelihood ratio bounded by $N$ outside a set of vanishing target mass, which is enough for the Metropolis correction.

\subsection{Proof strategy}\label{sec:proof-overview}

We describe the main steps before specifying the algorithm. Throughout this discussion, the unknown Gibbs sample appears only in the analysis. Given $J$, draw $X\sim\mu_J$, independently draw a standard $N$-dimensional Brownian motion $B$, and set
\begin{equation}\label{eq:observation-process}
 Y_t=tX+B_t,\qquad
 m_t=\bbE[X\mid J,Y_t],\qquad
 C_t=\Cov(X\mid J,Y_t).
\end{equation}
The conditional law of $X$ is $\mu_{J,Y_t}(x)\propto\exp\{H_J(x)+Y_t^{\mathsf T}x\}$. The observation process has innovation drift $m_t$. Thus a computable approximation to $m_t$ yields a sequential Gaussian proposal for its path. At time $T_N=4\log N$, coordinatewise sign rounding recovers $X$ with probability $1-O(N^{-1})$.

\emph{A computable exact target on paths.}
The unknown partition function in the observation-path density initially prevents direct Metropolis correction. We remove it by conditioning the joint experiment on $X=\sign Y_{T_N}$. The probability of this event is $\Phi(\sqrt{T_N})^N$, independently of $X$. Consequently the conditioned path law has exactly $\mu_J$ as its rounded endpoint distribution, and its unnormalized density can be evaluated from the Hamiltonian of $\sign Y_{T_N}$ and the endpoint itself. The Gaussian proposal density is a product of known transition densities. Their ratio is therefore computable without evaluating a partition function; see \eqref{sampler:log-weight}.

\emph{Integrated TAP accuracy.}
Let $ F_{\tap}(m,f)$ be the full TAP functional with total field $f$, defined in \S\ref{sec:algorithm}. Its gradient at the posterior mean need not vanish. We prove the dimension-free integrated estimate
\[
 \int_0^\infty\bbE\|\nabla_m F_{\tap}(m_t,h\bone+Y_t)\|^2\,dt
 \le C_{\beta,h}.
\]
This estimate holds for every fixed $\beta>0$ and finite $h$, without the AT condition. Its proof combines a one-row Gaussian response bound with the posterior covariance budget $\int_0^\infty\bbE\Tr(C_t^2)\,dt\le N$. The bound concerns the true posterior mean, so an optimization argument is still needed to turn it into a computable drift.

\emph{Initialization and full local curvature.}
We first identify the scalar law of the actual observation experiment. Put $U_0=h+\beta\sqrt q\,Z_0$, draw $\xi\in\{-1,1\}$ conditionally with mean $\tanh U_0$, and define
\[
 a(v)=\bbE\tanh^2(U_0+v\xi+\sqrt v\,Z_1).
\]
The strict AT condition implies $\sup_{v\ge0}\beta^2a'(v)<1$. The posterior self-overlap converges to the unique solution
\[
 r(t)=a\bigl(t+\beta^2(r(t)-q)\bigr).
\]
A cavity argument under the original quenched law then gives a fixed-depth AMP initializer for $m_t$, uniformly on every bounded time interval. The observation is correlated with the disorder through $X$; the proof retains its exact one-row likelihood and does not apply an independent-side-information state-evolution theorem to it.

We next prove positive curvature of the full TAP Hessian on a ball of radius $c\sqrt N$ around $m_t$ at each fixed positive time, with constants uniform on compact time intervals. This requires controlling the rank-one term from the derivative of the Onsager correction. Conditional Gaussian regression and the post-AMP comparison of \cite{celentano2024sudakov} reduce the problem to a scalar quadratic form. Its two strict margins follow from the original AT inequality and the scalar observation stability above. An exact conditional change of measure transfers the reference estimate to the Gibbs-generated observation. The quantitative Gaussian coupling theorem of \cite[Theorem~4]{reeves2025dimension} supplies an exponential reference bound whose rate is independent of the original transcript depth.

\emph{Fixed domains along the path.}
The posterior mean is a martingale. Together with its scalar norm law, this gives a uniform estimate for its distance from any fixed time anchor. We therefore choose finitely many deterministic anchors and retain a fixed convex neighborhood between consecutive anchors. At time zero we use the original-model local-curvature theorem of \cite[Proposition~4.5]{brennecke2023operator}. After a sufficiently large fixed time $S$, a single convex cap centered at $\sign Y_S$ contains all subsequent posterior means with exponentially high probability. A sparse-submatrix estimate makes the full TAP functional strongly convex on every such cap. Hence the additional posterior structure is needed only on a time interval independent of $N$.

\emph{Likelihood control and correction.}
At the ordinary fine-grid times we minimize the TAP functional on the stored domain. Strong convexity bounds the squared drift error by the squared TAP residual and the optimization tolerance. A randomly shifted grid converts the integrated residual estimate into a bound on the sum of these errors. Immediately after a deterministic anchor we use zero drift for one short interval, whose cost is negligible. This avoids requiring a pointwise residual estimate at the anchors.

The simultaneous domain assertions have probability $1-o(1)$, without a specified rate. To use that qualitative control, we stop the likelihood analysis at the first current failure of a domain condition. The stopped drift action has bounded expectation. Girsanov's formula then bounds the upper tail of the path log-likelihood at level $\log N$. An independence Metropolis chain with $O(N\log N)$ proposals yields \eqref{eq:main-tv}. All stopping tests involving the unknown posterior are confined to the proof; the implemented routines have prescribed iteration budgets and are defined on every input.

\subsection{Related sampling results}\label{sec:related}

For the zero-field SK model, algorithmic stochastic localization gives normalized Wasserstein sampling guarantees throughout $\beta<1$ \cite{el2022sampling,celentano2024sudakov}. Functional inequalities and mixing estimates for Glauber dynamics have been developed in \cite{bauerschmidt2019very,eldan2022spectral,anari2024trickle,anari2022entropic,adhikari2024spectral,anari2024universality}. More recent bounds reach $\beta<1/2$ \cite{wang2026optimal} and $\beta<1/2+\epsilon$ for an absolute $\epsilon>0$ \cite[Corollary~1.6]{boban2026rank}. Polynomial-time total-variation sampling throughout $\beta<1$ is established in \cite{lee2026potential,hou2026simulated}; the latter gives inverse-polynomial accuracy and disorder failure bounds.

Concurrent public manuscripts report an $\Omega_\beta(N^{-1})$ spectral gap for random-scan Glauber dynamics \cite{openai2026d} and worst-case total-variation cutoff on the $O(N\log N)$ scale throughout $0\le\beta<1$ \cite{openai2026c}. At $\beta=1$, \cite{openai2026b} gives a worst-case lower bound of $N^{5/3-o(1)}$, and \cite{openai2026a} gives the corresponding $N^{5/3+o(1)}$ mixing-time scale. We refer to the cited manuscripts for their precise probabilistic and time-normalization conventions.

At zero field and $\beta>1$, \cite{el2022sampling} proves an obstruction for stable sampling algorithms, and \cite{sellke2025exponentially} proves exponentially slow worst-case Glauber mixing for sufficiently large $\beta$. The stretched-exponential barriers in \cite{openai2026barriers} concern typical configurations drawn from equilibrium and then used as fixed initial states. With a positive field, Bolthausen's iteration converges in the asymptotic mean-square sense under the AT condition, with geometric decay under strict stability \cite[Theorem~2.1]{bolthausen2014iterative}. This convergence takes the system-size limit before the iteration limit. Polynomial-time Glauber mixing is known for every fixed $\beta$ at sufficiently large field \cite[Corollary~1.2]{bandeira2026mixing}; high-magnetization sampling is also treated in \cite{kumar2026high}.

The independence-correction principle is closely related to the sequential posterior sampler for $\mathbb Z_2$ synchronization in \cite{li2026efficient}. There the proposal reveals spin coordinates and approximates their conditional logits. Here the auxiliary Gaussian observation provides the covariance budget used by the integrated residual estimate and a fixed terminal cap. Our proposal is consequently defined on observation paths, and the correction uses the exact sign-conditioned path target. The low-temperature contribution is the treatment of every fixed positive-field pair satisfying strict AT stability.

The algorithm is specified in \S\ref{sec:algorithm}; \S\ref{sec:analysis} collects its probabilistic and optimization estimates, and \S\ref{sec:completion} proves the main theorem. The appendices contain the Gaussian response, quenched cavity, AMP, curvature, and numerical arguments. Known external theorems are stated with their sources where they are used.


\section{Overview of the algorithm}\label{sec:algorithm}

The sampler first generates a Gaussian observation path using approximate
posterior means. It then applies an independence Metropolis correction on
the space of these paths and returns the signs of the final observation.
This section specifies both steps. The estimates underlying the proposal
are established in \S\ref{sec:analysis}, and \S\ref{sec:completion}
combines them to prove Theorem~\ref{thm:main}.

Put $J=\beta\bW$ and write $H_J(x)=\frac12x^{\mathsf T}Jx+
h\mathbf1^{\mathsf T}x$ and $\mu_J(x)=Z_J^{-1}e^{H_J(x)}$.
Vector norms are Euclidean unless specified otherwise. The algorithm
uses a fixed choice of depths, radii, and time anchors depending on
$(\beta,h)$; none of these quantities depends on $N$. Exact Gaussian
draws specify its ideal law. The deterministic routines have a prescribed
polynomial operation budget on every input, and \S\ref{sec:numerical}
gives a finite-precision implementation with an additional $o(1)$ error
in the law of the returned spins.

\subsection{Metropolis correction and an exact path target}

We first recall a bound for the independence Metropolis sampler. Its useful
feature here is that the likelihood ratio need only be controlled outside
a set of small target mass.

\begin{lem}[Independence Metropolis with a likelihood cutoff]
\label{lem:imh}
Let $\pi$ and $\nu$ be equivalent probability measures, let
$w=d\pi/d\nu$, and let $K$ be the independence Metropolis kernel with
target $\pi$ and proposal $\nu$. For $B\ge1$, put
$\delta_B=\pi(w>B)$. If the initial state has law $\nu$, then
\begin{equation*}
 \|\nu K^k-\pi\|_{\TV}
 \le \delta_B+\frac{\delta_B}{B}
       +\exp\left\{-\frac{(1-\delta_B)k}{B}\right\}.
\end{equation*}
\end{lem}
\begin{proof}
Write $a(z)=\int\min\{1,w(z')/w(z)\}\,\nu(dz')$.
The independence-sampler bound of
\cite[Theorem~6]{roberts2011quantitative}, averaged over the initial
state, gives
\[
 \|\nu K^k-\pi\|_{\TV}
 \le \mathbb E_{Z\sim\nu,\,Z'\sim\pi}
       [\{1-\min(a(Z),a(Z'))\}^k].
\]
On $\{w\le B\}$, the acceptance probability is at least
$(1-\delta_B)/B$, and $\nu(w>B)\le\delta_B/B$.
Splitting the expectation according to whether either initial likelihood
ratio exceeds $B$ proves the claim. We apply this result to a continuous
path state with retained auxiliary random variables.
\end{proof}

To obtain an evaluable target on that state space, consider the exact
observation experiment
\[
 X\sim\mu_J,\qquad Y_t=tX+B_t,
 \qquad m_t=\mathbb E[X\mid J,(Y_s)_{s\le t}],
\]
where $B$ is standard Brownian motion independent of $(J,X)$.
The observation at time $t$ is sufficient for its entire history, and
the conditional spin law is
$\mu_{J,Y_t}(x)\propto e^{H_J(x)+Y_t^{\mathsf T}x}$.
Let $0=t_0<\cdots<t_L=T$ be a grid independent of this experiment and
write $\Delta_j=t_{j+1}-t_j$. If the grid is randomized, retain the
randomness specifying it in the state. Conditional on this randomness,
let $\mathcal W$ be the law of centered Gaussian increments with
covariances $\Delta_j I_N$, and let $P_J$ be the law of the observations
on the grid. Then
\[
 \frac{dP_J}{d\mathcal W}(y)
 =e^{-NT/2}\frac{Z_J(y_L)}{Z_J},
 \qquad Z_J(z)=\sum_xe^{H_J(x)+z^{\mathsf T}x}.
\]
This density contains the posterior partition function. The following
conditioning removes it from the target used for the correction.

\begin{thm}[Exact rounding-conditioned target]
\label{thm:sign-target}
Let $p_T=\Phi(\sqrt T)^N$, where $\Phi$ is the standard normal
distribution function. The probability law
\begin{equation}\label{eq:sign-target}
 \frac{d\Pi_J}{d\mathcal W}(y)
 =\frac{\exp\{H_J(\sign y_L)+\|y_L\|_1-NT/2\}}{Z_Jp_T}
\end{equation}
has pushforward $\mu_J$ under $y\mapsto\sign y_L$, and
$d\Pi_J/dP_J\le p_T^{-1}$. For $T=4\log N$,
$1-p_T\le N\Phi(-\sqrt T)=O(N^{-1})$.
\end{thm}
\begin{proof}
In the joint experiment condition on
$E_T=\{X=\sign Y_T\}$. For every fixed $X=x$, each coordinate is
rounded correctly with probability $\Phi(\sqrt T)$, independently of
the others. Thus $\mathbb P(E_T\mid X,J)=p_T$, so conditioning leaves
the spin marginal equal to $\mu_J$. The likelihood of the observations
conditional on $X=x$ is
$\exp\{x^{\mathsf T}y_L-NT/2\}$ relative to $\mathcal W$.
Restricting the sum over $x$ to $x=\sign y_L$ proves
\eqref{eq:sign-target}. Moreover,
\[
 \frac{d\Pi_J}{dP_J}(y)
   =\frac{\mu_{J,y_L}(\sign y_L)}{p_T}\le p_T^{-1}.
\]
The last assertion follows from the Gaussian tail bound and a union
bound. These statements remain valid with any independent grid offset
and algorithmic seeds retained in the state.
\end{proof}

Suppose a proposal $Q_J$ generates its observations sequentially through
\begin{equation*}
 y_{j+1}=y_j+\Delta_j b_j+\sqrt{\Delta_j}\,Z_j,
 \qquad Z_j\stackrel{\mathrm{iid}}{\sim}\mathcal N(0,I_N),
\end{equation*}
where $b_j\in[-1,1]^N$ is an evaluable function of $J$, the retained
randomness, and the observations up to $t_j$. Its density is
\[
 \frac{dQ_J}{d\mathcal W}(y)
 =\exp\left\{\sum_{j=0}^{L-1}b_j^{\mathsf T}(y_{j+1}-y_j)
                   -\frac12\sum_{j=0}^{L-1}\Delta_j\|b_j\|^2\right\}.
\]
Consequently, up to a constant independent of the path and all retained
randomness, the log likelihood ratio of $\Pi_J$ to $Q_J$ is
\begin{equation}\label{sampler:log-weight}
 \mathcal L_J(y)
 =H_J(\sign y_L)+\|y_L\|_1
  -\sum_{j=0}^{L-1}b_j^{\mathsf T}(y_{j+1}-y_j)
  +\frac12\sum_{j=0}^{L-1}\Delta_j\|b_j\|^2.
\end{equation}
The independence Metropolis acceptance probability is therefore
$\min\{1,\exp(\mathcal L_J(y')-\mathcal L_J(y))\}$.
Computing it requires the Hamiltonian and the proposal drifts, together
with the recorded grid; neither partition function is evaluated.
When a path is scored again, its own anchors, offset, and seeds are used
to reproduce exactly the deterministic maps defining its proposal.

\subsection{Sequential construction of the Gaussian proposal}

Set $T_N=4\log N$. Fix a terminal time $S>0$ and a deterministic mesh
$0=s_0<s_1<\cdots<s_K=S$, both independent of $N$.
The next subsection specifies the constraint set stored at each anchor.
We take
\[
 M=\lceil T_NN^4\rceil,\qquad \eta=T_N/M\le N^{-4},
 \qquad d_N=N^{-20}.
\]
Remove from $(0,\eta)$ the $d_N$-neighborhoods of its endpoints and of
the residues of the anchors $s_j$ modulo $\eta$. Denote the remaining
finite union of intervals by $\mathcal U_N$. For all sufficiently large
$N$, it has length at least $\eta/2$. Sample $U$ uniformly from
$\mathcal U_N$ and use the sorted union
\begin{equation}\label{sampler:grid}
 \{0,T_N\}\ \cup\ \{s_1,\ldots,s_K\}\ \cup\
 \{U,U+\eta,\ldots,U+(M-1)\eta\}.
\end{equation}
Increase a fixed small-$N$ cutoff so that $T_N>S$, the fixed anchors are
separated by more than $\eta$, and the preceding construction is valid.
Below this cutoff the algorithm returns independent uniform signs.
Every interval of \eqref{sampler:grid} then has length between $d_N$
and $\eta$.

The offset is retained as part of every proposed state. It is sampled
afresh for each independent proposal, and has the same distribution in
the reference and target experiments. Its role in the analysis is the
quadrature inequality
\begin{equation}\label{sampler:quadrature}
 \mathbb E_U\left[\eta\sum_{j=0}^{M-1}g(U+j\eta)\right]
 \le2\int_0^{T_N}g(t)\,dt
 \qquad(g\ge0),
\end{equation}
which follows directly from the density bound
$\mathbf1_{\mathcal U_N}(u)/|\mathcal U_N|\le2/\eta$.
The lower bound on interval lengths will be useful only for numerical
implementation.

At a deterministic anchor $s_j<S$, compute an AMP approximation
$A_j=A_j(J,y_{s_j})$ using the fixed depths in
\eqref{amp:two-phase}, and store a clipped center and its constraint
set. At $S$, store $\tau=\sign y_S$ and a terminal constraint set.
On the interval immediately following any deterministic anchor,
including zero and $S$, use $b_j=0$. At every other left endpoint,
which is a shifted fine-grid point, approximately minimize the TAP
functional on the currently stored set and use the resulting vector
as $b_j$. These instructions depend only on observations already
generated.

The zero drift after an anchor makes the random-grid argument exact:
the integrated residual estimate \eqref{eq:integrated-tap} controls
the evaluations covered by \eqref{sampler:quadrature}, while the
finitely many other outgoing intervals have total squared-drift cost
at most $(K+1)N\eta$. No pointwise residual estimate at a prescribed
anchor is needed.

We also impose explicit guards. An input with $\|J\|_F>N^2$ is
rejected and assigned a uniform spin output. Along every proposal,
if $\|y_j\|_\infty>N^2$ at a grid point, all subsequent drifts are
zero and all remaining anchor and optimization routines are skipped.
Every other routine has a fixed iteration budget; an unsuccessful
routine returns a bounded vector according to a deterministic fallback
rule. The implementation never tests whether a true posterior mean is
inside a constraint set.

\begin{breakablealgorithm}
\caption{Generation and scoring of one Gaussian-path proposal}
\label{alg:gaussian-proposal}
\begin{algorithmic}[1]
\Require $J$, the fixed anchor record, and the tolerances specified below.
\Ensure An augmented path $y\sim Q_J$ and its log weight $\mathcal L_J(y)$.
\State Sample and retain $U\sim\operatorname{Unif}(\mathcal U_N)$; form
the grid \eqref{sampler:grid} and set $y_0=0$.
\State Set the guard flag to inactive.
\For{$j=0,\ldots,L-1$}
  \If{$\|y_j\|_\infty>N^2$}
    \State Activate the guard flag permanently.
  \EndIf
  \If{the guard flag is active}
    \State $b_j\gets0$.
  \ElsIf{$t_j=s_a<S$ is a deterministic anchor}
    \State Compute $A_a(J,y_j)$ by \eqref{amp:two-phase}; store
    $c_a$ and $\mathcal K_a$ from \eqref{sampler:ball-domain}; set
    $\mathcal K\gets\mathcal K_a$.
    \State $b_j\gets0$.
  \ElsIf{$t_j=S$}
    \State Store $\tau\gets\sign y_j$ and
    $\mathcal K\gets K_\tau(\rho)$; set $b_j\gets0$.
  \Else
    \State $b_j\gets\mathsf{TAPApprox}(J,h\mathbf1+y_j,\mathcal K)$,
    using Algorithm~\ref{alg:tap-optimization} and the currently stored set.
  \EndIf
  \State Draw an independent $Z_j\sim\mathcal N(0,I_N)$ and set
  $y_{j+1}\gets y_j+\Delta_jb_j+\sqrt{\Delta_j}Z_j$.
  \State Retain $b_j$ and $\Delta_j$ for the score calculation.
\EndFor
\State Compute $\mathcal L_J(y)$ from \eqref{sampler:log-weight}.
\State \Return the offset, path, any retained seeds, and the score.
\end{algorithmic}
\end{breakablealgorithm}

The AMP routine can be deterministic. If it uses auxiliary seeds, sample
them independently in advance and retain them in the proposed state.
Their prescribed prior law is also used under $P_J$ and $\Pi_J$, so it
cancels in the likelihood ratio. All fallbacks are included in the
definition of the proposal and are reproduced when its score is
evaluated.

\subsection{TAP optimization and the fixed constraint sets}

For $m\in[-1,1]^N$ and a total field $f\in\mathbb R^N$, use the
minimization convention
\begin{equation*}
  F_{\tap}(m,f)=-\frac12m^{\mathsf T}Jm-f^{\mathsf T}m
 +\sum_i I(m_i)-\frac{N\beta^2}{4}
                 \left(1-\frac{\|m\|^2}{N}\right)^2,
\end{equation*}
where
$I(u)=\frac{1+u}{2}\log\frac{1+u}{2}
      +\frac{1-u}{2}\log\frac{1-u}{2}$,
with its continuous extension at $\pm1$.
On the open cube,
\[
 \nabla F_{\tap}(m,f)=\atanh m-f-Jm
             +\beta^2\left(1-\frac{\|m\|^2}{N}\right)m.
\]

The full Hessian is
\begin{equation}\label{eq:tap-hessian}
 \nabla_m^2 F_{\tap}(m,f)
 =\operatorname{diag}\bigl((1-m_i^2)^{-1}\bigr)-J
  +\beta^2\left(1-\frac{\|m\|^2}{N}\right)I_N
  -\frac{2\beta^2}{N}mm^{\mathsf T}.
\end{equation}
We retain both terms produced by differentiating the Onsager correction.

For fixed nonnegative integers $L_{\mathrm{old}},L_{\mathrm{obs}}$, the anchor routine uses
\begin{equation}\label{amp:two-phase}
\begin{split}
 a^{-1}&=0,\qquad a^0=\sqrt q\,\bone,\\
 a^{s+1}&=\tanh\left(h\bone+\mathbf1_{\{s\ge L_{\mathrm{old}}\}}y
 +Ja^s-\beta^2\left(1-\frac{\|a^s\|^2}{N}\right)a^{s-1}\right),
 \qquad 0\le s<L_{\mathrm{old}}+L_{\mathrm{obs}}.
\end{split}
\end{equation}
Its output is $A(J,y)=a^{L_{\mathrm{old}}+L_{\mathrm{obs}}}$. The first phase ignores $y$; the second inserts it without resetting the previous iterate. Both depths are fixed independently of $N$. In the initializer proofs we abbreviate $L_{\mathrm{old}}=L$ and $L_{\mathrm{obs}}=k$.

The Hessian \eqref{eq:tap-hessian} does not depend on $f$. We therefore retain a fixed
constraint set between consecutive anchors while the field changes.

Choose $\rho>0$ as in \S\ref{sec:terminal}, and put
$S=2\log(8/\rho)$. Theorem~\ref{curv:fixed-anchor},
Theorem~\ref{amp:actual}, and Lemma~\ref{sampler:anchors} permit fixed
constants $a>0$, $0<\gamma\le1/2$, fixed AMP depths, and a fixed
mesh $0=s_0<\cdots<s_K=S$ such that, with probability $1-o(1)$ under
the true observation law, simultaneously for every $j<K$,
\begin{align}
 \|A_j(J,Y_{s_j})-m_{s_j}\|&\le\tfrac a4\sqrt N,
       \label{sampler:anchor-accuracy}\\
  F_{\tap}(\cdot,f)&\text{ is $\gamma$-strongly convex on }
       [-1,1]^N\cap B(m_{s_j},4a\sqrt N),
       \label{sampler:anchor-curvature}\\
 \sup_{s_j\le t\le s_{j+1}}\|m_t-m_{s_j}\|
       &\le\tfrac a2\sqrt N.
       \label{sampler:anchor-increments}
\end{align}
The field $f$ in \eqref{sampler:anchor-curvature} is arbitrary.
The choices are made in this order: first obtain common positive
curvature and neighborhood radii on the bounded time interval, then
choose the anchor mesh by continuity of the scalar overlap and the
anchor-distance estimate. In particular, the radii are chosen before
the number of anchors. The original-field curvature input handles
time zero, as explained in the proof of the anchor construction.

Set
\[
 \varepsilon_N=N^{-8},\qquad
 u_N=\tfrac14\min(\rho,a,1)N^{-30},\qquad
 \mathcal B_N=[-1+u_N,1-u_N]^N.
\]
At an anchor $s_j<S$, store
\begin{equation}\label{sampler:ball-domain}
 c_j=\operatorname{clip}_{\mathcal B_N}(A_j),\qquad
 \mathcal K_j=[-1,1]^N\cap B(c_j,2a\sqrt N),
 \qquad s_j\le t<s_{j+1}.
\end{equation}
For all sufficiently large $N$, clipping moves the center by at most
$a\sqrt N/4$. Thus \eqref{sampler:anchor-accuracy}--
\eqref{sampler:anchor-increments} imply that every true mean in this
time interval is within $a\sqrt N$ of $c_j$ and lies in
$\mathcal K_j$. Moreover, every point of $\mathcal K_j$ is within
$\frac52a\sqrt N$ of $m_{s_j}$, so the set lies inside the
neighborhood in \eqref{sampler:anchor-curvature}.

At the terminal anchor store $\tau=\sign y_S$ and use
\begin{equation*}
 \mathcal K_{\mathrm{term}}=K_\tau(\rho)
   =\left\{m\in[-1,1]^N:
             \sum_i(1-\tau_i m_i)\le\rho N\right\},
 \qquad S\le t\le T_N.
\end{equation*}
The estimates in \S\ref{sec:terminal} show that this one cap contains
all later true means and is a domain of $1/2$-strong convexity, except
on an event of exponentially small probability. We freeze its sign
vector at $S$; later observations change the field but not the cap.
All sets used by the proposal are consequently computable from the
observations available when they are first stored.

At each ordinary drift-evaluation time, the optimization goal is
\begin{equation}\label{eq:algorithm-objective}
 \widehat m\in\mathcal K\cap\mathcal B_N,\qquad
  F_{\tap}(\widehat m,f)-\min_{z\in\mathcal K} F_{\tap}(z,f)
       \le\varepsilon_N.
\end{equation}
The minimum in this expression is over the original, unclipped set.
On a correct domain, Lemma~\ref{sampler:objective} then bounds
$\|\widehat m-m_t\|^2$ by a constant multiple of the squared TAP
residual at $m_t$, plus $O(\varepsilon_N)$.
In particular, the argument does not require the true mean to be in
the clipped cube.

Here are the deterministic choices needed to implement this goal.
Coordinate clipping onto $\mathcal B_N$ preserves a ball domain
because it fixes $c_j$ and is nonexpansive. For a cap, the map
\[
 z\longmapsto(1-\zeta_N)z+\zeta_N(1-\rho/2)\tau,
 \qquad \zeta_N=2u_N/\rho,
\]
takes $K_\tau(\rho)$ into $K_\tau(\rho)\cap\mathcal B_N$.
Under the input and field guards, the entropy continuity bound and
the explicit quadratic terms show that either modification changes
the objective by at most
\[
 C_{\beta,h,a,\rho}
       \{N^4u_N+Nu_N\log(e/u_N)\}=o(N^{-8}).
\]
Thus it is enough to minimize on the clipped domain to accuracy
$\varepsilon_N/2$.
On that domain a deterministic gradient-Lipschitz bound is
\[
 L_N=u_N^{-1}+\|J\|_F+3\beta^2.
\]
Choose $B_N=C_{\beta,h}N^4$ large enough to bound the relevant
objective range and the gradient bound times the domain diameter.
Projected gradient descent with step $1/L_N$ is run for
\[
 n_{\mathrm{opt}}
   =\left\lceil\frac{2L_N}{\gamma}
                    \log\frac{8B_N}{\varepsilon_N}\right\rceil
\]
iterations. Its convergence on a $\gamma$-strongly convex domain
follows directly from \cite[Theorem~3.10]{bubeck2015convex}; the
gradient bound converts the squared-distance estimate there into
the required objective accuracy. There is no convergence-dependent
stopping rule. The same number of iterations is prescribed when
the domain lacks the required convexity.

\begin{breakablealgorithm}
\caption{TAP approximation on a stored constraint set}
\label{alg:tap-optimization}
\begin{algorithmic}[1]
\Require Guarded $J,f$, a stored ball or cap $\mathcal K$, and its
fixed feasible center $z^0$.
\Ensure A bounded vector, satisfying \eqref{eq:algorithm-objective}
on a correct convex domain.
\State Set $\mathcal D\gets\mathcal K\cap\mathcal B_N$ and
$z\gets z^0$; use $z^0=c_j$ for a ball and
$z^0=(1-\rho/2)\tau$ for a cap.
\For{$r=1,\ldots,n_{\mathrm{opt}}$}
  \State $z\gets\operatorname{Proj}_{\mathcal D}
             (z-L_N^{-1}\nabla F_{\tap}(z,f))$,
  using the fixed feasible numerical maps in \S\ref{sec:numerical}.
\EndFor
\State \Return $z$; if a prescribed numerical routine fails within
its fixed budget, return the feasible center $z^0$ according to the
deterministic fallback rule.
\end{algorithmic}
\end{breakablealgorithm}

Both projections reduce to one-dimensional monotone searches. For a
ball intersected with the box, the projected coordinates are
\[
 z_i(\lambda)=\operatorname{clip}_{[-1+u_N,1-u_N]}
             \left(\frac{v_i+\lambda c_{j,i}}{1+\lambda}\right),
 \qquad \lambda\ge0,
\]
with $\lambda=0$ when the box projection lies in the ball and otherwise
chosen to give distance $2a\sqrt N$ from $c_j$. For a cap they are
$\operatorname{clip}_{[-1+u_N,1-u_N]}(v_i+\lambda\tau_i)$,
with $\lambda\ge0$ enforcing the cap constraint. Feasible bisection
and sufficiently small inward rounding give polynomial-time
projection routines. A gradient evaluation costs $O(N^2)$ arithmetic
operations. The tolerances and safeguards in \S\ref{sec:numerical}
make every call terminate within a deterministic polynomial budget,
while retaining \eqref{eq:algorithm-objective} on the structural event.


\section{Analysis of the Gaussian proposal}\label{sec:analysis}

We collect the estimates used in the proposal analysis. The first subsection records the external inputs, including their parameter and probability requirements. The subsequent statements concern the actual quenched observation experiment \eqref{eq:observation-process}. Their longer proofs are given in the appendices. We finish by constructing the fixed anchor neighborhoods and deriving the deterministic residual-to-mean estimate used in the likelihood argument.

\subsection{Probabilistic inputs}\label{sec:inputs}

Write $\langle\cdot\rangle$ for conditional Gibbs expectation. Independent replicas always mean independent samples conditional on the displayed disorder and observations. All unqualified expectations include the original Gaussian disorder. For a posterior mean $m_t$, its covariance is $C_t$, as in \eqref{eq:observation-process}. The following Gaussian-channel identities hold for an arbitrary prior on the discrete cube.

\begin{lem}[Gaussian observation identities]\label{lem:filtering-input}
The innovation process $\widetilde B_t=Y_t-\int_0^t m_s\,ds$ is Brownian motion in the observation filtration, and
\begin{equation}\label{eq:posterior-budget}
 dY_t=m_t\,dt+d\widetilde B_t,\qquad
 dm_t=C_t\,d\widetilde B_t,\qquad
 \int_0^\infty\bbE\Tr(C_t^2)\,dt
       =N-\bbE\|m_0\|^2\le N.
\end{equation}
If $\mathcal T_t$ is the third centered posterior moment tensor, then
\[
 d(C_t)_{ij}=\sum_k(\mathcal T_t)_{ijk}\,d\widetilde B_{t,k}
                  -(C_t^2)_{ij}\,dt.
\]
All these assertions also hold conditional on an environment determining the prior.
\end{lem}

We use the filtering theorem and vector-MMSE identities of \cite[Theorem~2 and Eq.~(2.25)]{elalaoui2022information}; see also the stochastic-localization formulation in \cite{montanari2023posterior}. The covariance equation is its application to the bounded observables $X_iX_j$. The infinite-time identity follows from the same identity on $[0,T]$ and $m_T\to X$ in $L^2$, since $Y_T/T\to X$. These statements require neither Gaussian interactions nor replica symmetry.

\begin{lem}[Original overlap and pressure]\label{lem:original-overlap}
Fix $(\beta,h)$ satisfying \eqref{eq:strict-at}, and let $R_{12}=N^{-1}(\sigma^1)^{\mathsf T}\sigma^2$. There is a constant $C<\infty$ such that
\[
 \bbE\langle(R_{12}-q)^2\rangle\le C/N,
 \qquad \bbE\Tr(C_0^2)\le CN.
\]
Moreover,
\[
 N^{-1}\log Z_J\xrightarrow{\bbP_J}f_{\mathrm{RS}},\qquad
 f_{\mathrm{RS}}=\bbE\log\bigl(2\cosh(h+\beta\sqrt q\,Z)\bigr)
                     +\frac{\beta^2}{4}(1-q)^2.
\]
The same conclusions hold uniformly after deleting any fixed number of sites, retaining the original interaction variance $\beta^2/N$.
\end{lem}

The overlap and expected-pressure assertions are \cite[Theorem~1.1]{kusuoka2026quantitative}. Gaussian concentration \cite[Section~5.4]{boucheron2013concentration} upgrades expected-pressure convergence to convergence in probability. To obtain the covariance bound, use the replica identity
\[
 \frac{\Tr(C_0^2)}{N^2}
 =\left\langle (R_{12}-q)^2
       -2(R_{12}-q)(R_{13}-q)
       +(R_{12}-q)(R_{34}-q)\right\rangle
\]
and Cauchy--Schwarz. A cavity model with $r=N-d$ sites has effective inverse temperature $\beta\sqrt{r/N}$. For fixed $d$ these parameters remain in a compact subset of strict AT, so the uniform version of the cited theorem applies. Its fixed point differs from $q$ by $O(N^{-1})$: the derivative in $q$ of the fixed-point map is
\[
 \beta^2\bbE[\sech^2U(1-3\tanh^2U)]\le\alpha<1,
 \qquad U=h+\beta\sqrt q\,Z.
\]
This also explains why all cavity estimates below retain the original entry variance.

\begin{lem}[Original TAP iterations and local curvature]\label{lem:original-tap-input}
Let $m^{(k)}$ be the canonical Bolthausen iterates for the original uniform-field model. Under strict AT they satisfy
\[
 \lim_{k\to\infty}\limsup_{N\to\infty}
       N^{-1}\bbE\|m^{(k)}-m_0\|^2=0.
\]
There are positive constants $\epsilon,c$ such that, at every sufficiently large fixed depth $k$, with probability tending to one,
\[
 \operatorname{diag}\bigl((1-w_i^2)^{-1}\bigr)-J
       +\beta^2(1-q)I_N-\frac{2\beta^2}{N}ww^{\mathsf T}
       \succeq cI_N
\]
for every $w\in(-1,1)^N$ with $\|w-m^{(k)}\|\le\epsilon\sqrt N$. The original finite-depth Gaussian transcript also has its canonical empirical state evolution and conditional Gaussian regression representation.
\end{lem}

The mean-consistency assertion follows from \cite[Theorem~3]{chen2021convergence}, whose local-uniform overlap hypothesis is supplied by Lemma~\ref{lem:original-overlap}. The contraction of the original iteration is \cite[Theorem~2.1]{bolthausen2014iterative}. The displayed curvature statement is \cite[Proposition~4.5 and Remark~4.5]{brennecke2023operator}; its proof requires strict AT, without the additional hypothesis attached to the posterior-covariance conclusions elsewhere in that paper. We use the finite-depth transcript representation from \cite[Section~4]{brennecke2023operator} and the Gaussian regression formulation in \cite[Appendix~A.1]{celentano2024sudakov}.

The cited curvature matrix uses $q$ in its scalar diagonal term. For $\|w-m^{(k)}\|\le\epsilon\sqrt N$,
\[
 \left|N^{-1}\|w\|^2-q\right|
 \le2\epsilon+\left|N^{-1}\|m^{(k)}\|^2-q\right|.
\]
The last term tends to zero at each fixed depth. Shrinking $\epsilon$ therefore gives local positivity of the full Hessian \eqref{eq:tap-hessian}. To use a full-GOE version of an external theorem, add an independent Gaussian diagonal. It changes the Hamiltonian by a spin-independent constant and has operator norm $o_{\bbP}(1)$, so the neighborhood estimate transfers to the hollow matrix $J$. The integrated residual estimate below always uses the hollow matrix itself.

\subsection{The integrated TAP residual}\label{sec:residual}

The first new estimate controls the TAP gradient along the entire observation path. It will allow us to accumulate approximation errors without proving a pointwise quantitative posterior-overlap bound at every time.

\begin{thm}[Integrated TAP residual]\label{thm:integrated-tap}
For every fixed $\beta>0$ and finite $h$, there is $C_{\beta,h}<\infty$
such that the true quenched observation process satisfies
\begin{equation}\label{eq:integrated-tap}
 \int_0^\infty\bbE_{J,X,B}
 \|\nabla_m F_{\tap}(m_t,h\bone+Y_t)\|_2^2\,\de t
 \le C_{\beta,h}.
\end{equation}
The constant is independent of $N$. No AT assumption is required.
\end{thm}

The proof is in \S\ref{sec:integrated-proof}. It compares one interaction row with an independent cavity observation experiment and then applies \eqref{eq:posterior-budget} to each cavity prior. The likelihood comparison has a constant independent of observation time, which is essential for integrating to infinity.

\subsection{Scalar parameters and posterior overlap concentration}\label{sec:scalar-summary}

Let $U=h+\beta\sqrt q\,Z_0$, draw $\xi\in\{-1,1\}$ conditionally with $\bbE[\xi\mid U]=\tanh U$, and let $Z_1$ be independent standard Gaussian. Define
\begin{equation}\label{eq:scalar-observation-law}
 U_v=U+v\xi+\sqrt v\,Z_1,
 \qquad a(v)=\bbE\tanh^2U_v,
 \qquad u(v)=\bbE\sech^4U_v.
\end{equation}

\begin{lem}[Uniform scalar stability]\label{lem:scalar-stability}
The function $a$ satisfies $a(0)=q$, $a'(v)=u(v)$, and
\[
 \beta^2\sup_{v\ge0}u(v)\le\kappa<1.
\]
One may take $\kappa=\alpha$ when $h\le\beta^2q$ and $\kappa=\beta^2(1-q)$ when $h>\beta^2q$.
\end{lem}
\begin{proof}
The derivative identity is the conditional scalar-MMSE formula \cite[Proposition~9 and Corollary~2]{guo2011estimation}. To identify the semigroup used for stability, Bayes' rule gives, for a bounded test $f$,
\[
 \bbE[f(U_v)\mid U=x]
 =\frac{e^{-v/2}}{\cosh x}\,
      \bbE[f(x+\sqrt v\,Z_1)\cosh(x+\sqrt v\,Z_1)].
\]
This is the tilted heat semigroup of \cite[Section~2.1, proof of Lemma~2.1]{kusuoka2026quantitative}. When $h\le\beta^2q$, the monotonicity argument there applies to every semigroup time $v\ge0$ and gives $u(v)\le u(0)$; the equality case in $h$ follows by continuity. We invoke that calculation directly.

For $h>\beta^2q$, average posterior variance decreases under further observation, so $u(v)\le\bbE\sech^2U_v\le1-q$. The Gaussian comparison and linear-estimator argument in the same proof give, with $s^2=\beta^2q$,
\[
 1-q\le\bbE\sech^2(s^2+sZ_1)<\frac1{1+s^2}.
\]
The strict inequality holds because the conditional-mean estimator of a uniform sign in this Gaussian channel is nonlinear and differs from the linear estimator with mean-square error $(1+s^2)^{-1}$ on a set of positive probability. Since $q>0$, rearranging gives $\beta^2(1-q)<1$.
\end{proof}

For each $t\ge0$, contraction gives a unique $r(t)\in[q,1)$ satisfying
\begin{equation}\label{eq:observed-fixed-point}
 r(t)=a(v(t)),\qquad v(t)=t+\beta^2(r(t)-q).
\end{equation}
It is continuously differentiable, $r(0)=q$, and
$r'(t)=u(v(t))/(1-\beta^2u(v(t)))\le(1-\kappa)^{-1}$.
Write $R_N(t)=N^{-1}\|m_t\|^2$.

\begin{prop}\label{prop:scalar-law}
For every $T\ge0$,
\[
 \int_0^T\mathbb E\bigl(R_N(t)-r(t)\bigr)^2dt
 \le C_{\beta,h}\left(\frac{T}{\sqrt N}+\sqrt{\frac TN}\right).
\]
For every fixed $T<\infty$, posterior replica overlaps satisfy
\[
 \sup_{0\le t\le T}
 \mathbb E\left\langle(R_{12}-r(t))^2\right\rangle_{Y_t}\longrightarrow0.
\]
If $\varphi:[-1,1]\to\mathbb R$ is $C^2$, then
\[
 \sup_{0\le t\le T}
 \left|\frac1N\sum_i\varphi(m_{t,i})
 -\mathbb E\varphi\!\left(\tanh(U+v(t)\xi+\sqrt{v(t)}Z')\right)\right|
 \xrightarrow{\mathbb P}0,
 \qquad v(t)=t+\beta^2(r(t)-q).
\]
The probability and expectations include the original, unplanted SK
disorder and its Gibbs-generated observation process.
\end{prop}

The proof in \S\ref{sec:scalar-law} retains the exact cavity-row likelihood of the observation. It first proves an integrated self-overlap estimate using the covariance budget, and then uses the posterior martingale to obtain the uniform qualitative statements. No independent-field overlap theorem is applied to $Y_t$.

\subsection{AMP initialization under the actual law}\label{sec:amp-summary}

The two-phase recursion \eqref{amp:two-phase} first approximates the original mean and then inserts the observed field while retaining its Onsager memory. It has the following approximation guarantee.

\begin{thm}[Actual-law AMP initializer]\label{amp:actual}
Assume the strict AT hypothesis. For every fixed $T<\infty$ and
$\eta>0$, there are fixed nonnegative integers $L,k$, depending only on
$\beta,h,\eta$, such that the explicit algorithm
\eqref{amp:two-phase} satisfies
\[
 \mathbb P\left\{\sup_{0\le t\le T}
     \frac{\|a^{L+k}(Y_t)-m_t\|_2}{\sqrt N}\le\eta\right\}
 \longrightarrow1.
\]
Its running time per observation vector is $O((L+k)N^2)$ arithmetic
operations. Both depths are fixed independently of $N$.
\end{thm}

The proof is in \S\ref{sec:initializer-proof}. A finite self-avoiding cavity recursion identifies the relevant state evolution under the true observation law. A fixed-depth comparison then transfers that result to AMP. The matrix derivatives in this comparison hold the observation fixed and differentiate its likelihood together with the Gaussian matrix density. This is the additional step needed beyond original-model AMP state evolution.

\subsection{Curvature of the full TAP functional}\label{sec:curvature-summary}

Let $\mathcal H_J(w)=\nabla_m^2 F_{\tap}(w,f)$; the expression is independent of $f$ by \eqref{eq:tap-hessian}.

\begin{thm}[Uniform fixed-anchor curvature]\label{curv:fixed-anchor}
Fix $0<\delta<T_0<\infty$. There are constants $\rho,\gamma>0$, depending
only on $\beta,h,\delta,T_0$, such that for every deterministic finite set
$\mathcal T\subset[\delta,T_0]$ fixed independently of $N$,
\begin{equation}\label{curv:actual-neighborhood}
 \mathbb P\left\{
   \mathcal H_J(w)\succeq\gamma I_N
   \text{ for every }t\in\mathcal T\text{ and every }
   w\in(-1,1)^N\text{ with }\|w-m_t\|\le\rho\sqrt N
                 \right\}\longrightarrow1.
\end{equation}
The same constants work for every such finite set. In particular, they
are chosen before the anchor mesh.
\end{thm}

We prove this theorem in \S\ref{curv:section}. The original-model result in Lemma~\ref{lem:original-tap-input} is used only at time zero and for the original transcript. At positive observation times the proof requires a new conditional comparison. In particular, the longitudinal matrix calculation retains the negative rank-one term $-2\beta^2ww^{\mathsf T}/N$ in the full Hessian.

The quantitative transfer has a specific order of choices. For a prescribed fixed curvature and mean tolerance, choose the new AMP depth and a conditional reference exponent $c>0$ before choosing a sufficiently large original transcript depth. At each fixed original depth the original good event has probability $1-o(1)$; on it, fixed-tolerance reference failures have probability at most $\exp\{C(K+L)-cN\}$, with $c$ independent of $K$. The system-size threshold may depend on both depths. This permits the true/reference likelihood loss to be made smaller than $cN/2$ and gives the claimed true-law neighborhood estimate.

\subsection{A fixed terminal cap}\label{sec:terminal}

The compact-time estimates are sufficient because the end of the observation path has a simpler geometry.

\begin{prop}[Terminal convexity and containment]\label{prop:terminal-cap}
For every fixed $\beta>0$ there are $\rho,c>0$ such that, with probability at least $1-2e^{-cN}$ over $J$, the full TAP functional is $1/2$-strongly convex on every cap
\[
 K_\tau(\rho)=\{m\in[-1,1]^N:\ \sum_i(1-\tau_i m_i)\le\rho N\},
 \qquad \tau\in\{-1,1\}^N,
\]
for every external field. Set $S=2\log(8/\rho)$ and $\tau=\sign Y_S$. For every prior on the cube, conditionally on its environment,
\[
 \bbP\{\exists t\ge S:\ m_t\notin K_\tau(\rho)\}
       \le e^{-\rho N/4}.
\]
\end{prop}

The proof is in \S\ref{sec:terminal-proof}. Small principal submatrices of $J$ control the entropy-weighted interaction on each cap, simultaneously over all its sign centers. For containment, $d_H(X,\sign Y_S)$ is binomial after averaging the observation noise; its conditional expectation is a martingale. The exponential maximal inequality therefore gives one cap valid for all later times. Neither assertion requires the AT condition.

\subsection{Simultaneous correctness of the fixed domains}\label{sec:anchor-domains}

\begin{lem}[Distances from a fixed time anchor]\label{sampler:anchors}
For any fixed $s\in[0,S]$,
\[
 \sup_{s\le t\le S}
 \left|\frac{\|m_t-m_s\|_2^2}{N}-\{r(t)-r(s)\}\right|
 \xrightarrow{\bbP}0.
\]
This holds simultaneously for every fixed finite collection of anchors.
\end{lem}
\begin{proof}
The process $Z_t=N^{-1}m_s^{\mathsf T}(m_t-m_s)$, $t\ge s$, is a
martingale. Conditional total variance and Doob's inequality give
\[
 \bbE\sup_{s\le t\le S}|Z_t|^2
 \le\frac4{N^2}\bbE[m_s^{\mathsf T}C_sm_s]
 \le\frac4N\bbE\|C_s\|_{\op}
 \le\frac4N\sqrt{\bbE\Tr(C_s^2)}=o(1).
\]
The last equality uses the qualitative posterior overlap conclusion in
\S\ref{sec:scalar-law}. The uniform empirical law there also gives
$N^{-1}\|m_t\|^2\to r(t)$ uniformly on $[0,S]$. Expand
$\|m_t-m_s\|^2/N=\|m_t\|^2/N-\|m_s\|^2/N-2Z_t$.
A finite union proves the final statement.
\end{proof}

We now justify the fixed record used in \eqref{sampler:anchor-accuracy}--\eqref{sampler:anchor-increments}. At time zero, Lemma~\ref{lem:original-tap-input} gives a positive curvature radius around a sufficiently late original iterate. The fixed-depth empirical comparison in \S\ref{sec:initializer-proof} strengthens the double-limit mean approximation to concentration of the normalized error at each fixed depth. Choose that depth so its limiting error is smaller than a fixed fraction of the curvature radius. This gives a neighborhood of the true $m_0$ with probability tending to one.

Next choose a short deterministic initial interval on which $r(t)-q$ is sufficiently small. Lemma~\ref{sampler:anchors} then keeps all its true means inside the initial neighborhood with probability $1-o(1)$. On the remaining positive compact interval ending at $S$, Theorem~\ref{curv:fixed-anchor} supplies common positive radii and curvature constants. Decrease $a,\gamma$ so that the balls of radius $4a\sqrt N$ fit in these neighborhoods. These constants are chosen before the anchor mesh.

Choose a fixed mesh with $r(s_{j+1})-r(s_j)<a^2/16$. Applying Lemma~\ref{sampler:anchors} at its finitely many anchors gives \eqref{sampler:anchor-increments}. Choose the fixed AMP depths in Theorem~\ref{amp:actual} for normalized error at most $a/4$. The same depths work at every anchor; a finite union gives \eqref{sampler:anchor-accuracy} and \eqref{sampler:anchor-curvature} simultaneously. The clipped ball domains in \eqref{sampler:ball-domain} therefore contain every intervening true mean and retain $\gamma$-strong convexity. Proposition~\ref{prop:terminal-cap} supplies the same conclusions after $S$.

Thus, with probability $1-o(1)$ under the true observation experiment, all the stored domains are correct throughout the path. Each domain is constructed from its current anchor observation only. This simultaneous assertion will be used through the adapted stopping rule in \S\ref{sec:completion}.

\subsection{From objective accuracy to drift accuracy}

\begin{lem}[Objective accuracy is sufficient]\label{sampler:objective}
Suppose $ F_{\tap}(\cdot,f)$ is $\gamma$-strongly convex on a convex domain
$\mathcal K\subseteq[-1,1]^N$, and $m\in\mathcal K\cap(-1,1)^N$.
If $\widehat m\in\mathcal K$ has objective gap at most $\varepsilon_N$,
then
\begin{equation*}
 \|\widehat m-m\|^2
 \le4\gamma^{-2}\|\nabla F_{\tap}(m,f)\|^2
                      +4\gamma^{-1}\varepsilon_N.
\end{equation*}
\end{lem}
\begin{proof}
The objective gap implies $ F_{\tap}(\widehat m,f)\le F_{\tap}(m,f)+\varepsilon_N$.
Strong convexity at $m$ gives, with $d=\|\widehat m-m\|$,
$\gamma d^2/2\le\|\nabla F_{\tap}(m,f)\|d+\varepsilon_N$.
Use $rd\le\gamma d^2/4+r^2/\gamma$.
\end{proof}

Apply this lemma with $f=h\bone+Y_t$, $m=m_t$, and the stored ball or cap. The clipped numerical optimizer of \S\ref{sec:algorithm} has the required objective gap on the original domain; its accuracy and finite operation budget are verified in \S\ref{sec:numerical}. Consequently, before the first current failure of the structural conditions,
\[
 \|\widehat m_t-m_t\|^2
 \le4\gamma^{-2}\|\nabla F_{\tap}(m_t,h\bone+Y_t)\|^2
                    +4\gamma^{-1}N^{-8}.
\]
This is the sole quantitative mean-approximation estimate needed to complete the likelihood argument.


\section{Completion of the proof}\label{sec:completion}

We now prove that the proposal of Algorithm~\ref{alg:gaussian-proposal}
has a polynomial likelihood cutoff under the exact target
\eqref{eq:sign-target}. The structural estimates are required with
probability $1-o(1)$ only. An adapted stopping argument allows us to use
this qualitative probability bound without multiplying it by the
dimension or by the observation horizon.

Throughout this section, probabilities without a conditioning subscript
include the original Gaussian disorder, the true observation experiment,
and the independent retained offset and seeds. The proposal is always
evaluated using its own prescribed deterministic maps, including guards
and fallbacks. To discuss path laws on an input rejected by the global
norm guard, define an auxiliary proposal with zero drift and put the
analysis stopping time equal to zero. The implemented algorithm instead
returns uniform spins on that input. Since
\[
 \mathbb E\|J\|_F^2=\beta^2(N-1),\qquad
 \mathbb P(\|J\|_F>N^2)=O_\beta(N^{-3}),
\]
this convention changes the final average error by $o(1)$.

\subsection{A stopped likelihood estimate}

Write $R_t=\nabla F_{\tap}(m_t,h\mathbf1+Y_t)$. At every ordinary
fine-grid point at which the stored domain contains $m_t$ and is
$\gamma$-strongly convex, the objective guarantee
\eqref{eq:algorithm-objective} and Lemma~\ref{sampler:objective} give
\begin{equation}\label{eq:completion-drift-error}
 \|b_j-m_{t_j}\|^2
    \le4\gamma^{-2}\|R_{t_j}\|^2
                   +4\gamma^{-1}\varepsilon_N.
\end{equation}
The estimate uses the residual at the true mean, although that mean
does not enter the implementation.

For the analysis, let $\tau$ be the first current grid point before
$T_N$ at which a guard fails or, at an ordinary drift-evaluation point,
one of the following properties fails: the current domain contains the
current true mean, the TAP functional is $\gamma$-strongly convex on
that domain, or the computed vector has the prescribed feasible
objective accuracy. Put $\tau=T_N$ if no failure occurs. These tests
depend on $J$, the retained independent randomness, and the observations
up to the current time. Thus $\tau$ is an adapted stopping time in the
observation filtration enlarged by that randomness. In particular, it
does not test whether a stored domain will contain future means.
The algorithm neither computes nor uses $\tau$.

The anchor and terminal-cap conclusions give simultaneous correctness
of all the domains with probability $1-o(1)$. On this event the
prescribed optimization budget gives the required accuracy at every
ordinary grid point. The field guard also has vanishing failure
probability. Indeed, a process with drift in $[-1,1]^N$ satisfies
\[
 \sup_{t\le T_N}|Y_{t,i}|
 \le T_N+\sup_{t\le T_N}|\widetilde B_{t,i}|,
\]
where $\widetilde B$ is its driving Brownian motion; the reflection
bound applies both to the true process and to every proposal.
Consequently,
\begin{equation*}
 b_N:=\mathbb P(\tau<T_N)=o(1).
\end{equation*}
No quantitative rate for $b_N$ is used below.

Lift the proposal to continuous time by keeping its drift constant on
each interval of the augmented grid. Under $P_J$, let
$\widetilde B$ be the innovation Brownian motion, so that
$dY_t=m_t\,dt+d\widetilde B_t$. Define
\[
 V_\tau=\int_0^\tau\|m_t-b_{\lfloor t\rfloor}\|^2\,dt,
 \qquad
 M_\tau=\int_0^\tau(m_t-b_{\lfloor t\rfloor})^{\mathsf T}
                                      d\widetilde B_t,
\]
where $\lfloor t\rfloor$ denotes the preceding grid point.
The posterior-mean martingale and the covariance budget
\eqref{eq:posterior-budget} imply, conditional on any realization of
the grid,
\begin{align*}
 \sum_j\int_{t_j}^{t_{j+1}}
        \mathbb E\|m_t-m_{t_j}\|^2\,dt
 &=\sum_j\int_{t_j}^{t_{j+1}}
        \int_{t_j}^{t}\mathbb E\operatorname{Tr}(C_s^2)\,ds\,dt\\
 &\le\eta\int_0^{T_N}\mathbb E\operatorname{Tr}(C_s^2)\,ds
 \le\eta N.
\end{align*}
The grid is independent of the true observation experiment, so the
same inequality holds after averaging over its offset.

Before $\tau$, split the drift error on an ordinary interval into
$m_t-m_{t_j}$ and $m_{t_j}-b_j$. For the second term use
\eqref{eq:completion-drift-error}. Each outgoing interval has length
at most $\eta$, and discarding its stopping indicator only increases
the resulting nonnegative residual sum. The random-offset inequality
\eqref{sampler:quadrature} and the integrated TAP estimate
\eqref{eq:integrated-tap} therefore give
\[
 \mathbb E\sum_{\substack{j:\,t_j<\tau\\
                         t_j\ \mathrm{ordinary}}}
       \Delta_j\|R_{t_j}\|^2
 \le2\int_0^{T_N}\mathbb E\|R_t\|^2\,dt
 \le C_{\beta,h}.
\]
On each interval immediately following a deterministic anchor the
drift is zero, so its contribution is at most $N\eta$. There are
only $K+1$ such intervals. Including the objective errors yields
\begin{equation*}
 \mathbb E V_\tau
  \le C_{\beta,h,\gamma}
       +C(K+1)N\eta+C_\gamma T_N\varepsilon_N
  \le C_{\beta,h,\gamma,K},
 \qquad \mathbb E M_\tau^2=\mathbb E V_\tau.
\end{equation*}
The last identity is It\^o's isometry. All integrands are bounded on
the finite horizon, so no additional integrability assumption is needed.

Girsanov's theorem, applied to the bounded true and proposal drifts,
gives on $\{\tau=T_N\}$
\[
 \log\frac{dP_J}{dQ_J}=M_\tau+\frac12V_\tau;
\]
see \cite[Chapter~III, Section~5]{karatzas1991brownian}.
The same identity holds for the retained finite-dimensional state.
Indeed, relative to the common zero-drift reference, the true
observation likelihood depends only on the endpoint and the proposal
likelihood is the product of its Gaussian transition densities.
Both are measurable functions of the retained grid observations.
The offset and seed densities are the same and cancel.

By Chebyshev's and Markov's inequalities, for every $b>0$,
\begin{equation*}
 \mathbb P\left\{\log\frac{dP_J}{dQ_J}>b\right\}
  \le b_N+\frac{C}{b^2}+\frac{C}{b},
\end{equation*}
where $C$ depends only on the fixed parameter pair and anchor record.
This is the required likelihood control. The proof estimates the
stopped energy, and therefore never charges the full possible drift
cost on the event of structural failure.

\subsection{Metropolis correction and the returned spin law}

Let $\Pi_J$ be the exact rounding-conditioned target on the augmented
grid, and put
\[
 p_N=\Phi(\sqrt{T_N})^N,
 \qquad w_J=\frac{d\Pi_J}{dQ_J},
 \qquad\delta_N(J)=\Pi_J\{w_J>N\}.
\]
Theorem~\ref{thm:sign-target} gives
$p_N=1-O(N^{-1})$ and $d\Pi_J/dP_J\le p_N^{-1}$.
Hence
\begin{align*}
 \mathbb E_J\delta_N(J)
 &\le p_N^{-1}\,
       \mathbb P\left\{\log\frac{dP_J}{dQ_J}>\log(Np_N)\right\}\\
 &\le p_N^{-1}\left\{b_N+
                 \frac{C}{\log(Np_N)}+
                 \frac{C}{\log^2(Np_N)}\right\}
   =o(1).
\end{align*}
In particular, the target-to-proposal likelihood ratio is at most $N$
outside a set whose target mass tends to zero in probability over the
disorder. This conclusion concerns the exact target used by the
Metropolis step, including all its retained auxiliary variables.

Start the independence Metropolis chain from one draw of $Q_J$ and run
\[
 k_N=\lceil2N\log N\rceil
\]
correction steps. When $\delta_N(J)\le1/2$,
Lemma~\ref{lem:imh} bounds its distance from $\Pi_J$ by
\[
 \delta_N(J)+\frac{\delta_N(J)}N
       +e^{-2(1-\delta_N(J))\log N}
 \le\delta_N(J)+\frac2N.
\]
When $\delta_N(J)>1/2$, the trivial bound by one is at most
$2\delta_N(J)$. Thus, uniformly in $J$, the augmented chain's error
is at most $2\delta_N(J)+2/N$. Total variation cannot increase under
the endpoint-sign map, and the pushforward of $\Pi_J$ is exactly
$\mu_J$. After adding the probability of the input-guard failure,
we obtain for the ideal algorithm
\[
 \mathbb E_J\|\widehat\mu_J-\mu_J\|_{\TV}
  \le2\mathbb E_J\delta_N(J)+\frac2N+O_\beta(N^{-3})=o(1).
\]
The implementation in \S\ref{sec:numerical} adds $o(1)$ to this bound.
Markov's inequality now proves Theorem~\ref{thm:main}.

\begin{breakablealgorithm}
\caption{The complete SK sampler}
\label{alg:complete-sampler}
\begin{algorithmic}[1]
\Require $J=\beta\bW$, fixed $(\beta,h)$ satisfying the strict AT
condition, and the fixed valid anchor record.
\Ensure A spin configuration with conditional law $\widehat\mu_J$.
\If{$N$ is below the fixed cutoff or the input norm guard fails}
  \State \Return independent uniform signs.
\EndIf
\State Draw an augmented path $y$ and its score $\ell=\mathcal L_J(y)$
using Algorithm~\ref{alg:gaussian-proposal}.
\For{$r=1,\ldots,\lceil2N\log N\rceil$}
  \State Independently draw $y'$ and
  $\ell'=\mathcal L_J(y')$ using Algorithm~\ref{alg:gaussian-proposal},
  including a fresh retained grid offset.
  \State Draw an independent $V\sim\operatorname{Unif}(0,1)$.
  \If{$\log V\le\min\{0,\ell'-\ell\}$}
    \State $(y,\ell)\gets(y',\ell')$, retaining the proposal's own
    offset, grid, and seeds.
  \EndIf
\EndFor
\State \Return $\sign y_{T_N}$.
\end{algorithmic}
\end{breakablealgorithm}

The correction uses only the two log weights. Its state is an entire
generated proposal, so keeping the current score avoids any need to
recompute it. If a score is recomputed to higher precision, the
deterministic maps that defined the proposal remain fixed. In
particular, refinement of a score does not change the optimization
tolerance or the underlying proposal distribution.

\subsection{Computational cost and the role of the estimates}

Each path has $O(N^4\log N)$ intervals and a fixed number of anchors.
The AMP depths are fixed. Under the guards,
$L_N=O_{\beta,h,a,\rho}(N^{30})$ and
$n_{\mathrm{opt}}=O_{\beta,h,a,\rho,\gamma}(N^{30}\log N)$.
Each projected-gradient iteration uses a dense $O(N^2)$ gradient
evaluation and a one-dimensional projection search; the latter has
$O(N\log N)$ arithmetic cost at any prescribed inverse-polynomial
accuracy. Together with the $O(N\log N)$ independent proposals,
these choices give the loose arithmetic bound
\[
 O_{\beta,h}(N^{37}(\log N)^3).
\]
The bound is for the displayed conservative implementation and is not
an optimized complexity claim. The numerical appendix establishes a
polynomial bit-operation bound, without identifying it with this
arithmetic exponent. Guards, fixed iteration counts, and bounded-cost
random approximations make these bounds deterministic on every input.

Only the bounded interval $[0,S]$ requires the scalar posterior law,
AMP initialization, and the strict AT curvature estimate. After $S$,
one fixed convex cap supplies the required mean approximation through
$T_N=4\log N$. The integrated residual estimate and the covariance
budget control its accumulated drift cost uniformly in this horizon.
The likelihood cutoff then converts that bounded cost, together with
qualitative structural success, into vanishing total-variation error.
No quantitative rate for the observed-path overlap convergence or structural failure is required for the conclusion of Theorem~\ref{thm:main}; the original-model input remains the quantitative estimate of Lemma~\ref{lem:original-overlap}.


\appendix

\section{Gaussian response and the integrated TAP residual}\label{sec:integrated-proof}

We prove Theorem~\ref{thm:integrated-tap}. The argument uses only one-row Gaussian response estimates and the posterior budget in Lemma~\ref{lem:filtering-input}; it requires no overlap concentration or AT assumption.

\subsection{Response to a Gaussian cavity row}

\begin{lem}[Gaussian response]\label{lem:response}
Fix $B<\infty$. Let $\nu$ be a law on $\{-1,1\}^r$, with mean $m$ and
covariance $C$. For $z\in\bbR^r$, let
\[
 \nu_z(\sigma)=\frac{e^{z\cdot\sigma}\nu(\sigma)}{
                    \bbE_\nu e^{z\cdot\sigma}},
 \qquad m_z=\bbE_{\nu_z}\sigma,\qquad C_z=\Cov_{\nu_z}(\sigma).
\]
Suppose $vr\le B$, $g\sim N(0,vI_r)$, and
$\|\delta\|_2\le v\sqrt r$. Set $z=g+\delta$. There is a constant $C_B$
such that
\begin{align*}
 \bbE_g\|m_z-m\|_2^2&\le C_Bv\Tr(C^2),\\
 \bbE_g\Tr(C_z^2)&\le C_B\Tr(C^2),\\
 \bbE_g\left|\log\bbE_\nu e^{z\cdot\sigma}
       -z\cdot m-\frac v2\Tr C\right|^2
       &\le C_Bv^2\Tr(C^2),\\
 \bbE_g\left|z\cdot(m_z-m)-v\Tr C\right|^2
       &\le C_Bv^2\Tr(C^2).
\end{align*}
In particular, for
$L_\nu(z)=\frac12\log(\bbE_\nu e^{z\cdot\sigma}/\bbE_\nu e^{-z\cdot\sigma})$,
\[
 \bbE_g|L_\nu(z)-z\cdot m|^2\le C_Bv^2\Tr(C^2).
\]
\end{lem}

\begin{proof}
Write $\chi=\sigma-m$, let $\chi'$ be an independent copy under $\nu$,
and put $u=\chi\cdot\chi'$. Jensen's inequality gives
$D(z):=\bbE_\nu e^{z\cdot\chi}\ge1$. The identity
\[
 \|\chi\|_2^2=\Tr C-2m\cdot\chi
\]
implies
\[
 w(\chi):=\exp\left(\frac v2\|\chi\|_2^2+\delta\cdot\chi\right)
   =c\exp(a\cdot\chi),\quad
 c=e^{v\Tr C/2},\quad a=\delta-vm.
\]
Here $1\le c\le e^{B/2}$, $\|a\|_2\le2v\sqrt r$, and
$|a\cdot\chi|,|vu|$ are bounded by constants depending only on $B$.
Gaussian pair integration gives
\[
 \bbE_g\|m_z-m\|_2^2
 \le \bbE_{\chi,\chi'}u\,w(\chi)w(\chi')e^{vu}.
\]
This follows by first removing the denominator $D(z)^2\ge1$ from the
nonnegative squared norm. Replacing $e^{vu}$ by $1$ leaves
$\|\bbE_\nu[\chi w(\chi)]\|_2^2$, and the error is at most
$C_Bv\bbE u^2=C_Bv\Tr(C^2)$.
For each unit vector $b$, centering and Cauchy--Schwarz show that
\[
 |\bbE(b\cdot\chi)w(\chi)|
 \le C_B(b^{\mathsf T}Cb)^{1/2}(a^{\mathsf T}Ca)^{1/2}.
\]
Consequently
\begin{equation}\label{eq:weighted-vector}
 \|\bbE\chi w(\chi)\|_2^2
 \le C_Bv^2r\|C\|_{\op}^2
 \le C_Bv\Tr(C^2).
\end{equation}
This proves the first assertion.

For the second, $C_z\preceq M_z:=\bbE_{\nu_z}\chi\chi^{\mathsf T}$, so
$\Tr(C_z^2)\le\Tr(M_z^2)$. Removing the positive denominator and
integrating the Gaussian vector bounds the expectation by
\[
 \bbE_{\chi,\chi'}u^2w(\chi)w(\chi')e^{vu}
 \le C_B\Tr(C^2).
\]

To prove the log-partition estimate, note that
$\bbE_gD(z)=c\bbE_\nu e^{a\cdot\chi}$. Since $\bbE\chi=0$ and
$a\cdot\chi$ is bounded,
\[
 |\bbE_gD(z)-c|\le C_Ba^{\mathsf T}Ca
 \le C_Bv^2r\|C\|_{\op}.
\]
The square of this bound is at most $C_Bv^2\Tr(C^2)$.
Moreover,
\begin{align*}
 \Var_g D(z)
 &=\bbE_{\chi,\chi'}w(\chi)w(\chi')(e^{vu}-1)\\
 &\le v\|\bbE\chi w(\chi)\|_2^2+C_Bv^2\Tr(C^2)
 \le C_Bv^2\Tr(C^2),
\end{align*}
where the stronger first bound in \eqref{eq:weighted-vector} and $vr\le B$
are used in the last step. Thus $\bbE|D(z)-c|^2\le C_Bv^2\Tr(C^2)$.
Both $D(z)$ and $c$ are at least $1$, and the logarithm is
$1$-Lipschitz on $[1,\infty)$. This proves the third assertion. Applying
it to $z$ and $-z$ proves the assertion about $L_\nu$.

Finally, let $F(g)=m_{g+\delta}-m$. Since $\nabla_gF=C_z$, Gaussian
integration by parts gives the divergence isometry
\[
 \bbE_g[g\cdot F-v\operatorname{div}F]^2
 =v\bbE_g\|F\|_2^2+v^2\bbE_g\Tr[(\nabla_gF)^2]
 \le C_Bv^2\Tr(C^2).
\]
The trace identity
$\Tr C_z-\Tr C=-2m\cdot F-\|F\|_2^2$, together with
$\|F\|_2\le2\sqrt r$, shows that
\[
 \bbE|v(\Tr C_z-\Tr C)|^2
   \le C v^2r\bbE\|F\|_2^2\le C_Bv^2\Tr(C^2).
\]
The same bound holds for $\bbE|\delta\cdot F|^2$.
Adding these two errors to the divergence estimate proves the final
assertion. Finite support and Gaussian integrability justify every
differentiation and integration by parts.
\end{proof}

\begin{lem}[One-row TAP residual]\label{lem:row-residual}
Under the assumptions of Lemma~\ref{lem:response}, fix any $a\in\bbR$ and
consider the law
\[
 \widehat\nu_{z,a}(x,\sigma)\propto
 e^{ax+xz\cdot\sigma}\nu(\sigma),\qquad x\in\{-1,1\}.
\]
Write $\widehat m_i=\bbE x$ and $\widehat m_R=\bbE\sigma$. Uniformly in $a$,
\[
 \bbE_g\left|
 \atanh\widehat m_i-a-z\cdot\widehat m_R
       +v(r-\|\widehat m_R\|_2^2)\widehat m_i
 \right|^2\le C_Bv^2\Tr(C^2).
\]
\end{lem}

\begin{proof}
With $p=(1+\widehat m_i)/2$, the exact identities are
\[
 \atanh\widehat m_i=a+L_\nu(z),\qquad
 \widehat m_R=p m_z+(1-p)m_{-z}.
\]
Lemma~\ref{lem:response} implies that
$z\cdot(m_z-m)=v\Tr C+e_+$ and
$z\cdot(m_{-z}-m)=-v\Tr C+e_-$, where
$\bbE e_+^2+\bbE e_-^2\le C_Bv^2\Tr(C^2)$.
Therefore
\[
 z\cdot\widehat m_R=z\cdot m+v\Tr C\widehat m_i
                          +p e_++(1-p)e_-.
\]
The dependence of $p$ on $g$ causes no difficulty: $0\le p\le1$.
Also $\bbE\|\widehat m_R-m\|_2^2\le C_Bv\Tr(C^2)$, so
\[
 \bbE\left|v(\|\widehat m_R\|_2^2-\|m\|_2^2)\right|^2
 \le C v^2r\bbE\|\widehat m_R-m\|_2^2
 \le C_Bv^2\Tr(C^2).
\]
Combine these estimates with the odd log-partition estimate.
\end{proof}


\subsection{Proof of Theorem~\ref{thm:integrated-tap}}

\begin{proof}
Fix a site $i$ and condition on the interactions not incident to it.
Let $\mu^{-i}$ be the cavity law on the remaining $r=N-1$ spins and let
$g=(J_{ij})_{j\ne i}\sim N(0,vI_r)$, $v=\beta^2/N$.
Consider the following reference experiment, in which the row is independent
of all observations:
\[
 X_R^c\sim\mu^{-i},\qquad
 X_i^c\sim p_h(x):=\frac{e^{hx}}{2\cosh h},\qquad
 Y_R=tX_R^c+\sqrt t Z_R,
 \quad Y_i=tX_i^c+\sqrt t Z_i.
\]
The cavity spin, one-site spin, row $g$, and Gaussian noises are independent
except for the displayed observations. Conditional on the cavity disorder
and $Y$, put $\nu=\mu^{-i}_{Y_R}$ and $a=h+Y_i$.

Summing over the spin at $i$ shows that the density of the true law of
$(g,Y)$ relative to this reference law is exactly
\begin{equation*}
 R(g,Y)=
 \frac{\bbE_\nu\cosh(a+g\cdot\sigma)/\cosh a}
      {\bbE_{\mu^{-i}}\cosh(h+g\cdot\sigma)/\cosh h}.
\end{equation*}
The denominator is at least $1/\cosh h$. For any nonnegative row function
$F$, Gaussian exponential tilting therefore yields
\begin{align*}
 \bbE_g[R(g,Y)F(g)]
 &\le \cosh h\,
     \bbE_{\xi\sim p_a,\sigma^0\sim\nu}
             \bbE_g[e^{\xi g\cdot\sigma^0}F(g)]\\
 &=\cosh h\,e^{vr/2}
     \bbE_{\xi\sim p_a,\sigma^0\sim\nu}
             \bbE_gF(g+v\xi\sigma^0).
\end{align*}
For the squared residual in Lemma~\ref{lem:row-residual}, the shift has
norm $v\sqrt r$. Hence the true one-row residual has mean square at most
$C_{\beta,h}v^2\bbE\Tr[(C_t^{c,i})^2]$, where $C_t^{c,i}$ is the cavity
posterior covariance in the reference experiment. The estimate is uniform
in $t$ and $Y_i$; the probabilities $p_a$ sum to one. In particular,
there is no factor that grows like $e^t$.

Apply the covariance budget \eqref{eq:posterior-budget} to each fixed
cavity law:
\[
 \int_0^\infty\bbE\Tr[(C_t^{c,i})^2]\,\de t\le N-1.
\]
Summing over $i$ bounds the integrated squared residual with the
site-deleted Onsager coefficient
$v\sum_{j\ne i}(1-m_{t,j}^2)$ by
$C_{\beta,h}Nv^2(N-1)\le C_{\beta,h}\beta^4$.
The conventional TAP coefficient includes the $j=i$ term. Its additional
residual is $v(1-m_{t,i}^2)m_{t,i}$, and
\[
 \int_0^\infty\bbE\sum_i[v(1-m_{t,i}^2)m_{t,i}]^2\,\de t
 \le v^2\int_0^\infty\bbE\Tr(C_t^2)\,\de t\le v^2N.
\]
This proves \eqref{eq:integrated-tap}.
\end{proof}


\section{The scalar law under the actual observation experiment}\label{sec:scalar-law}

We prove Proposition~\ref{prop:scalar-law}. The initial overlap estimate is Lemma~\ref{lem:original-overlap}. The scalar parameters are defined in \eqref{eq:scalar-observation-law}--\eqref{eq:observed-fixed-point}, and their strict contraction is Lemma~\ref{lem:scalar-stability}. Every cavity model retains entry variance $\beta^2/N$.

\subsection{Proof of the posterior scalar and overlap laws}

\begin{proof}
We give the cavity estimates, the scalar contraction, and the upgrade from
integrated to uniform qualitative concentration separately.

We use Lemma~\ref{lem:response}, including the direct estimate
$\bbE_g|D(g)-\exp(v\Tr C/2)|^2\le C_Bv^2\Tr(C^2)$ proved there.
The lemma also allows shifts of norm at most $K v\sqrt r$ for any fixed
$K$, with a constant depending on $K$: the proof is unchanged.

\medskip\noindent\emph{The exact one-row observation density.}
Delete vertex $i$ and condition on the remaining disorder. Write
$\mu^c=\mu^{-i}$ and $g=(J_{ij})_{j\ne i}$, so $v=\beta^2/N$ and $r=N-1$.
Under a reference law, independently draw $g$, a cavity spin $X_R^c\sim\mu^c$,
and a one-site spin $X_i^c\sim p_h$, where
$p_h(x)=e^{hx}/(2\cosh h)$; observe each through its independent Gaussian
channel of time $t$.
For fixed observations put $\nu=\mu^c_{Y_R}$, $a=h+Y_i$, and denote the
means and covariances of $\mu^c,\nu$ by $m^0,C^0,m^t,C^t$.
The true joint law of $(g,Y_R,Y_i)$ divided by this reference law is
\[
 R(g,Y)=\frac{A_t(g,a)}{A_0(g,h)},\qquad
 A_t(g,a)=\frac{\mathbb E_\nu\cosh(a+g\cdot\sigma)}{\cosh a},\quad
 A_0(g,h)=\frac{\mathbb E_{\mu^c}\cosh(h+g\cdot\sigma)}{\cosh h}.
\]
This follows directly by writing the two partition functions after summing
out vertex $i$. Notice $A_0\ge1/\cosh h$.
The true one-site mean and the remaining means satisfy exactly
\[
 m_{t,i}=\tanh(a+L_\nu(g)),\qquad
 m_{t,R}=p m_g+(1-p)m_{-g},\quad p=(1+m_{t,i})/2.
\]
The vector response estimate remains valid under the row density:
\[
 \mathbb E_g R(g,Y)\|m_{t,R}-m^t\|^2
 \le C_{\beta,h}v\operatorname{Tr}[(C^t)^2].
\]
To see this, use $R\le\cosh h A_t$, expand the numerator into
$\xi\sim p_a,\sigma^0\sim\nu$, and use
$\mathbb E_g[e^{\xi g\cdot\sigma^0}F(g)]
=e^{vr/2}\mathbb E_gF(g+v\xi\sigma^0)$.
The shifted vector estimate applies with $K=1$ to each of $m_g,m_{-g}$;
only $0\le p\le1$ is needed.
Also $\mathbb E_gR(g,Y)\le C_{\beta,h}$ uniformly in $Y,t$.

Set
\[
 b=N^{-1}\|m^0\|^2,\quad c=N^{-1}m^0\cdot m^t,\quad
 s=N^{-1}\|m^t\|^2.
\]
The preceding response bound yields, conditionally on cavity data,
\[
 \mathbb E_g R\,|R_N(t)-s|
 \le \frac{C}{N}\left(1+\sqrt{\operatorname{Tr}[(C^t)^2]}\right),
\]
and
\[
 \mathbb E_g R\,|R_N(t)-s|^2
 \le \frac{C}{N^2}\left(1+\operatorname{Tr}[(C^t)^2]\right).
\]
Indeed, $|R_N-s|\le N^{-1}+2\sqrt r\,\|m_{t,R}-m^t\|/N$.

\medskip\noindent\emph{Base and mixed overlaps.}
Under the reference law,
\[
 \mathbb E(b-q)^2+\mathbb E(c-q)^2\le C/N.
\]
The bound for $b$ is Jensen applied to the original cavity overlap.
For $c$, draw an independent original cavity replica $S^0$ and a replica
$S^t$ from the posterior given $Y_R$. After averaging $Y_R$, the pair
$(S^0,S^t)$ is distributed as two independent original cavity replicas.
Conditional expectation of their overlap is $c$, so Jensen again applies.
These statements are uniform in $t$.
As $c^2\le bs$, and $q>0$, one has
\[
 (q-s)_+\le C_q(|b-q|+|c-q|).
\]
For example, on $|b-q|+|c-q|\le q/2$, use $s\ge c^2/b$ and the bounded
derivatives of $c^2/b$; on the complement the right side bounds $1$ after
enlarging the constant.

\medskip\noindent\emph{Approximating the row density in $L^1$.}
Define
\[
 \widetilde A_t(g,a)=e^{v\operatorname{Tr}C^t/2}
              \frac{\cosh(a+g\cdot m^t)}{\cosh a},\qquad
 \widetilde A_0(g,h)=e^{v\operatorname{Tr}C^0/2}
              \frac{\cosh(h+g\cdot m^0)}{\cosh h},
\]
and $\widetilde R=\widetilde A_t/\widetilde A_0$.
Both $A_0$ and $\widetilde A_0$ are at least $1/\cosh h$.
For fixed $a$,
\[
 A_t-\widetilde A_t
 =\sum_{\xi=\pm1}p_a(\xi)e^{\xi g\cdot m^t}
       [D_\nu(\xi g)-e^{v\operatorname{Tr}C^t/2}].
\]
Cauchy--Schwarz, the Gaussian $D$ estimate, and the bounded Gaussian
exponential moments of $g\cdot m^t$ imply
$\mathbb E_g|A_t-\widetilde A_t|
\le Cv\sqrt{\operatorname{Tr}[(C^t)^2]}$.
The identical calculation, with the additional bounded exponential factor
from $\widetilde A_t$, gives
\[
 \mathbb E_g\widetilde A_t|A_0-\widetilde A_0|
 \le Cv\sqrt{\operatorname{Tr}[(C^0)^2]}.
\]
Here one expands both normalized cosh factors, then applies Cauchy--Schwarz
to $e^{g\cdot(\xi m^t+\zeta m^0)}$ and the $D$ error; its exponential
second moment is bounded independently of $a,t$.
It follows that
\[
 \mathbb E_g|R-\widetilde R|
 \le Cv\left(\sqrt{\operatorname{Tr}[(C^0)^2]}
             +\sqrt{\operatorname{Tr}[(C^t)^2]}\right).
\]
Moreover, the odd-log-partition estimate and
$\mathbb E_g A_t^2\le C$ give
\[
 \mathbb E_g R\left|m_{t,i}^2-\tanh^2(a+g\cdot m^t)\right|
 \le Cv\sqrt{\operatorname{Tr}[(C^t)^2]}.
\]
The same conclusion holds for a bounded Lipschitz test function of the
one-site mean.

\medskip\noindent\emph{Identification of the scalar experiment.}
We have
\[
 \widetilde R=e^{-\beta^2(s-b)/2}
 \frac{\cosh(a+g\cdot m^t)}{\cosh a}
 \frac{\cosh h}{\cosh(h+g\cdot m^0)}.
\]
The pair $(g\cdot m^0,g\cdot m^t)$ is centered Gaussian with covariance
$\beta^2\bigl(\begin{smallmatrix}b&c\\c&s\end{smallmatrix}\bigr)$.
Compare it, by linear interpolation of the covariance matrices, with the
positive semidefinite matrix
$\beta^2\bigl(\begin{smallmatrix}q&q\\q&s\vee q\end{smallmatrix}\bigr)$,
and replace the prefactor by $e^{-\beta^2((s\vee q)-q)/2}$.
For any bounded $C^2$ test function of the mean, Gaussian covariance
differentiation bounds the resulting change in its unnormalized weighted
expectation by
\[
 C\bigl(|b-q|+|c-q|+(q-s)_+\bigr).
\]
For completeness, write $x_0=g\cdot m^0$ and
$x_t=g\cdot m^t$. For a bounded $C^2$ test function $\varphi$, every
second partial derivative, in $(x_0,x_t)$, of
\[
 \frac{\cosh(a+x_t)}{\cosh a}
 \frac{\cosh h}{\cosh(h+x_0)}\,
 \varphi\bigl(\tanh(a+x_t)\bigr)
\]
is bounded in absolute value by $C_{h,\varphi}e^{|x_t|}$, uniformly in
$a,x_0$. This follows from
$\cosh(a+x_t)/\cosh a\le e^{|x_t|}$ and the bounded first two
derivatives of $1/\cosh(h+x_0)$.
The variance of $x_t$ is at most $\beta^2$ throughout the covariance
interpolation, so the derivative bound is uniformly integrable.
The change of prefactor satisfies the same bound.

In the matched covariance case write
$U=h+\beta\sqrt q Z_0$ and
$V=\beta\sqrt{(s-q)_+}Z_1$, independently.
The reference singleton observation has density
\[
 e^{-t/2}\frac{\cosh(h+y)}{\cosh h}\,\gamma_t(dy).
\]
After multiplication by the matched row density, the joint factor becomes
\[
 e^{-[t+\beta^2(s-q)_+]/2}
 \frac{\cosh(U+V+y)}{\cosh U}\,\gamma_t(dy).
\]
Integration over $V+y$, a centered Gaussian of variance
$t+\beta^2(s-q)_+$, is precisely the scalar Gaussian-observation kernel
from initial field $U$. Its total mass is $1$, and its squared posterior
mean has expectation
\[
 f_t(s)=a\bigl(t+\beta^2(s-q)_+\bigr).
\]

Combining the last estimates gives the useful tested identity: for any
bounded function $\psi$ of cavity disorder and $Y_R$,
\[
 \left|\mathbb E_{\rm full}\psi\,[m_{t,i}^2-f_t(s)]\right|
 \le C\|\psi\|_\infty
 \left(\frac1{\sqrt N}+
       \frac1N\mathbb E_{\rm cav}\sqrt{\operatorname{Tr}[(C^t)^2]}\right).
\]
The expectation on the left is under the full quenched observation law;
the right side uses the independent cavity observation experiment.
This is legitimate because the density $R$ is retained in both terms
before Gaussian integration. No independence of the full row and the
full observation is asserted.

\medskip\noindent\emph{Scalar contraction and the covariance budget.}
Use $\psi=s-r(t)$, average the tested identity over $i$, and use the row
removal estimate to replace $s$ by $R_N(t)$ in bounded quadratic terms.
Since $f_t(r(t))=r(t)$ and $f_t$ is $\kappa$-Lipschitz on $[0,1]$,
\[
 \begin{split}
 \mathbb E(R_N-r(t))^2
 &=\frac1N\sum_i\mathbb E(R_N-r(t))(m_{t,i}^2-r(t))\\
 &\le\kappa\,\mathbb E(R_N-r(t))^2+
 C\left(\frac1{\sqrt N}+
 \frac1{N^2}\sum_i\mathbb E_{\rm cav,i}
                    \sqrt{\operatorname{Tr}[(C_i^t)^2]}\right).
 \end{split}
\]
For each fixed cavity disorder its posterior mean is a bounded martingale,
and the It\^o isometry gives
\[
 \int_0^\infty\mathbb E_{\rm cav,i}\operatorname{Tr}[(C_i^t)^2]dt
 \le N-1.
\]
Cauchy--Schwarz in time therefore proves the claimed integrated bound.

\medskip\noindent\emph{General empirical tests.}
For $\varphi\in C^2([-1,1])$, put
\[
 b_{\varphi,t}(s)=\mathbb E\varphi\!
   \left(\tanh(U+w\xi+\sqrt w Z')\right),\qquad
 w=t+\beta^2(s-q)_+.
\]
The scalar posterior mean is a martingale with diffusion coefficient
$1-M^2$, so $b_{\varphi,t}$ is Lipschitz in $s$ with constant at most
$\beta^2\|\varphi''\|_\infty/2$.
The preceding one-site identity applies to $\varphi(m_{t,i})$ and
$b_{\varphi,t}(s)$ with the same error bound up to constants.
Let $V_N(t)=N^{-1}\sum_i\varphi(m_{t,i})$ and
$V_i^c(t)=N^{-1}\sum_{j\ne i}\varphi(m_{j}^t)$.
The row response estimate gives
\[
 \mathbb E_g R|V_N-V_i^c|
 \le\frac C N\left(1+\sqrt{\operatorname{Tr}[(C_i^t)^2]}\right).
\]
Test the one-site identity against
$V_i^c-b_{\varphi,t}(r(t))$ and average over $i$.
Cauchy--Schwarz and $2ab\le a^2/2+2b^2$ absorb one half of
$\mathbb E(V_N-b_{\varphi,t}(r(t)))^2$, leaving a constant multiple of
the scalar squared error and the preceding row errors. Hence
\[
 \int_0^T\mathbb E\left[\bigl(V_N(t)-b_{\varphi,t}(r(t))\bigr)^2\right]dt
 \le C_{\beta,h,\varphi}\left(\frac T{\sqrt N}+\sqrt{\frac TN}\right).
\]

\medskip\noindent\emph{From integrated to uniform qualitative control.}
Put $A_N(t)=\mathbb E R_N(t)$.
The posterior martingale identities imply
\[
 A_N'(t)=N^{-1}\mathbb E\operatorname{Tr}(C_t^2),\qquad
 N^{-1}\mathbb E\|m_t-m_s\|^2=A_N(t)-A_N(s),\quad s\le t.
\]
Thus $A_N$ is nondecreasing. Its integrated convergence to the continuous
function $r$ implies uniform convergence on each compact interval (apply the integrated bound on a slightly larger interval); at
$t=0$ use the original overlap bound.  For $R_N$ itself,
\[
 \mathbb E|R_N(t)-R_N(s)|^2\le4[A_N(t)-A_N(s)].
\]
For each $t$ choose a time $s$ within a short preceding interval where the
integrated squared error is small. The last inequality and uniform
convergence of $A_N$ show
$\sup_{t\le T}\mathbb E(R_N(t)-r(t))^2\to0$.
For $t$ near zero take $s=0$. The same argument applies to $V_N$, using
$\mathbb E|V_N(t)-V_N(s)|^2
\le\|\varphi'\|_\infty^2[A_N(t)-A_N(s)]$.

Let $\chi_N(t)=\mathbb E\operatorname{Tr}(C_t^2)$ and let $T_t$ be the
third centered posterior moment tensor. The covariance filtering SDE is
$d(C_t)_{ij}=\sum_k(T_t)_{ijk}d\widetilde B_{t,k}-(C_t^2)_{ij}dt$.
Consequently
\[
 \chi_N'(t)=\mathbb E\|T_t\|_{\rm F}^2
               -2\mathbb E\operatorname{Tr}(C_t^3)
 \ge-2N\chi_N(t),
\]
because $\|C_t\|_{\rm op}\le\operatorname{Tr}C_t\le N$.
Integrating this differential inequality over $[t,t+1/N]$ gives
\[
 A_N(t+1/N)-A_N(t)
 \ge\frac{1-e^{-2}}2\frac{\chi_N(t)}{N^2}.
\]
Uniform convergence of $A_N$ to continuous $r$ yields
$\sup_{t\le T}\chi_N(t)/N^2\to0$.
For two posterior replicas,
\[
 \mathbb E\left[\left.(R_{12}-R_N(t))^2\right|J,Y_t\right]
 =\frac{\operatorname{Tr}(C_t^2)+2m_t^{\mathsf T}C_tm_t}{N^2}.
\]
The expectation of the right side is at most
$\chi_N(t)/N^2+2\sqrt{\chi_N(t)}/N$, which tends uniformly to zero.
This proves the replica-overlap assertion.

Finally, for $s\le t$, posterior averaging and Jensen give
\[
 \mathbb E\left[\left(\frac{m_s\cdot m_t}{N}-r(s)\right)^2\right]
 \le\mathbb E\langle(R_{12}-r(s))^2\rangle_{Y_s}.
\]
The right side tends uniformly to zero. Expanding the squared distance
therefore shows
\[
 \frac1{N^2}\mathbb E\|m_t-m_s\|^4
 =(r(t)-r(s))^2+o_N(1)
\]
uniformly for $s,t$ in a fixed compact interval.
Apply the Hilbert-space $L^4$ Doob inequality to $m_u-m_s$ on each interval
of a fixed time grid. Since $r$ is Lipschitz, the sum of the resulting
fourth-moment bounds over a grid of mesh $\delta$ is $O(T\delta)+o_N(1)$.
Letting $N\to\infty$ and then $\delta\downarrow0$ gives a modulus of
continuity in probability for $m_t/\sqrt N$.
The empirical averages are Lipschitz functions of this normalized vector.
Their uniform-in-time convergence in probability now follows from the
fixed-time $L^2$ convergence on the grid and continuity of the scalar limit.
\end{proof}


\section{Cavity initialization and AMP under the observation law}\label{sec:initializer-proof}

We use Proposition~\ref{prop:scalar-law}.  Throughout, $\beta,h$ are fixed and satisfy
strict AT, $q$ is their original fixed point, and
$\alpha=\beta^2\mathbb E\operatorname{sech}^4(h+\beta\sqrt q Z)<1$.
Let $\kappa<1$ be the uniform scalar observation stability constant.
All cavity models use the original entry variance $\beta^2/N$.
Deleting finitely many sites changes the corresponding inverse temperature
by $O(N^{-1})$, so the preceding original and observed overlap results
apply uniformly to these models.

\subsection{The explicit finite-depth algorithm}

For a deleted set $S$ and $i\notin S$, define the original-field cavity
tree by
\[
 B^0_{S,i}=\sqrt q,\qquad
 B^{\ell+1}_{S,i}=
 \tanh\left(h+\sum_{j\notin S\cup\{i\}}J_{ij}B^\ell_{S\cup\{i\},j}\right).
\]
Choose fixed nonnegative integers $L,k$.  Given an observation vector $y$,
define
\[
 M^{0,L}_{S,i}(y)=B^L_{S,i},\qquad
 M^{j+1,L}_{S,i}(y)=
 \tanh\left(h+y_i+\sum_{a\notin S\cup\{i\}}
                 J_{ia}M^{j,L}_{S\cup\{i\},a}(y)\right).
\]
Only coordinates outside the deleted set are used by each recursive call.
The output is $M^{k,L}(y)=M^{k,L}_{\varnothing}(y)$.
Every branch is self-avoiding.  Direct evaluation or memoization takes
$N^{O(L+k)}$ arithmetic operations.  In particular, the algorithm is
polynomial-time for fixed $L,k$; no assertion is made for depths growing
with $N$.

\begin{thm}\label{thm:initializer}
For every fixed $T<\infty$ and $\eta>0$, there are fixed integers $L,k$
depending only on $\beta,h,\eta$ such that
\[
 \mathbb P\left\{
 \sup_{0\le t\le T}\frac{\|M^{k,L}(Y_t)-m_t\|_2}{\sqrt N}\le\eta
 \right\}\longrightarrow1.
\]
Thus the actual quenched posterior mean has a polynomial-time coarse
initializer, uniformly over every fixed observation interval.
This theorem does not assert local positivity of the TAP Hessian at that
mean or at the initializer.
\end{thm}

\subsection{A finite-step state evolution}

The state evolution uses three scalar coordinates.  For any bounded
initializer or iterate $M$, these are the limits of
\[
 p_N=N^{-1}\|M\|^2,\qquad
 c_N=N^{-1}m_t\cdot M,\qquad
 b_N=N^{-1}m_0\cdot M.
\]
Let $U=h+\beta\sqrt q Z_0$, draw $\xi$ conditionally with law
$p_U(x)=e^{Ux}/(2\cosh U)$, and let $Z$ be independent standard Gaussian.
Here feasibility means that
$\bigl(\begin{smallmatrix}q&q&b\\q&r(t)&c\\b&c&p\end{smallmatrix}\bigr)$
is positive semidefinite and $0\le p\le1$.
For a feasible triple $\theta=(p,c,b)$ define
\[
 V_{t,\theta}
 =h+\frac bq(U-h)+\bigl[t+\beta^2(c-b)\bigr]\xi
       +\sqrt{t+\beta^2(p-b^2/q)}\,Z,
\]
and
\[
 \Gamma_t(p,c,b)=
 \left(
 \mathbb E\tanh^2V_{t,\theta},\quad
 \mathbb E\xi\tanh V_{t,\theta},\quad
 \mathbb E\tanh U\tanh V_{t,\theta}
 \right).
\]
The square root is defined because covariance feasibility gives
$p\ge b^2/q$.  The map is continuous, uniformly on compact feasible
sets and compact time intervals; it is also locally Lipschitz there by
Gaussian heat differentiation and the bounded derivatives of $\tanh$.

Define $\rho_0=\sqrt q\,\mathbb E\tanh U$ and
\[
 \rho_{\ell+1}=\Delta(\rho_\ell),\qquad
 \Delta(a)=\mathbb E\tanh U_1\tanh U_2,
\]
where $(U_1,U_2)$ is Gaussian, both means are $h$, both variances are
$\beta^2q$, and their covariance is $\beta^2a$.
For fixed integers $R>L$, the joint cavity state evolution of
\cite[Theorem~1]{chen2021convergence}, with the deterministic index
shift matching the present initialization, gives
\[
 \frac{\|B^L\|^2}{N}\longrightarrow q,\qquad
 \frac{(B^L)^{\mathsf T}B^R}{N}\longrightarrow\rho_L
 \quad\text{in }L^2.
\]
The covariance recursion in that theorem identifies the second limit
by induction on $L$; boundedness upgrades convergence in probability
to $L^2$. Moreover,
\cite[Theorems~2 and~3]{chen2021convergence} and
Lemma~\ref{lem:original-overlap} give
\[
 \lim_{R\to\infty}\limsup_{N\to\infty}
       \frac1N\bbE\|B^R-m_0\|^2=0.
\]
Since $\|B^L\|\le\sqrt N$, Cauchy--Schwarz yields
\[
 \bbE\left|\frac{(B^L)^{\mathsf T}(m_0-B^R)}N\right|^2
 \le \frac1N\bbE\|m_0-B^R\|^2.
\]
Taking $N\to\infty$ and then $R\to\infty$ proves
$N^{-1}(B^L)^{\mathsf T}m_0\to\rho_L$ in $L^2$.
We use these original-model inputs directly; the argument below
establishes their observed-law extension.
Gaussian covariance differentiation gives
\[
 |\Delta'(a)|
 =\beta^2\mathbb E[\operatorname{sech}^2U_1\operatorname{sech}^2U_2]
 \le\alpha<1,\qquad -q<a<q.
\]
Since $\Delta(q)=q$ and $0\le\rho_0\le q$,
$0\le q-\rho_L\le q\alpha^L$.

Put $\theta_{0,L}=(q,\rho_L,\rho_L)$ and recursively
$\theta_{j+1,L}(t)=\Gamma_t(\theta_{j,L}(t))$.
For every fixed $L,j,T$, the empirical triple associated to
$M^{j,L}(Y_t)$ converges to $\theta_{j,L}(t)$, uniformly in
$t\in[0,T]$ in probability. We now prove this statement.

\subsection{Deletion estimates and the correct row law}

For the true posterior mean, the Gaussian response estimate and the
already proved qualitative covariance bound give, for each deleted row,
\[
 \frac1N\mathbb E\|m_{t,R}-m_t^{\,-i}\|^2
 \le\frac C{N^2}\mathbb E\operatorname{Tr}[(C_t^{\,-i})^2]=o(1),
\]
uniformly on fixed time intervals.  For original means the same estimate
is $O(N^{-1})$, by the original overlap bound.

For the cavity algorithms, deleting one additional vertex changes a
fixed coordinate in mean square by $O_{L,j}(N^{-1})$.
Here is a direct induction. The assertion is trivial at the constant
leaves of the original tree.  At an original-field layer, condition on
everything outside the tested row, whose Gaussian entries are independent
of its child vectors; Gaussian isometry and the Lipschitz property of $\tanh$ propagate
the estimate, with an additional $O(N^{-1})$ from the deleted edge.
At an observed-field layer use the exact row density from the preceding
proof.  If $e$ is any child error vector independent of the tested row
under the cavity reference experiment, then
\[
 \mathbb E_g R(g,Y)(g\cdot e)^2
 \le \cosh h\,e^{vr/2}
       \mathbb E_{\xi,\sigma^0}
          \left[v\|e\|^2+v^2(\sigma^0\cdot e)^2\right]
 \le C_{\beta,h}v\|e\|^2.
\]
The missing edge has second moment $O(v)$ under the same density.
The reference child observation is exactly that of the smaller original
cavity model, so the induction applies. Constants may depend on fixed
depth, but not on $N,t$.

These estimates also imply that deleting a tested vertex changes every
empirical product used in $(p_N,c_N,b_N)$ by $o(1)$ in $L^1$, uniformly in
time. Products are bounded, and Cauchy--Schwarz transfers the normalized
vector estimates to empirical products. No union bound over all the
polynomially many cavity calls is used or needed.

For the state evolution it is convenient to retain the true spin $X_i$.
Under the independent cavity/singleton reference experiment,
$X_i\sim p_h$, $Y_i=tX_i+\sqrt t Z_i$; the reference cavity spin is drawn
from the original cavity law and generates $Y_R$.
Conditional on the cavity observation, let $\nu=\mu^{-i}_{Y_R}$.
The exact density of $(g,X_i,Y_R,Y_i)$ under the full model relative to
this reference experiment is
\[
 R_{X_i}(g,Y_R)=
 \frac{\mathbb E_\nu e^{X_i g\cdot\sigma}}
 {\mathbb E_{\mu^{-i}}\cosh(h+g\cdot\sigma)/\cosh h}.
\]
It does not involve $Y_i$ once $X_i$ is retained.
Write $m^0,m^t$ for the cavity original and posterior means, and $M$ for
the child algorithm vector.  Conditional on cavity data,
\[
 (g\cdot m^0,g\cdot m^t,g\cdot M)
\]
is centered Gaussian. Its limiting covariance matrix, under the induction
hypothesis, is
\[
 \beta^2
 \begin{pmatrix}
 q&q&b\\q&r(t)&c\\b&c&p
 \end{pmatrix}.
\]
The full Gaussian partition estimate approximates the displayed row
density in $L^1$ by
\[
 \widetilde R_{X_i}
 =e^{X_i g\cdot m^t-\beta^2(r(t)-q)/2}
       \frac{\cosh h}{\cosh(h+g\cdot m^0)},
\]
up to errors tending to zero uniformly in $t$ on compact intervals.
Indeed, the error is bounded by the same $v\sqrt{\operatorname{Tr}C^2}$
terms as before, together with deviations of the empirical covariance
parameters; these tend to zero by the preceding scalar theorem and the
induction hypothesis. Gaussian covariance interpolation is legitimate
also at singular limiting covariance matrices, by a small independent
Gaussian regularization followed by dominated convergence.

\subsection{The tilted Gaussian calculation}

We compute the limiting row law explicitly.
Write $A_0=g\cdot m^0$, $A_t=g\cdot m^t$, and $A_M=g\cdot M$.
Under the limiting reference Gaussian law,
\[
 A_t=A_0+G_t,\qquad A_M=(b/q)A_0+G_M,
\]
where $A_0\sim N(0,\beta^2q)$ is independent of $(G_t,G_M)$, and
\[
 \operatorname{Var}G_t=\beta^2(r(t)-q),\quad
 \operatorname{Var}G_M=\beta^2(p-b^2/q),\quad
 \operatorname{Cov}(G_t,G_M)=\beta^2(c-b).
\]
Multiplying the reference mass $p_h(X_i)$ by the approximate row density
gives
\[
 \frac{e^{X_i(h+A_0)}}{2\cosh(h+A_0)}
 \exp\left(X_iG_t-\frac12\operatorname{Var}G_t\right).
\]
Thus $U=h+A_0$ retains its original Gaussian law,
$X_i\mid U\sim p_U$, and Gaussian exponential tilting shifts $G_M$ by
$\beta^2(c-b)X_i$, without changing its variance.
Adding $Y_i=tX_i+\sqrt t Z_i$ gives exactly the scalar field
$V_{t,\theta}$ defined above.

The algorithm output at the tested coordinate is exactly
$\tanh(h+Y_i+g\cdot M)$.
The original true mean at that coordinate is
$\tanh(h+g\cdot m^0)+o_{L^2}(1)$ under the row law, by the original cavity
estimate and a Gaussian shift. Therefore the one-row limiting
expectations of its squared algorithm output and its product with the
original mean are the first and third coordinates of $\Gamma_t$.
For the second coordinate use Bayes orthogonality:
\[
 \mathbb E[X_iF(J,Y_t)]=\mathbb E[m_{t,i}F(J,Y_t)]
\]
for every bounded measurable $F$. This identity continues to hold after
multiplication by any bounded cavity-observation-measurable test.
The second coordinate is consequently $\mathbb E\xi\tanh V_{t,\theta}$.

All these row identities hold after multiplication by an arbitrary
bounded function of the cavity disorder and cavity observation.  To prove
concentration, let $A_N$ be one of the three empirical products at the
new iteration, let $a$ be its proposed limit, and let $A_i^c$ denote its
version with the tested row deleted.  Then
\[
 \mathbb E(A_N-a)^2
 =\frac1N\sum_i\mathbb E(A_N-a)(F_i-a)
 =\frac1N\sum_i\mathbb E(A_i^c-a)(F_i-a)+o(1)=o(1).
\]
The middle equality uses deletion stability, and the last equality uses
the tested row identity with the bounded test $A_i^c-a$.
This proves finite-step joint $L^2$ state evolution without an independence
assumption on the full row and full observation.

At the initial observed layer, the original tree has
$p_N\to q$ and $b_N\to\rho_L$ in $L^2$.
For its product with the current mean, if $a=B^L(J)$ then
\[
 \mathbb E\left[\left(\frac{a\cdot(m_t-m_0)}N\right)^2\right]
 \le\frac1{N^2}\mathbb E[a^{\mathsf T}C_0a]
 \le\frac{\mathbb E\sqrt{\operatorname{Tr}(C_0^2)}}N
 =O(N^{-1/2}).
\]
This holds uniformly in $t$ by total variance and $\|a\|^2\le N$.
Hence $c_N\to\rho_L$, establishing the initial triple.
The required original-tree initialization is the cited Chen--Tang state evolution.

\subsection{Uniformity in time}

For every fixed observed depth $j$ and $T<\infty$,
\[
 \mathbb E\left[
 \left(\frac{\|M^{j,L}(Y_t)-M^{j,L}(Y_s)\|}{\sqrt N}\right)^4\right]
 \le C_{j,T}|t-s|^2,\qquad 0\le s\le t\le T,
\]
with a constant independent of $L,N$.
The leaves do not depend on the observation, so their increment is zero.
At each observed layer, $\tanh$ is Lipschitz and the field increment is a
singleton increment plus a Gaussian row applied to the child increment.
Use the exact row density at the later time $t$ for the entire path up to
$t$; the Gaussian channel endpoint is sufficient, so this is the same
row density as above.  Conditional on the reference cavity path, the
child increment $e$ is independent of the tested row, and Gaussian shifting
gives
\[
 \mathbb E_gR(g,Y_t)|g\cdot e|^4\le C_{\beta,h}v^2\|e\|^4.
\]
The singleton increment has fourth moment at most $C_T|t-s|^2$.
The inequality $\|e\|^4\le r\sum_a e_a^4$ closes an induction on depth,
first for a single coordinate and then for the normalized vector.

The standard dyadic proof of the Kolmogorov continuity criterion gives a
modulus of continuity in probability, uniformly in $N,L$, for each fixed
$j$. The true mean has the corresponding modulus by the scalar-law proof.
Therefore the three empirical state-evolution quantities have a uniform
modulus on fixed time intervals. Combining a finite grid with their
fixed-time $L^2$ convergence proves uniform-in-time convergence in
probability to the continuous deterministic recursion $\theta_{j,L}(t)$.

\subsection{Convergence of the initializer and explicit depth choices}

Consider first the ideal initial triple $(q,q,q)$, corresponding to the
original true means at the leaves. If $(p,c,b)=(p,p,q)$ with $q\le p\le r(t)$,
then
\[
 V_{t,\theta}=U+w\xi+\sqrt w Z,\qquad
 w=t+\beta^2(p-q).
\]
The scalar posterior identities give
\[
 \Gamma_t(p,p,q)=(a(w),a(w),q).
\]
Thus, with $p_0=q$ and
$p_{j+1}=a(t+\beta^2(p_j-q))$,
\[
 q\le p_j\le r(t),\qquad
 0\le r(t)-p_j\le\kappa^j.
\]
The normalized squared error of the ideal depth-$j$ tree converges to
$r(t)-p_j(t)$.

One can choose the finite original depth by a direct stability bound.
For two child vectors differing by $e$, the row comparison above gives
\[
 \mathbb E_gR(g,Y)(g\cdot e)^2
 \le \cosh h\,e^{\beta^2/2}v(1+\beta^2)\|e\|^2.
\]
Consequently, an original-leaf normalized mean-square error is multiplied
at each observed layer by at most
\[
 K=\beta^2\cosh h\,e^{\beta^2/2}(1+\beta^2).
\]
The original depth-$L$ error has limiting value
$2(q-\rho_L)\le2q\alpha^L$.
Minkowski's inequality therefore bounds the limiting normalized error by
\[
 \sqrt{r(t)+p_{k,L}(t)-2c_{k,L}(t)}
 \le \kappa^{k/2}+K^{k/2}\sqrt{2q\alpha^L},
\]
uniformly in $t$.
Choose fixed $k,L$ so that each term on the right is at most $\eta/4$.
Uniform state evolution then implies
$\sup_{t\le T}\|M^{k,L}(Y_t)-m_t\|/\sqrt N\le\eta$ with probability
tending to one. This proves the theorem.

This theorem supplies a fully specified polynomial-time initializer under
the actual disorder law. \S\ref{sec:amp} reduces its cost to a fixed
number of matrix-vector products. The local-curvature estimate in
Theorem~\ref{curv:fixed-anchor} then permits refinement by constrained
convex optimization, with drift action controlled by
Theorem~\ref{thm:integrated-tap}.


\subsection{AMP comparison at fixed observation}\label{sec:amp}

This section upgrades the finite-depth self-avoiding initializer to a
fixed-depth AMP algorithm.  We use the true scalar law of
\S\ref{sec:scalar-law} and the coarse cavity initializer in
Theorem~\ref{thm:initializer}.  In particular, the observations are generated
from the original quenched Gibbs law throughout this proof.  There is no
planted change of the disorder law.

Write $J_{ij}=\beta A_{ij}/\sqrt N$, with $A_{ij}$ independent standard
Gaussian for $i<j$, $A_{ii}=0$.  Given $J$, let $X\sim\mu_J$ and
$Y_t=tX+B_t$.  For fixed $t$, let $\mathbb P_t$ denote the joint law of
$(A,Y_t)$.  Relative to the product of the original Gaussian matrix law and
an independent $N(0,tI)$ observation, its density is
\[
 \rho_t(A,y)=e^{-Nt/2}\frac{Z_J(h+y)}{Z_J(h)}.
\]
All matrix derivatives below hold $y$ fixed.  For every multiset $P$ of
$d$ edges, the spin cumulant formula and the product rule give
\begin{equation}\label{amp:density-derivatives}
 \left|\frac{\partial_P\rho_t}{\rho_t}\right|
 \le C_{d,\beta}N^{-d/2}.
\end{equation}
Indeed, a derivative of $\log Z_J$ is $(\beta/\sqrt N)^d$ times a joint
cumulant of $d$ random variables bounded by one.  The cumulant bound is
uniform in the field, disorder, and Gibbs measure.  The same statement
holds at $t=0$, with the deterministic observation zero.

\subsection{A conditional Gaussian-row moment bound}

Conditional on all entries outside a row and on $Y_t=y$, the row
$a=(A_{ij})_{j\ne i}$ has a density proportional to
$e^{\ell(a)}$ against independent standard Gaussians, where
\[
 \left|\partial_{A_{ij}}\ell\right|
 =\frac\beta{\sqrt N}
   \left|\langle\sigma_i\sigma_j\rangle_{h+y}
               -\langle\sigma_i\sigma_j\rangle_h\right|
 \le\frac{2\beta}{\sqrt N},\qquad
 \|\nabla\ell\|_2\le2\beta.
\]
The Gaussian exponential-moment inequality for Lipschitz functions gives
\[
 \frac{\mathbb E_\gamma e^{2\ell}}
      {(\mathbb E_\gamma e^\ell)^2}\le e^{8\beta^2}.
\]
Consequently, for every $p\ge1$ and every vector $e$ independent of that
row when $y$ and the other entries are held fixed,
\begin{equation}\label{amp:row-moment}
 \mathbb E_t\bigl[|a\cdot e|^p\mid A_{-i,-i},Y_t\bigr]
 \le C_{p,\beta}\|e\|_2^p.
\end{equation}
The same bound applies to a subvector of a row and to a row containing
some omitted entries.  Constants are uniform in $t,N,y$ and the conditioned
matrix entries.  This is the only conditional row bound needed below.

\subsection{Cavity and AMP recursions}

Fix a finite depth $d$.  For $0\le s\le d$, let $f_{s,i}(z;y)$ be bounded
smooth functions whose derivatives in $z$ of every fixed order are
uniformly bounded in $N,i,y$.  Their dependence on $y$ is coordinatewise;
in the application below it is just an additive external field.
Take $f_{0,i}\equiv\sqrt q$ and define
\[
 w^0_{S,i}=0,\qquad
 w^{s+1}_{S,i}=\frac1{\sqrt N}
       \sum_{j\notin S\cup\{i\}} A_{ij}
          f_{s,j}(w^s_{S\cup\{i\},j};y),\qquad
 x^s_{S,i}=f_{s,i}(w^s_{S,i};y).
\]
Write $w^s=w^s_\varnothing$, $x^s=x^s_\varnothing$.
The comparison AMP is
\begin{align*}
 u^0&=0,\qquad u^1=N^{-1/2}A f_0,\\
 u^{s+1}&=N^{-1/2}A f_s(u^s;y)
      -b_s^{\rm AMP}f_{s-1}(u^{s-1};y),\\
 b_s^{\rm AMP}&=\frac1N\sum_i f'_{s,i}(u_i^s;y).
\end{align*}

\begin{prop}\label{amp:cavity-comparison}
For every fixed depth $d$ there is a constant $C_d$, independent of
observation time $t$, such that
\[
 \sup_{0\le s\le d}\frac1N
  \mathbb E_t\|f_s(u^s;Y_t)-x^s(Y_t)\|_2^2\le\frac{C_d}{N}.
\]
For the original/observed $\tanh$ denoisers specified below, the difference
converges to zero uniformly over each fixed time interval in probability.
\end{prop}

\subsection{Elementary cavity moment and derivative estimates}

We give all estimates under the full joint law $\mathbb P_t$, including
for algorithms defined on deleted sets.  This avoids replacing the
observation by an independently sampled side field.

The conditional row estimate also
gives the following explicit $L^p$ induction rule. If $p\ge2$ and $e$ is
measurable with respect to the row complement and $Y_t$, then
\[
 \left\|\frac1{\sqrt N}\sum_j A_{ij}e_j\right\|_{L^p(\mathbb P_t)}
 \le\frac{C_{p,\beta}}{\sqrt N}
       \left(\sum_j\|e_j\|_{L^p(\mathbb P_t)}^2\right)^{1/2}.
\]
Indeed, condition first in \eqref{amp:row-moment} and then use
Minkowski's inequality in $L^{p/2}$ for $\sum_j e_j^2$.
The inductions below are proved simultaneously for every fixed $p$.
In a chain-rule product of at most $m$ factors, H\"older's inequality
uses the already available bounds at a larger fixed moment order
(for instance $mp$). This closes the induction at every fixed derivative
and moment order.

For each fixed $s,p$, deleting one additional vertex $a$ gives
\begin{equation}\label{amp:delete}
 \|w^s_{S,i}-w^s_{S\cup\{a\},i}\|_{L^p(\mathbb P_t)}
 +\|x^s_{S,i}-x^s_{S\cup\{a\},i}\|_{L^p(\mathbb P_t)}
 \le C_{s,p}N^{-1/2}.
\end{equation}
For $p\ge2$, this follows by induction from
\eqref{amp:row-moment}: the deleted edge contributes $O(N^{-1/2})$,
and the remaining row is applied to a child error vector whose individual
$L^p$ norms are $O(N^{-1/2})$.  Its $\ell_2$ norm has $L^p$ norm $O(1)$.
The smaller $p$ follow by Jensen.  The induction is uniform over every
fixed number of deleted vertices.

For a multiset $P$ of $m$ edges, one similarly has
\begin{equation}\label{amp:derivative}
 \|\partial_P w^s_{S,i}\|_{L^p(\mathbb P_t)}
 +\|\partial_P x^s_{S,i}\|_{L^p(\mathbb P_t)}
 \le C_{s,m,p}N^{-m/2}.
\end{equation}
For the fields this is especially simple.  If no edge of $P$ meets the
root $i$, every differentiated child is independent of row $i$ at fixed
$y$, so \eqref{amp:row-moment} and induction apply.  If exactly one
edge of $P$ meets $i$, differentiation selects its neighbor and contributes
$N^{-1/2}$, followed by a derivative of order $m-1$ of the child.
If at least two differentiations meet $i$, the differentiated field is
zero, since the field is linear in row $i$.  The activation estimate
follows from the chain rule and H\"older's inequality.  The same reasoning
applies if one term is omitted from the outermost row sum.

We will also use a sharper first-derivative estimate.  For an edge
$e=\{a,b\}$ not incident to the root,
\begin{equation}\label{amp:nonincident}
 \|\partial_{A_{ab}} w^s_{S,i}\|_{L^p(\mathbb P_t)}
 +\|\partial_{A_{ab}} x^s_{S,i}\|_{L^p(\mathbb P_t)}
 \le C_{s,p}N^{-1},\qquad i\notin\{a,b\}.
\end{equation}
If an endpoint is deleted, the derivative is zero.  Otherwise, among the
children of root $i$, only the two roots $a,b$ are incident to $e$; their
derivatives are $O(N^{-1/2})$ by the exact linear-row formula.  All other
children have derivatives $O(N^{-1})$ by induction.  Their derivative
vector therefore has $L^p(\ell_2)$ norm $O(N^{-1/2})$, and
\eqref{amp:row-moment}, including the outer $N^{-1/2}$, proves
\eqref{amp:nonincident}.  Again the estimate remains valid with an
omitted term in the outer row sum.  This refinement is what controls the
loop error below.

\subsection{The cavity recursion satisfies the AMP equation}

Set $b_s^c=N^{-1}\sum_j f'_{s,j}(w_j^s;y)$.  We prove the coordinate bound
\begin{equation}\label{amp:defect}
 \mathbb E_t\left|w_i^{s+1}
       -N^{-1/2}(Ax^s)_i+b_s^c x_i^{s-1}\right|^2
 \le C_s/N.
\end{equation}
The depth-one identity is exact.  Fix $s\ge1$ and a root $i$.  For $j\ne i$
put
\[
 z_j=\frac1{\sqrt N}\sum_{\ell\ne i,j}A_{j\ell}
          x^{s-1}_{\{j\},\ell},\qquad
 z'_j=\frac1{\sqrt N}\sum_{\ell\ne i,j}A_{j\ell}
          x^{s-1}_{\{i,j\},\ell},
\]
\[
 L_j=f_{s,j}(z_j;y),\quad K_j=f_{s,j}(z'_j;y),\quad
 B_{ij}=f'_{s,j}(z_j;y),\quad D_{ij}=x^{s-1}_{\{j\},i}.
\]
Then $w_i^{s+1}=N^{-1/2}\sum_{j\ne i}A_{ij}K_j$, while Taylor's formula
gives
\[
 N^{-1/2}(Ax^s)_i
 =N^{-1/2}\sum_{j\ne i}A_{ij}L_j
   +\frac1N\sum_{j\ne i}A_{ij}^2B_{ij}D_{ij}+E_i,
 \qquad \|E_i\|_{L^2(\mathbb P_t)}\le C_sN^{-1/2}.
\]
The last bound follows directly from
$|E_i|\le C_sN^{-3/2}\sum_{j\ne i}|A_{ij}|^3$ and the original Gaussian
marginal of $A$.

The quadratic correction can be centered.  Let $H_j=B_{ij}D_{ij}$.
By \eqref{amp:derivative}, every fixed derivative of $H_j$ obeys the
corresponding $N^{-m/2}$ moment bound.  For $j\ne\ell$, two Gaussian
Hermite integrations by parts in each of $A_{ij},A_{i\ell}$ yield
\[
 \mathbb E_t[(A_{ij}^2-1)(A_{i\ell}^2-1)H_jH_\ell]
 =\mathbb E_t\left[
  \rho_t^{-1}\partial_{A_{ij}}^2\partial_{A_{i\ell}}^2
                  (\rho_tH_jH_\ell)\right]=O(N^{-2}).
\]
Every term has a total of four derivatives distributed between the
density and the two smooth factors; \eqref{amp:density-derivatives}
and \eqref{amp:derivative} give the stated bound.  The diagonal terms
are bounded.  Hence
\[
 \left\|\frac1N\sum_{j\ne i}(A_{ij}^2-1)B_{ij}D_{ij}\right\|_{L^2}
 \le C_sN^{-1/2}.
\]
The mean value theorem gives
$\|B_{ij}-f'_{s,j}(w_j^s;y)\|_{L^2}\le C_sN^{-1/2}$, and
\eqref{amp:delete} gives
$\|D_{ij}-x_i^{s-1}\|_{L^2}\le C_sN^{-1/2}$.  Thus the entire quadratic
correction equals $b_s^cx_i^{s-1}+O_{L^2}(N^{-1/2})$.

It remains to bound the loop expression
\[
 F_i=\frac1{\sqrt N}\sum_{j\ne i}A_{ij}E_j,
 \qquad E_j=L_j-K_j.
\]
By \eqref{amp:delete} and a conditional moment bound in row $j$,
$\|E_j\|_{L^p}\le C_{s,p}N^{-1/2}$ for every fixed $p$.
When $y$ is held fixed, $E_j$ is independent of $A_{ij}$, and $K_j$ is
independent of the entire row $i$.  For $\ell\ne i,j$,
\eqref{amp:nonincident} gives
\[
 \|\partial_{A_{i\ell}}E_j\|_{L^p}
 =\|\partial_{A_{i\ell}}L_j\|_{L^p}\le C_{s,p}N^{-1}.
\]
For a distinct pair $j,\ell$, write $a=A_{ij}$, $b=A_{i\ell}$,
$S_a=(\partial_a\rho_t)/\rho_t$, and
$T_{ab}=(\partial_a\partial_b\rho_t)/\rho_t$.
Gaussian integration by parts twice, using
$\partial_aE_j=\partial_bE_\ell=0$, gives exactly
\[
 \mathbb E_t[abE_jE_\ell]
 =\mathbb E_t\left[
  (\partial_bE_j)(\partial_aE_\ell)
  +S_a(\partial_bE_j)E_\ell
  +S_bE_j(\partial_aE_\ell)+T_{ab}E_jE_\ell\right].
\]
Each term is $O(N^{-2})$: use respectively the nonincident derivative
bound twice, that bound together with $S_a=O(N^{-1/2})$ and
$\|E_\ell\|_2=O(N^{-1/2})$, or
$T_{ab}=O(N^{-1})$ together with the two deletion errors.  The diagonal
term $\mathbb E_t[A_{ij}^2E_j^2]$ is $O(N^{-1})$ by H\"older's inequality.
Expanding $\mathbb E_tF_i^2$ therefore proves
$\mathbb E_tF_i^2\le C_s/N$.  This establishes \eqref{amp:defect}.

\subsection{Transfer to AMP and the explicit initializer}

The Gaussian matrix marginal is unchanged under $\mathbb P_t$, so
$\|A/\sqrt N\|_{\rm op}$ is bounded by a fixed constant with an
exponentially small exceptional probability.  On this event, the finite
AMP recursion is Lipschitz in the previous fields and in its additive
defects.  Apply this deterministic bound successively to
\eqref{amp:defect}.  The resulting mean square difference of normalized
activation vectors is $O_d(N^{-1})$.  The activations are bounded on the
exceptional event.  This proves Proposition~\ref{amp:cavity-comparison}.

For the original/observed initializer take
\[
 f_{0,i}\equiv\sqrt q,\qquad
 f_{s,i}(z;y)=\tanh\bigl(h+\mathbf1_{\{s>L\}}y_i+\beta z\bigr),
 \quad s\ge1.
\]
In activation variables the algorithm is particularly simple:
This is exactly the two-phase recursion \eqref{amp:two-phase}.
The $s=0$ correction is zero.  The original Onsager memory is retained
when the observation field is first inserted.  The output is $a^{L+k}$,
and computing it takes $O((L+k)N^2)$ arithmetic operations.

For fixed depth, on the same bounded matrix norm event, this AMP output
is a Lipschitz function of $y$ in normalized Euclidean norm, with a
constant depending only on depth and $\beta$.  The process
$Y_t=tX+B_t$ has a uniform normalized fourth-moment modulus on fixed
time intervals.  The cavity algorithm has the corresponding modulus by
the preceding initializer proof.  A finite time grid and these two
moduli therefore upgrade the fixed-time comparison to
\[
 \sup_{0\le t\le T}\frac{\|a^{L+k}(Y_t)-M^{k,L}(Y_t)\|_2}{\sqrt N}
 \xrightarrow{\mathbb P}0
\]
for fixed $L,k,T$.  This proves the following algorithmic consequence.

\medskip\noindent\emph{Proof of Theorem~\ref{amp:actual}.}
\begin{proof}
Choose the depths in Theorem~\ref{thm:initializer} for normalized error
$\eta/2$. The uniform AMP--cavity comparison just proved has error
$o_{\mathbb P}(1)$ in normalized norm. Their sum is at most $\eta$ with
probability tending to one.
\end{proof}

\paragraph{Joint empirical transcript information.}
The comparison holds simultaneously at every fixed collection of old and
new depths.  The tested one-row state-evolution proof in the cavity
initializer section also applies to bounded smooth functions of a finite
tuple of cavity fields, by keeping their full Gaussian covariance matrix.
Thus the entire finite old/new AMP transcript has the corresponding joint
empirical state evolution under the actual quenched observation law.
Explicitly, if $p_{ab}$ are child-algorithm overlaps and $b_a,c_a$ are
their overlaps with $m_0,m_t$, use $U=h+\beta\sqrt q Z_0$, $\xi\mid U\sim p_U$,
and a Gaussian vector $(G_a)$, independent of $U,\xi,Z$, with covariance
$\mathbb E G_aG_b=\beta^2(p_{ab}-b_ab_b/q)$.  The next fields are
\[
 H_a=h+\frac{b_a}{q}(U-h)+\beta^2(c_a-b_a)\xi+G_a
           +\iota_a(t\xi+\sqrt t Z),
\]
where $\iota_a$ is zero for an original-field layer and one for an
observed-field layer.  Pairwise new overlaps are
$\mathbb E\tanh H_a\tanh H_b$, with the analogous base and teacher
products.  This describes empirical transcript limits; it does not assert
that the matrix conditioned on the transcript is an unweighted Gaussian
matrix under the actual law.


\section{Full TAP curvature along the observation path}\label{curv:section}

We prove Theorem~\ref{curv:fixed-anchor}. The argument conditions on an
original Gaussian transcript, constructs a product-prior reference law on
the remaining matrix, and transfers an exponentially unlikely curvature
failure to the actual quenched observation experiment. Gaussian completion
is used to keep the conditional exponential rate independent of the
original transcript depth.

Throughout this appendix, $J$ is the hollow interaction matrix with
off-diagonal variance $\beta^2/N$, $X\sim\mu_J$, and $Y_t=tX+B_t$.
For a symmetric interaction $G$ and a total field $b$, write
\begin{align}
 \mathcal F_b^G(w)
 &=\sum_i I(w_i)-\frac12w^{\mathsf T}Gw-b^{\mathsf T}w
       -\frac{\beta^2N}{4}(1-q_w)^2,
 \qquad q_w=N^{-1}\|w\|^2,\label{curv:tap-functional}\\
 \mathcal H_G(w)
 &=\operatorname{diag}\bigl((1-w_i^2)^{-1}\bigr)-G
       +\beta^2(1-q_w)I_N-\frac{2\beta^2}{N}ww^{\mathsf T}.
 \label{curv:hessian}
\end{align}
Here $I$ is the negative binary entropy used in the definition of
$F_{\mathrm{TAP}}$. Thus
$\mathcal F_b^J(w)=F_{\mathrm{TAP}}(w,b)$ and
$\mathcal H_G=\nabla^2\mathcal F_b^G$ is independent of $b$.

Fix $0<\delta<T_0<\infty$. We use the scalar quantities from
Proposition~\ref{prop:scalar-law}:
\begin{equation}\label{curv:scalar-parameters}
 \alpha=\beta^2\mathbb E\sech^4(h+\beta\sqrt q Z)<1,
 \qquad \sup_{v\ge0}\beta^2a'(v)=:\kappa<1,
 \qquad r(t)=a\bigl(t+\beta^2(r(t)-q)\bigr).
\end{equation}
All quantitative statements below concern a prescribed positive error
tolerance. At a fixed new depth $L$, their conditional failure probabilities
are bounded by $\exp\{C(K+L)-cN\}$, where $c>0$ is independent of the
original depth $K$. The threshold for $N$ may depend on $K$.
We choose $L$ and $c$ before choosing $K$. No exponential bound is required
for the original transcript to satisfy its empirical good conditions.

\subsection{The scalar inequality for longitudinal directions}\label{sec:longitudinal}
We first prove the finite-dimensional inequality used to control the
longitudinal Hessian directions. The subsequent subsections reduce the
random Hessian to this matrix, retaining both strict AT margins.

\paragraph{The scalar Gaussian observation.}
Let $U$ be any finite real random variable satisfying
$\mathbb P(\tanh U\ne0)>0$. Set $m=\tanh U$, draw
$X\in\{-1,1\}$ conditionally on $U$ with $\bbE[X\mid U]=m$, and let $Z$ be an
independent standard Gaussian. For $v>0$, put
\[
 Y=vX+\sqrt v Z,\qquad n=\tanh(U+Y),\qquad s=1-n^2.
\]
Bayes' rule gives $n=\bbE[X\mid U,Y]$. In particular,
\begin{equation}\label{long:eq:mmse}
 \bbE[n\mid U]=m,
 \qquad \bbE[s\mid U]=1-\bbE[n^2\mid U]\le 1-m^2.
\end{equation}
Define
\begin{align*}
 A&=\bbE[s m^2], & B&=\bbE[s m(n-m)], & C&=\bbE[s(n-m)^2],\\
 u&=\bbE[s^2], & u_0&=\bbE[(1-m^2)^2], &
 Q&=\begin{pmatrix}A&B\\B&C\end{pmatrix}.
\end{align*}
The standard SK specialization is $U=h+\beta\sqrt q\,Z_0$ and
$\lambda=\beta^2$. The algebra below does not require this specialization.

\begin{lem}[Gram positivity and the second AT margin]\label{long:lem:scalar}
For every $v>0$, $Q\succ0$. Moreover, with $k=2B+C$,
\begin{equation*}
 u+k=\bbE[s(1-m^2)]\le u_0.
\end{equation*}
Consequently, if $\lambda\ge0$, $a=1-\lambda u$ and
$\delta_0=1-\lambda u_0>0$, then
\[
 a-\lambda k\ge\delta_0>0.
\]
\end{lem}
\begin{proof}
The matrix $Q$ is the Gram matrix of $\sqrt s\,m$ and
$\sqrt s\,(n-m)$. Since $v>0$, conditional on each finite $U$, the variable
$n$ has a density with support $(-1,1)$ and $s>0$ almost surely. If
$\alpha m+\gamma(n-m)=0$ almost surely, this conditional nondegeneracy
first forces $\gamma=0$, and $\mathbb P(m\ne0)>0$ then forces $\alpha=0$.
Thus $Q\succ0$. Finally,
\[
 u+2B+C
 =\bbE[s\{s+2m(n-m)+(n-m)^2\}]
 =\bbE[s(1-m^2)].
\]
Multiply the conditional inequality in \eqref{long:eq:mmse} by the nonnegative
$U$-measurable variable $1-m^2$ and take expectations.
\end{proof}

\paragraph{The matrix inequality.}
Define
\begin{align*}
 T&=\begin{pmatrix}
 2A-B-C&2A+2B\\
 4B+2C&2B+2C
 \end{pmatrix},
 &L&=(1-\lambda u)I+\lambda T,\\
 M&=\begin{pmatrix}2&2\\2&3\end{pmatrix},
 &K&=\begin{pmatrix}
 -2u+4A+2B&-2u+4A+6B+2C\\
 -2u+4A+6B+2C&4A+8B+4C-3u
 \end{pmatrix},
\end{align*}
and
\begin{equation}\label{long:eq:F}
 F=\lambda M+\lambda^2K-L^\top Q^{-1}L.
\end{equation}

\begin{thm}[Algebraic negativity]\label{long:thm:matrix}
Let $Q\succ0$, $\lambda\ge0$, and $u\in\mathbb R$. If
\[
 a:=1-\lambda u>0,
 \qquad a-\lambda(2B+C)>0,
\]
then the symmetric matrix in \eqref{long:eq:F} satisfies $F\prec0$.
No sign condition on $B$ is needed.
\end{thm}
\begin{proof}
Write
\[
 \Delta=AC-B^2>0,\qquad k=2B+C,
 \qquad P=a^2+2A\lambda a-\lambda^2(2AB+k^2+B^2).
\]
Direct multiplication gives the two identities
\begin{align}
 F_{22}&=-a\lambda-\frac{Aa^2}{\Delta}<0,
 \label{long:eq:diag}\\
 \det F&=\frac{a(a-\lambda k)P}{\Delta}.
 \label{long:eq:det}
\end{align}
We prove $P>0$ by a sign split on $k$.

If $k\ge0$, then
\begin{equation}\label{long:eq:posk}
 P=(a-\lambda k)\{a+\lambda(k+2A)\}
     +\lambda^2(\Delta+Ak)>0.
\end{equation}
Indeed $A>0$, and the first product is strictly positive.

If $k\le0$, then $B\le-C/2<0$, and
\begin{equation}\label{long:eq:negk}
 2AB+k^2+B^2
 =\frac{2B\Delta+k(B+C)^2}{C}<0.
\end{equation}
The first term in the numerator is strictly negative and the second is
nonpositive. Therefore $P\ge a^2+2A\lambda a>0$ (strictly greater if
$\lambda>0$). Equation \eqref{long:eq:det} now gives $\det F>0$.
Together with \eqref{long:eq:diag}, this is the negative-definite criterion for
a symmetric two-by-two matrix.
\end{proof}

\begin{cor}[The scalar AT implication]\label{long:cor:AT}
For the scalar Gaussian observation above, assume
\[
 \lambda u_0<1,\qquad \lambda u<1.
\]
Then $F\prec0$ for every $v>0$.
\end{cor}
\begin{proof}
The first inequality and Lemma~\ref{long:lem:scalar} give
$a-\lambda k>0$, while the second gives $a>0$.
Apply Theorem~\ref{long:thm:matrix}.
\end{proof}

\paragraph{An explicit uniform gap.}
The strict inequality has a quantitative version that is independent of the
smallest eigenvalue of $Q$. This matters when $v$ approaches zero and $Q$
becomes nearly singular.

\begin{cor}[A uniform numerical margin]\label{long:cor:gap}
Under the scalar assumptions above, suppose
\[
 1-\lambda u_0\ge\eta,
 \qquad 1-\lambda u\ge\eta,
 \qquad 0<\eta\le1.
\]
Then, for every $v>0$,
\begin{equation*}
 F\preceq-\frac{\eta^4}{20+20\lambda+5200\lambda^2}\,I.
\end{equation*}
\end{cor}
\begin{proof}
Put $b=a-\lambda k$. Lemma~\ref{long:lem:scalar} gives $b\ge\eta$.
Equations \eqref{long:eq:posk} and \eqref{long:eq:negk} imply in both cases
\[
 P\ge\min(a,b)(a+2\lambda A)\ge\eta^2,
 \qquad abP\ge\eta^4.
\]
Let $G=\lambda M+\lambda^2K$ and
\[
 H=-\Delta\Tr F
   =\Tr\{L^\top\operatorname{adj}(Q)L\}
      -\Delta\Tr G>0.
\]
The scalar definitions imply
\[
 0<A\le1,\quad 0<C\le4,\quad |B|\le2,
 \quad 0\le u\le1,\quad 0<a\le1.
\]
Hence $\|\operatorname{adj}(Q)\|_{\rm op}\le\Tr Q\le5$,
$\Delta\le4$, $\|T\|_F^2\le500$, and
\[
 |\Tr G|\le5\lambda+49\lambda^2.
\]
It follows that
\begin{align*}
 H&\le5(\sqrt2+\sqrt{500}\lambda)^2
          +4(5\lambda+49\lambda^2)\\
  &\le20+20\lambda+5196\lambda^2
   \le20+20\lambda+5200\lambda^2.
\end{align*}
For a positive-definite two-by-two matrix, its least eigenvalue is at least
its determinant divided by its trace. Applying this to $-F$ and using
\eqref{long:eq:det} yields
\[
 \lambda_{\min}(-F)
 \ge\frac{\det(-F)}{\Tr(-F)}
 =\frac{abP}{H}
 \ge\frac{\eta^4}{20+20\lambda+5200\lambda^2}.
\]
\end{proof}


\subsection{The original Gaussian transcript and its measurable good events}
\label{curv:old-transcript}

Add independent diagonal entries $D_{ii}\sim N(0,2\beta^2/N)$ and set
$G=J+D$. This is GOE at scale $\beta$. Its Gibbs measure equals $\mu_J$,
because $x_i^2=1$, while
\begin{equation}\label{curv:diagonal}
 \|D\|_{\mathrm{op}}\longrightarrow0
 \quad\text{in probability},\qquad
 \mathcal H_J(w)-\mathcal H_G(w)=D.
\end{equation}
We work with $G$ until the final transfer.

Use the canonical original Bolthausen iteration, with its deterministic
scalar coefficients, and record a finite adaptive linear transcript
$\mathcal C_K$. Its queried span contains
$m^{(1)},\ldots,m^{(K)}$; in particular record the answer to $Gm^{(K)}$.
Let $P$ be the orthogonal projection onto that span and let $Q=I-P$.
An orthonormalized implementation gives the same transcript information.
Set $m=m^{(K)}$ and $U_i=\operatorname{arctanh}m_i$.

The original Gaussian facts used here are the finite-depth transcript law
and conditional decomposition in \cite[Section 4]{brennecke2023operator}. For clarity, the
precise consequences needed below are listed here. With
\[
 M_0=(m^{(1)},\ldots,m^{(K)}),\qquad
 Z_0=(z^{(2)},\ldots,z^{(K+1)}),
\]
the limiting Gram matrices of $M_0$ and $Z_0$ coincide and are positive
definite at every fixed depth. The empirical joint law is its scalar
Gaussian transcript law, and
\begin{equation}\label{curv:old-map}
 T_0=Z_0(M_0^{\mathsf T}M_0)^{-1}M_0^{\mathsf T},\qquad
 G_{\mathcal C}:=GP+PG-PGP
       =\beta(T_0+T_0^{\mathsf T})+o_{\mathrm{op}}(1).
\end{equation}
The selected prior satisfies exactly
\begin{equation}\label{curv:old-generating-field}
             m_i^{(K)}=\tanh(h+\beta z_i^{(K)}).
\end{equation}
At fixed $K$, $N^{-1}\|m^{(K)}\|^2\to q$ and the normalized distance
between late consecutive iterates is at most $\varepsilon_K+o(1)$ in
probability, for a deterministic sequence $\varepsilon_K\to0$.
The original TAP residual and TAP value satisfy
\begin{equation}\label{curv:old-tap-errors}
 \frac{\|R_G(m)\|}{\sqrt N}\le\varepsilon_K+o_{\mathbb P}(1),\qquad
 \left|\frac{\mathcal F^G_{h\mathbf1}(m)}N+f_{\mathrm{RS}}\right|
       \le\varepsilon_K+o_{\mathbb P}(1),
\end{equation}
after increasing $\varepsilon_K$ if necessary, where
\[
 R_G(m)=U-Gm-h\mathbf1+\beta^2(1-q_m)m,
 \qquad
 f_{\mathrm{RS}}
   =\mathbb E\log\bigl(2\cosh(h+\beta\sqrt q Z)\bigr)
                         +\frac{\beta^2}{4}(1-q)^2.
\]
The residual assertion follows from the ordinary AMP recursion and the
vanishing distance between late consecutive iterates. Fixed-depth
equivalence to the canonical iteration is part of the original Gaussian
transcript result. To see the TAP-value assertion, use the residual
identity and Gaussian integration by parts to obtain, in the successive
limits,
$N^{-1}m^{\mathsf T}Gm\to2\beta^2q(1-q)$ and
$\mathbb E[(U-h)\tanh U]=\beta^2q(1-q)$, and substitute into
\eqref{curv:tap-functional}.

Each of the conditions in \eqref{curv:old-map}--\eqref{curv:old-tap-errors}
is transcript-measurable: the extra query $Gm$ is recorded. After fixing
the new depth and the tolerances used below, let $\mathcal G_{N,K}$ impose
these conditions, the corresponding whitened Gram bounds, the finite
families of weighted scalar tests specified in
\S\ref{curv:profiles}, and a fixed bound on $\|GP\|_{\mathrm{op}}$.
For every sufficiently large fixed $K$,
\begin{equation*}
             \mathbb P(\mathcal G_{N,K})\longrightarrow1.
\end{equation*}
No exponential probability is required in this assertion. The original
true mean and the actual partition function are not included in
$\mathcal G_{N,K}$; their cutoffs are handled separately under the true
law in \S\ref{curv:likelihood-transfer}.

Here is an explicit choice of the
whitened Gram conditions. At each fixed $K$, let
\[
 \Sigma_K=\lim_{N\to\infty}N^{-1}M_0^{\mathsf T}M_0
         =\lim_{N\to\infty}N^{-1}Z_0^{\mathsf T}Z_0,
 \qquad
 \widetilde M_0=M_0\Sigma_K^{-1/2},\quad
 \widetilde Z_0=Z_0\Sigma_K^{-1/2}.
\]
For a prescribed $0<\vartheta<1/2$, include in $\mathcal G_{N,K}$ the
two bounds
\[
 \left\|N^{-1}\widetilde M_0^{\mathsf T}\widetilde M_0-I_K\right\|_{\op}
 \le\vartheta,\qquad
 \left\|N^{-1}\widetilde Z_0^{\mathsf T}\widetilde Z_0-I_K\right\|_{\op}
 \le\vartheta.
\]
These hold with probability $1-o(1)$ at every fixed $K$. The raw
smallest eigenvalue of $\Sigma_K$ may affect how large $N$ must be to
enter this original event, but it does not enter the subsequent
conditional exponential rate. In fact,
\[
 T_0=\widetilde Z_0
       (\widetilde M_0^{\mathsf T}\widetilde M_0)^{-1}
       \widetilde M_0^{\mathsf T},\qquad
 \|T_0\|_{\op}\le\frac{1+\vartheta}{1-\vartheta},
\]
uniformly in $K$. This is a relative Gram estimate; an absolute
entrywise tolerance independent of $K$ for the raw Gram matrix would
not by itself suffice for the depth-independent rate. All $o_N(1)$ profile errors are understood at
fixed $K$. For each prescribed fixed accuracy the conditional estimates
below have a rate independent of $K$; no exponential estimate is
asserted for the original-good event itself.

\subsection{Exact conditional planting and the likelihood in TAP coordinates}
\label{curv:likelihood}

For any finite adaptive linear Gaussian transcript, Gaussian regression
gives, conditionally on that transcript,
\begin{equation}\label{curv:gaussian-regression}
                  G=G_{\mathcal C}+QWQ,
\end{equation}
where $W$ is a fresh GOE matrix with off-diagonal variance $\beta^2/N$.
Adaptivity causes no extra density factor: each input was chosen from
earlier answers, and conditioning successively restricts the Gaussian to
the recorded linear fibers. This is also the conditional representation
in \cite[Appendix A.1]{celentano2024sudakov}.

Define $\pi_m(x)=\prod_i e^{U_ix_i}/(2\cosh U_i)$. The reference law
$\mathsf R_K$ uses the original marginal of $\mathcal C_K$, and then draws
\begin{equation*}
 X\sim\pi_m,\qquad
 G=G_{\mathcal C}+Q\widetilde WQ+\frac{\beta^2}{N}QXX^{\mathsf T}Q,
\end{equation*}
with fresh $\widetilde W$. The shift preserves all the recorded answers.
Attach the same channel $Y_t=tX+B_t$ under both laws. Denote the actual law
of $(G,X,Y)$ by $\mathsf P$.

\begin{lem}[Exact likelihood identity]\label{curv:rn-lemma}
With $a_0=PX-m$, the likelihood ratio satisfies
\begin{align}
 \log\frac{d\mathsf P}{d\mathsf R_K}
  ={}&\frac12X^{\mathsf T}G_{\mathcal C}X+h\mathbf1^{\mathsf T}X
        +\frac{\beta^2}{4N}\|QX\|^4
        -\log\pi_m(X)-\log Z_G \label{curv:raw-rn}\\
  ={}&-\mathcal F^G_{h\mathbf1}(m)-\log Z_G
        -(X-m)^{\mathsf T}R_G(m)\notag\\
     &+\frac12a_0^{\mathsf T}Ga_0+a_0^{\mathsf T}GQX
       -\frac{\beta^2}{2}(1-q_m)\|a_0\|^2
       +\frac{\beta^2}{4N}(2a_0^{\mathsf T}m+\|a_0\|^2)^2.
 \label{curv:tap-rn}
\end{align}
These identities are unchanged by attaching independent auxiliary seeds
or any common observation channel.
\end{lem}

\begin{proof}
Relative to the unshifted conditional Gaussian in
\eqref{curv:gaussian-regression}, the shifted conditional matrix has
likelihood
\[
 \exp\left\{\frac12X^{\mathsf T}QGQX
                   -\frac{\beta^2}{4N}\|QX\|^4\right\}.
\]
Dividing the true Gibbs mass by this likelihood and by $\pi_m$ cancels
the residual interaction and gives \eqref{curv:raw-rn}. In particular,
the coefficient $\beta^2/(4N)$ follows from
$\operatorname{Var}(X^{\mathsf T}QWQX/2)
     =\beta^2\|QX\|^4/(2N)$.

For the second identity use $G_{\mathcal C}m=Gm$,
$QG_{\mathcal C}Q=0$, and
\[
 \tfrac12X^{\mathsf T}G_{\mathcal C}X
 =X^{\mathsf T}Gm-\tfrac12m^{\mathsf T}Gm
              +\tfrac12a_0^{\mathsf T}Ga_0+a_0^{\mathsf T}GQX,
 \qquad
 \|QX\|^2=N(1-q_m)-2a_0^{\mathsf T}m-\|a_0\|^2.
\]
Substitute these expressions, $I(m_i)=m_iU_i-\log(2\cosh U_i)$, and
$Gm+h\mathbf1-U=\beta^2(1-q_m)m-R_G(m)$ into
\eqref{curv:raw-rn}. Expanding the squared norm cancels the terms linear
in $a_0^{\mathsf T}m$ and gives \eqref{curv:tap-rn}.
\end{proof}

Consequently, on $\|G\|_{\mathrm{op}}\le M$ and
$\|PX-m\|\le\zeta\sqrt N$,
\begin{align}
 \frac1N\log\frac{d\mathsf P}{d\mathsf R_K}
 \le{}&\frac{-\mathcal F^G_{h\mathbf1}(m)-\log Z_G}{N}
          +\frac{2\|R_G(m)\|}{\sqrt N}\notag\\
       &+M(\zeta+\zeta^2/2)
          +\frac{\beta^2}{4}(2\zeta+\zeta^2)^2.
 \label{curv:rn-bound}
\end{align}
The negative term in \eqref{curv:tap-rn} has simply been discarded.

\subsection{Observable Gaussian completion}
\label{curv:completion}

Under both laws draw an independent GOE $W'$ at scale $\beta$ and form
the observable matrix
\begin{equation*}
             A_{\mathrm{obs}}
             =G-G_{\mathcal C}+PW'+W'P-PW'P.
\end{equation*}
Under $\mathsf R_K$, conditionally on $\mathcal C_K$,
\[
 W_*=Q\widetilde WQ+PW'+W'P-PW'P
\]
is a fresh full GOE independent of the product teacher. Indeed the two
summands are independent Gaussian blocks on orthogonal matrix subspaces
and their covariances sum to the full GOE covariance. Write $e=X-m$ and
introduce the analytical matrix
\begin{equation*}
                       A_*=W_*+\frac{\beta^2}{N}ee^{\mathsf T}.
\end{equation*}
Since $Qm=0$, we have the exact identities
\begin{equation}\label{curv:completion-exact}
 A_{\mathrm{obs}}=W_*+\frac{\beta^2}{N}Qee^{\mathsf T}Q,
 \qquad QA_{\mathrm{obs}}Q=QA_*Q,
 \qquad G=G_{\mathcal C}+QA_*Q.
\end{equation}
The following bound compares the AMP outputs from these two matrices:
\begin{equation}\label{curv:completion-error}
 \|A_{\mathrm{obs}}-A_*\|_{\mathrm{op}}
 \le\frac{\beta^2}{N}\|Pe\|(2\|e\|+\|Pe\|).
\end{equation}

Conditional on the transcript, the coordinates of $e$ are independent,
centered, and bounded by two. Hoeffding's inequality
\cite{boucheron2013concentration} for each direction
of a $1/2$-net in $\operatorname{ran}P$ gives
\begin{equation}\label{curv:product-projection}
 \mathsf R_K(\|Pe\|>\zeta\sqrt N\mid\mathcal C_K)
                  \le 2\cdot5^K e^{-c_0\zeta^2N},
\end{equation}
where $c_0>0$ is universal. For a fixed number of AMP steps, the difference
between the outputs on $A_{\mathrm{obs}}$ and $A_*$ is at most
$C_L\sqrt N\|A_{\mathrm{obs}}-A_*\|_{\mathrm{op}}$ on a fixed
matrix-operator-norm event. This follows by induction: the denoisers are
$1$-Lipschitz, all iterates are coordinatewise bounded, and the empirical
Onsager coefficient is $C/\sqrt N$-Lipschitz in an iterate. The exceptional
operator-norm event is exponentially small under the reference. Thus
completion stability has an exponential rate depending on $L$ and the
desired accuracy, but not on $K$.

\subsection{The centered-prior AMP and its finite-depth exponential coupling}
\label{curv:new-amp}

Put $d(t)=r(t)-q$. The centered scalar prior has mean zero, and for a
statistic $s=ve+\sqrt v Z$ its posterior mean is
\begin{equation*}
            \mathbb E[e\mid U,s]
              =\tanh(U+s+v m)-m.
\end{equation*}
This deterministic $vm$ shift is retained in the new AMP. With
$d_0(t)=a(t)-q$, $a^{-1}=0$, and $Y=Y_t$, define
\begin{align*}
 a^0&=\tanh(U+Y)-m,\\
 z^j&=Aa^j-\beta^2b_j a^{j-1},
      &b_j&=N^{-1}\sum_i[1-(m_i+a_i^j)^2],\\
 a^{j+1}&=\tanh(U+Y+z^j+\beta^2d_j(t)m)-m,
      &d_{j+1}(t)&=a\bigl(t+\beta^2d_j(t)\bigr)-q.
\end{align*}
The observable witness is the output $\widehat n_L=m+a^L$ with
$A=A_{\mathrm{obs}}$. An asterisk will denote the corresponding output
on $A_*$.

The scalar comparison fields are
\begin{equation*}
 z^j=\beta^2d_j(t)e+\beta H_j,
 \qquad (H_0,\ldots,H_{L-1})\sim N(0,D_L(t)),
 \qquad (D_L(t))_{ij}=d_{\min(i,j)}(t),
\end{equation*}
independently of the product teacher and observation noise. To verify the
covariance, couple the Gaussian statistics by one Brownian motion at the
times $\beta^2d_j(t)$. The corresponding conditional means are a
martingale with respect to the nested observation sigma fields. Therefore
the comparison input variables $a^j$ satisfy
\begin{equation}\label{curv:new-gram}
 \mathbb E[a^j\mid U]=0,\qquad
 \mathbb E[e a^j]=\mathbb E[(a^j)^2]=d_j(t),\qquad
 \mathbb E[a^ia^j]=d_{\min(i,j)}(t).
\end{equation}
The teacher loss of the comparison output is $1-q-d_L(t)$.

The scalar contraction in \eqref{curv:scalar-parameters} gives
\begin{equation}\label{curv:depth-error}
 0\le d(t)-d_j(t)\le\kappa^j,
 \qquad
 \mathbb E(a^{j+1}-a^j)^2=d_{j+1}(t)-d_j(t).
\end{equation}
For $t>0$, $d_0(t)>0$ and every finite increment is strictly positive,
since $a'(v)>0$ at finite $v$. Hence $D_L(t)\succ0$. Continuity and
compactness imply
\begin{equation}\label{curv:new-gram-uniform}
     \inf_{t\in[\delta,T_0]}\lambda_{\min}(D_L(t))>0
\end{equation}
at every fixed $L$.

We give the quantitative coupling step explicitly, because its constants
must be chosen before the original depth. The Gaussian coupling theorem
\cite[Theorem 4]{reeves2025dimension} applies to a finite generalized first-order
recursion driven by full GOE and to an arbitrary specified Gaussian
comparison with positive definite covariance. Its error is bounded by a
constant times its empirical mean/Onsager mismatch, $\sqrt N$ times its
empirical covariance mismatch, and $\sqrt s$, except on an event of
probability $3e^{-s}$, $0\le s\le N$. The constant depends only on the
number of steps, the map Lipschitz constants, and the extreme eigenvalues
of the comparison covariance.

Here condition on $\mathcal C_K,X,Y$, and use $W_*/\beta$ as the driving
GOE. In that theorem's recursion the matrix-multiplied map is $\beta a^j$
and the remaining map is
\[
        \frac{\beta^2}{N}e(e^{\mathsf T}a^j)
                        -\beta^2b_j a^{j-1}.
\]
These maps have dimension-free Lipschitz constants depending only on
$L,\beta$: the spike has operator norm at most $4\beta^2$, the denoisers
are bounded, and their empirical derivatives have the Lipschitz bound
used after \eqref{curv:product-projection}. The comparison mean is
$\beta^2d_j(t)e$ and the covariance is $\beta^2D_L(t)$.

In the notation of the cited theorem, its two mismatches are
\[
 \Delta_1=\|\mathcal M-\mathcal G-\mathcal F B_N\|_{2,1}
                      \|\Omega\|_{2,1}^{-1},\qquad
 \Delta_2=C_L^{\mathrm R}\sqrt N\|S_N-I\|_{\mathrm{op}},
\]
where $C_L^{\mathrm R}$ is the numerical depth-dependent factor in that
theorem, $\Omega^{\mathsf T}\Omega=\beta^2D_L(t)$,
\[
 B_N=N^{-1}(\beta^2D_L(t))^{-1}
                (\mathcal Y-\mathcal M)^{\mathsf T}\mathcal F,
 \qquad
 S_N=N^{-1}\Omega^{-\mathsf T}\mathcal F^{\mathsf T}
                                           \mathcal F\Omega^{-1}.
\]
The columns of $\mathcal F$ are $\beta a^j$ evaluated on the comparison.
Gaussian integration by parts gives the limiting nonzero Onsager entries
$\beta\mathbb E[1-(m+a^j)^2]$ in the preceding-input position. Their
contribution $\mathcal F B_N$ cancels the Onsager term of
$\mathcal G$, up to its empirical deviation. The remaining mean mismatch
is $\beta^2e\{N^{-1}e^{\mathsf T}a^j-d_j(t)\}$.
Thus both mismatches are controlled by finitely many empirical overlaps,
empirical derivatives, and products of a bounded denoiser with at most
two Gaussian factors.

Conditional on the old transcript, the comparison rows are independent,
although not identically distributed. Bounded products satisfy Hoeffding
bounds. Products with Gaussian factors have uniformly bounded
subexponential moment generating functions. Their conditional expectations
are the scalar quantities in \eqref{curv:new-gram}, up to the old empirical
profile errors included in $\mathcal G_{N,K}$. Choose these errors below a
fixed tolerance and apply the arbitrary-comparison theorem with
$s=c\epsilon^2N$. For every fixed $L$ and $\epsilon>0$ this gives
\begin{equation}\label{curv:coupling-rate}
 \sup_{t\in[\delta,T_0]}
 \mathsf R_K\left(\max_{j<L}
       N^{-1/2}\|z^{j,*}-z^{j,\mathrm{comp}}\|>\epsilon
                  \,\middle|\,\mathcal C_K\right)
       \le C_L e^{-c_{L,\epsilon}N},
 \qquad \mathcal C_K\in\mathcal G_{N,K}.
\end{equation}
The constants are independent of $K$. This use of the arbitrary-comparison
version is necessary: the specified scalar moments need not match after
conditioning on a particular realization of $X,Y$.

The same coupling also controls the observable teacher loss. The
comparison output has loss $1-q-d_L(t)$, and its empirical loss is an
average of independent bounded variables conditional on the old
transcript. Hoeffding's inequality, the $L^2$ coupling, and completion
stability therefore give, for every fixed $L$ and $\varepsilon>0$,
\begin{equation}\label{curv:teacher-loss}
 \sup_{t\in[\delta,T_0]}
 \mathsf R_K\left(
   \frac{\|\widehat n_L-X\|^2}{N}
       >1-q-d_L(t)+\varepsilon
              \,\middle|\,\mathcal C_K\right)
     \le \exp\{C(K+L)-c_{L,\varepsilon}N\}.
\end{equation}
The prescribed original empirical tolerances may be reduced for this
estimate. Its rate is independent of $K$.

\subsection{Old profiles, weighted Chernoff estimates, and uniform directions}
\label{curv:profiles}

The remaining tests involve a linear combination of all the old Gaussian
fields. We require constants independent of the original depth, even
though that depth remains fixed as $N\to\infty$.

In the canonical innovation construction, let $\phi_j$ be an adapted unit
vector orthogonal to the previously queried span. The normalized residual
innovation is
\[
          g^j=\beta^{-1}\sqrt N\,Q_{j-1}GQ_{j-1}\phi_j.
\]
Conditionally on the preceding transcript it is centered Gaussian with
covariance $Q_{j-1}+\phi_j\phi_j^{\mathsf T}$. In particular every
coordinate has conditional variance at most two. For deterministic
$\theta$ with $\sum_j\theta_j^2=1$, iterated conditioning gives
\begin{equation}\label{curv:innovation-mgf}
 \mathbb E e^{u\sum_j\theta_jg_i^j}\le e^{u^2},\qquad
 \mathbb E e^{c(\sum_j\theta_jg_i^j)^2}
                         \le(1-4c)^{-1/2},\quad c<1/4.
\end{equation}
The second inequality follows by integrating the first against an
independent standard Gaussian. The old effective fields are deterministic
linear combinations of these innovations. Whitening by their limiting
covariance makes the coefficient squared sum equal to one. Therefore
\eqref{curv:innovation-mgf} applies uniformly to every deterministic unit
direction in the whitened old-output block.

Explicitly, at each site $i$, put
$z_i=(z_i^{(2)},\ldots,z_i^{(K+1)})^{\mathsf T}$ and
$g_i=(g_i^1,\ldots,g_i^K)^{\mathsf T}$. The deterministic canonical
coefficients give $z_i=R_Kg_i$, with limiting covariance
$\Sigma_K=R_KR_K^{\mathsf T}$. For every unit $\theta\in\mathbb R^K$,
\[
 \theta^{\mathsf T}\Sigma_K^{-1/2}z_i
   =(R_K^{\mathsf T}\Sigma_K^{-1/2}\theta)^{\mathsf T}g_i,
 \qquad
 \|R_K^{\mathsf T}\Sigma_K^{-1/2}\theta\|_2^2=1.
\]
Thus the subgaussian constant is independent of $K$, even when
$\Sigma_K$ is ill-conditioned. The number of directions in the fixed
nets contributes the factor $\exp\{C(K+L)\}$, and does not reduce the
coefficient of $N$ in the conditional exponential bound.

Take, for example, $c=1/16$. The square of the exponential in
\eqref{curv:innovation-mgf} has uniformly bounded expectation. Fixed-depth
empirical Gaussian convergence, first for bounded truncations and then
using this uniform integrability, implies that on an old event of
probability $1-o(1)$,
\begin{equation}\label{curv:old-profile}
 \max_{\theta\in\mathcal N_K}
      \frac1N\sum_i e^{c f_{\theta,i}^2}\le C,
 \qquad |\mathcal N_K|\le e^{C'K}.
\end{equation}
Here $\mathcal N_K$ is any fixed-resolution sphere net and $f_\theta$ is
the corresponding whitened old Gaussian field. All constants in the
exponential bound are independent of $K$. The old-good event also imposes
the limiting joint Gaussian tests involving $f_\theta$ and the particular
old field in \eqref{curv:old-generating-field}. These are finite families
of tests at fixed $K$ and require no exponential original probability.
Uniformity in $t\in[\delta,T_0]$ follows by a fixed parameter net and the
continuity of the conditional scalar expectations; their growth in the
old Gaussian field is at most quadratic and is controlled by
\eqref{curv:old-profile}.

Here are the conditional exponential estimates for the weighted tests.
Let $V_i$ be independent centered variables with $|V_i|\le M$, conditional
on the old transcript, and set $w_i=f_i^2$. Taylor's formula and centering
give
\[
 \log\mathbb E e^{\lambda w_iV_i}
  \le \frac{\lambda^2M^2}{2}w_i^2e^{|\lambda|Mw_i}.
\]
If $N^{-1}\sum_i e^{cw_i}\le C$ and
$|\lambda|\le c/(2M)$, then
$w_i^2e^{|\lambda|Mw_i}\le C_c e^{cw_i}$. Hence
\begin{equation*}
 \log\mathbb E\exp\left(\lambda\sum_iw_iV_i\right)
             \le C_1N\lambda^2,
 \qquad
 \mathbb P\left(\left|N^{-1}\sum_iw_iV_i\right|>\epsilon
                          \,\middle|\,\mathcal C_K\right)
             \le2e^{-c_\epsilon N}.
\end{equation*}
The constants depend on $c,C,M,\epsilon$, not on $K$. A bounded function
of a Gaussian vector times one Gaussian coordinate is subgaussian after
centering, with uniform parameter. Thus mixed old-field/new-Gaussian
terms have conditional log-mgf at most
$C\lambda^2\sum_i f_i^2$. A bounded function times two new Gaussian
coordinates is subexponential with uniform parameter. Polarization
handles products of two different old directions.

For orthogonality to the old input span, take an orthonormal basis in
Euclidean norm. Conditional centering in \eqref{curv:new-gram} and
Hoeffding give an $e^{-c_\epsilon N}$ tail for every normalized projection
of a new bounded input. New Gaussian-output projections have the same
bound directly. A sphere net then controls the norm of the projection.
This argument requires only the squared coefficient norm, so the basis
vectors need no coordinatewise bound. Applied to both old input and
output spans, these estimates show that all new input and effective
output vectors have normalized projection at most $\epsilon$, except
with probability $e^{C(K+L)}e^{-c_\epsilon N}$.

We also record the truncation needed to transport weighted quadratic
tests through the $L^2$ coupling. For bounded diagonal weights
$s,\widehat s\in[0,1]^N$ and a Gaussian linear form $H$,
\begin{equation}\label{curv:weight-truncation}
 \frac1N\sum_i|s_i-\widehat s_i|H_i^2
 \le R^2\frac{\|s-\widehat s\|}{\sqrt N}
          +2\frac1N\sum_iH_i^2\mathbf1_{|H_i|>R}.
\end{equation}
The truncation radius $R$ is a fixed constant, chosen before the coupling
accuracy. To verify an exponential bound for the last term, let
$H_i=f_i+g_i$, with $f$ satisfying \eqref{curv:old-profile} and independent
Gaussian $g_i$ of bounded variance. For sufficiently small constants
$c'>v>0$,
\[
 \mathbb E[(e^{vH_i^2}-1)\mathbf1_{|H_i|>R}]
 \le e^{-(c'-v)R^2}\mathbb E e^{c'H_i^2}
 \le C_2e^{-c_2R^2}e^{cf_i^2}.
\]
Summing logarithmic moment generating functions proves that the average
square tail is arbitrarily small, except with probability $e^{-c_\epsilon
N}$, after choosing $R$ large. In particular no exponential tail bound
for an empirical Gaussian fourth moment is used.

For uniformity over directions, write $H_\theta=\mathsf H\theta$ and
choose $\theta_0$ in an $\epsilon$-net. Pointwise,
\[
 H_\theta^2\mathbf1_{|H_\theta|>R}
 \le4H_{\theta_0}^2\mathbf1_{|H_{\theta_0}|>R/2}
                       +4(H_\theta-H_{\theta_0})^2.
\]
The average of the last term is bounded by
$4\epsilon^2\|\mathsf H\|_{\mathrm{op}}^2/N$. The whitened-output Gram
bounds control this operator norm. Therefore a finite net extends all
the preceding quadratic estimates to every direction, at a cost
$e^{C(K+L)}$ in the prefactor. The same estimates apply to
$s/(1-\ell s)$ for any fixed $0\le\ell<1$, since that function and its
derivative are bounded on $[0,1]$.

\subsection{The second conditional regression and Onsager cancellation}
\label{curv:second-regression}

Condition on the old transcript, the teacher, the observation, and the
linear-query transcript of the analytical AMP on $A_*$. Denote this
information by $\mathcal D_{K,L}$. It determines the center
$n_L^*=m+a^{L,*}$ and the matrices
\[
 R_1=(a^{0,*},\ldots,a^{L-1,*}),\qquad
 S_1=W_*R_1,\qquad \Gamma_1=R_1^{\mathsf T}R_1.
\]
Indeed, the answers $S_1$ are obtained by subtracting the known spike
from the recorded products with $A_*$. Let $B_1$ be the new Onsager
matrix, so that its only nonzero entries are
$(B_1)_{j-1,j}=\beta^2 b_j$, and put $G_1=S_1-R_1B_1$.
Here rows and columns are indexed by $0,\ldots,L-1$.

Gaussian regression gives the exact conditional mean of $W_*$:
\begin{equation}\label{curv:new-exact-mean}
 \overline W_1
 =S_1\Gamma_1^{-1}R_1^{\mathsf T}
   +R_1\Gamma_1^{-1}S_1^{\mathsf T}
   -R_1\Gamma_1^{-1}(R_1^{\mathsf T}S_1)
                         \Gamma_1^{-1}R_1^{\mathsf T}.
\end{equation}
The inverses are used on the full-rank event implied by
\eqref{curv:new-gram-uniform} and the finite-depth coupling. The remaining
part is $Q_1W_{\mathrm f}Q_1$, where $P_1$ projects onto
$\operatorname{ran}R_1$, $Q_1=I-P_1$, and $W_{\mathrm f}$ is a fresh GOE
matrix at scale $\beta$, independent of $\mathcal D_{K,L}$.

The Stein identity for the Gaussian comparison is
\begin{equation}\label{curv:stein-gram}
 N^{-1}\bigl(R_1^{\mathsf T}G_1-B_1^{\mathsf T}\Gamma_1\bigr)
 \longrightarrow0.
\end{equation}
We use the following quantitative interpretation of this convergence.
For each fixed $L$ and each $\tau>0$, the probability that the operator
norm in \eqref{curv:stein-gram} exceeds $\tau$, conditional on an old
transcript satisfying its prescribed profile tolerances, is at most
$C_L\exp(-c_{L,\tau}N)$. The constants do not depend on $K$; the profile
tolerances and the threshold for $N$ may depend on $L$ and $K$,
respectively. To check the identity, column $j$ of $G_1$ differs from
$z^{j,*}-\beta^2d_j(t)e$ by
\[
 \beta^2e\bigl\{d_j(t)-N^{-1}e^{\mathsf T}a^{j,*}\bigr\}.
\]
The overlap tests in \eqref{curv:coupling-rate} control this difference.
For the comparison itself, Gaussian integration by parts gives
\[
 \mathbb E[a^i\beta H_j]
  =\beta^2\mathbb E[1-(m+a^i)^2]\,
       (D_L(t))_{i-1,j},\qquad i\ge1,
\]
and the expression is zero for $i=0$. These are exactly the entries of
$N^{-1}B_1^{\mathsf T}\Gamma_1$ in the scalar limit. Concentration of the
bounded--Gaussian products and the $L^2$ coupling prove the stated
fixed-tolerance bound.

Substituting $S_1=G_1+R_1B_1$ into
\eqref{curv:new-exact-mean} now gives
\begin{equation}\label{curv:new-mean-map}
 \overline W_1=T_1+T_1^{\mathsf T}+E_1,
 \qquad T_1=G_1\Gamma_1^{-1}R_1^{\mathsf T}.
\end{equation}
In fact,
\[
 E_1=-R_1\Gamma_1^{-1}
       \bigl(R_1^{\mathsf T}G_1-B_1^{\mathsf T}\Gamma_1\bigr)
                  \Gamma_1^{-1}R_1^{\mathsf T}.
\]
Thus, for every prescribed $\tau>0$, $\|E_1\|_{\mathrm{op}}\le\tau$
outside an event of probability $C_L e^{-c_{L,\tau}N}$ after decreasing
the preceding Gram tolerance. This uses the uniform positive
finite-depth bound in \eqref{curv:new-gram-uniform}. In particular,
the two Onsager terms cancel in the conditional mean; there is no
additional mixed correction involving the old transcript.

Combining this regression with \eqref{curv:completion-exact}, define
\begin{align}
 \overline G&:=\mathbb E_{\mathsf R_K}[G\mid\mathcal D_{K,L}]
       =G_{\mathcal C}+Q\overline W_1Q
                   +\frac{\beta^2}{N}Qee^{\mathsf T}Q,\notag\\
 G&=\overline G+QQ_1W_{\mathrm f}Q_1Q.\label{curv:conditional-G}
\end{align}
Both identities are exact. The approximation of $\overline W_1$ by the
two maps in \eqref{curv:new-mean-map} has the error just specified.

The projection estimates in \S\ref{curv:profiles} have the same
fixed-tolerance form: for each $\tau>0$, outside an event of conditional
probability at most $\exp\{C(K+L)-c_{L,\tau}N\}$,
\[
 \|QT_1Q-T_1\|_{\mathrm{op}}\le\tau,\qquad
 \sup_{\|v\|=1}
 \bigl|\|Q_1Qv\|^2-1+\|Pv\|^2+\|P_1v\|^2\bigr|\le\tau.
\]
The spike is treated in the same way using
\eqref{curv:product-projection}. Consequently the conditional mean is,
to any fixed accuracy, the sum of the original Gaussianized map, the
new Gaussianized map, and the unprojected spike. The new Gaussian output
block is $Z_1=G_1/\beta$; its limiting Gram matrix agrees with that of
$R_1$.

\subsection{Conditional expectation comparison and the covariance reduction}
\label{curv:sf}

For each unit vector $v$, write $b_v=Q_1Qv$. Conditional on both
transcripts, compare the centered Gaussian processes
\[
        X_v=b_v^{\mathsf T}W_{\mathrm f}b_v,
        \qquad
        Z_v=\frac{2\beta}{\sqrt N}\|b_v\|\,\xi^{\mathsf T}b_v,
        \qquad \xi\sim N(0,I_N).
\]
For $r=\|b_v\|$, $s=\|b_{v'}\|$, and $a=b_v^{\mathsf T}b_{v'}$,
the increment variance of $Z$ minus that of $X$ is
\[
       \frac{2\beta^2}{N}\{(r^2-s^2)^2+2(rs-a)^2\}\ge0.
\]
Sudakov--Fernique, with an arbitrary common deterministic offset indexed
also by $w$, therefore gives the conditional expectation upper bound for
the supremum of the negative Hessian by the supremum with $X_v$ replaced
by $Z_v$. This is the expectation form used in
\cite[Eq.~(4.5)]{brennecke2023operator}; it does not require a finite-sample stochastic
domination at every threshold. Open-cube index sets are treated by compact
exhaustion.

Both suprema are $C\beta/\sqrt N$-Lipschitz in their remaining standard
Gaussian coordinates, independently of the two transcript ranks.
Gaussian concentration \cite{boucheron2013concentration} consequently
gives an $e^{-cN u^2}$ tail above
the respective conditional mean. It remains to bound the auxiliary
supremum uniformly by a negative constant.

We spell out the covariance reduction. Scale the direction by
$V_i=\sqrt N v_i$, so $N^{-1}\sum_iV_i^2=1$. At the limiting empirical
level let $\mathsf M$ consist of the old inputs and new centered inputs,
and let $\mathsf Z$ consist of their Gaussian output blocks, with the
new outputs divided by $\beta$. Their Gram matrices agree, the two
blocks are orthogonal, and $\mathsf Z$ is jointly Gaussian. Therefore
\begin{equation}\label{curv:variance-one}
 H=\mathsf Z^{\mathsf T}
          (\mathbb E\mathsf M\mathsf M^{\mathsf T})^{-1}
                                      \mathbb E[\mathsf M V]
      +\left(1-\mathbb E[V\mathsf M^{\mathsf T}]
          (\mathbb E\mathsf M\mathsf M^{\mathsf T})^{-1}
                                      \mathbb E[\mathsf M V]\right)^{1/2}
                                                   \Xi
\end{equation}
is a centered Gaussian of variance one. Here $\Xi$ is independent of
the transcript Gaussian variables, teacher, and observation noise.
This statement holds for every possible joint law of $V$ with those
variables: the coefficients in \eqref{curv:variance-one} are numerical
moments of that law.

There is an indexing point in selecting its two relevant Gaussian
coordinates. The old coordinate must be $z^{(K)}$, the field that
actually generated the prior $m^{(K)}$ in
\eqref{curv:old-generating-field}. It is generally the penultimate
recorded output after the extra query $Gm^{(K)}$. Its paired input in
the covariance identity is $m^{(K-1)}$, which differs from $m^{(K)}$
by normalized norm at most $\varepsilon_K+o(1)$. The newly recorded
answer is retained among the other old Gaussian directions. Similarly,
the selected new output is the field that generated $a^{L,*}$; its paired
input is $a^{L-1,*}$. Equations~\eqref{curv:depth-error} and
\eqref{curv:coupling-rate} replace these by their common fixed-point
values with an error tending uniformly to zero as $L$ grows.

The reference teacher depends on the whole old transcript only through
$m^{(K)}=\tanh(h+\beta z^{(K)})$. The new Gaussian comparison block is
independent of the teacher and side channel. Conditional Gaussian
projection onto the selected old and new coordinates therefore leaves
a Gaussian independent of both selected coordinates, the teacher, and
the observation noise. Write the limiting selected coordinates as
$\sqrt q Z_0$ and $\sqrt d Z_1$. With
\begin{equation}\label{curv:scalar-coupling}
 \begin{gathered}
 U=h+\beta\sqrt q Z_0,\quad m=\tanh U,\quad
       \mathbb E[X\mid U]=m,\quad d=r(t)-q,\\
 n=\tanh\bigl(U+(t+\beta^2d)X
                       +\beta\sqrt d Z_1+\sqrt t Z_2\bigr),\quad
 a=n-m,\quad \epsilon=X-n,\quad s=1-n^2,
 \end{gathered}
\end{equation}
the reduction is
\[
 H=\frac{x_0}{\sqrt q}Z_0+\frac{x_1}{\sqrt d}Z_1
             +\sqrt{1-x_0^2/q-x_1^2/d}\,\Xi,
 \qquad x_0=\mathbb E[mV],\quad x_1=\mathbb E[aV].
\]
The square root is real since $\mathbb E[ma]=0$,
$\mathbb E[m^2]=q$, and $\mathbb E[a^2]=d$.
The scalar upper comparison for the negative full Hessian is consequently
\begin{align}
 \Psi_t(V)={}&2\beta\mathbb E[VH]
       +\beta^2\bigl(x_1+z\bigr)^2
       +2\beta^2\bigl(x_0+x_1\bigr)^2
       -\mathbb E[V^2/s]-\beta^2\mathbb E s,
 \qquad z=\mathbb E[\epsilon V],\quad \mathbb EV^2=1.
 \label{curv:scalar-comparison}
\end{align}
The first rank-one term comes from the residual planted spike and the
second from the unfrozen TAP term in \eqref{curv:hessian}.

\subsection{A finite-sample bound for the quadratic dual}
\label{curv:finite-dual}

We formulate explicitly the empirical approximation needed for the
Gaussian comparison. This separates it from the scalar matrix
calculation. In the scalar experiment \eqref{curv:scalar-coupling}, put
\[
 f=(m,a,\epsilon)^{\mathsf T},\qquad b=\mathbb E s,
 \qquad w_\ell=\frac{s}{1-\ell s},\qquad
 Q_\ell=\mathbb E[w_\ell ff^{\mathsf T}],\quad 0\le\ell<1.
\]
For $t\in[\delta,T_0]$, let
\[
 \mathcal X_t=
 \bigl\{x=(x_0,x_1,z)^{\mathsf T}:
             x_0^2/q+x_1^2/d(t)\le1,\ |z|\le2\bigr\},
\]
and associate to $x$ the variance-one Gaussian variable
\[
 H_x=\frac{x_0}{\sqrt q}Z_0+\frac{x_1}{\sqrt{d(t)}}Z_1
       +\sqrt{1-x_0^2/q-x_1^2/d(t)}\,\Xi.
\]
The weighted Gram matrix $Q_\ell$ is positive definite. Indeed, the
first two coordinates are independent in the weighted $L^2$ sense by
Lemma~\ref{long:lem:scalar}, and the conditional variance of the teacher
error is positive. This also follows directly from
\eqref{curv:weighted-gram} below when $\ell=0$ and from $w_\ell\ge s$.
By continuity, $\inf_{t\in[\delta,T_0]}\lambda_{\min}(Q_\ell)>0$.

Define the scalar quadratic dual by
\begin{align}
 \mathfrak D_{t,\ell}(x)
 ={}&-\ell-\beta^2b+\beta^2\mathbb E[w_\ell H_x^2]
       +\beta^2(x_1+z)^2+2\beta^2(x_0+x_1)^2\notag\\
 &-\bigl(x-\beta\mathbb E[w_\ell fH_x]\bigr)^{\mathsf T}
       Q_\ell^{-1}\bigl(x-\beta\mathbb E[w_\ell fH_x]\bigr).
 \label{curv:explicit-dual}
\end{align}
For every scalar direction $V$ in \eqref{curv:scalar-comparison}, its
moments $x=(\mathbb E[mV],\mathbb E[aV],\mathbb E[\epsilon V])$ belong
to $\mathcal X_t$ and
$\Psi_t(V)\le\mathfrak D_{t,\ell}(x)$. To see this, add
$\ell(\mathbb EV^2-1)+2\lambda^{\mathsf T}(x-\mathbb E[fV])$ to the
comparison functional and maximize pointwise over $V$. The terms
involving $V$ become
\[
 -\ell+\mathbb E\bigl[w_\ell(\beta H_x-\lambda^{\mathsf T}f)^2\bigr]
                +2\lambda^{\mathsf T}x.
\]
Minimizing in $\lambda$ gives \eqref{curv:explicit-dual}. The minimizing
multiplier is uniformly bounded for $t\in[\delta,T_0]$ and
$x\in\mathcal X_t$.

Using the exact conditional mean from \eqref{curv:conditional-G}, set
\begin{align*}
 \mathscr U_{N,K,L}(\rho)
 :=\sup_{\substack{\|v\|=1,\ w\in(-1,1)^N\\
                         \|w-n_L^*\|\le3\rho\sqrt N}}
 \bigg\{&v^{\mathsf T}\overline Gv
       +\frac{2\beta}{\sqrt N}\|Q_1Qv\|\,\xi^{\mathsf T}Q_1Qv
       +\frac{2\beta^2}{N}(v^{\mathsf T}w)^2\\
       &-\sum_i\frac{v_i^2}{1-w_i^2}
       -\beta^2(1-q_w)\bigg\},
\end{align*}
where $\xi$ is independent standard Gaussian. This is the auxiliary
supremum obtained by conditional Sudakov--Fernique.

\begin{lem}[Finite-sample dual reduction]\label{curv:finite-dual-lemma}
Fix $0<\ell<1$ and an error tolerance $\varepsilon>0$. There are
$L_0<\infty$ and $C_\varepsilon<\infty$, depending only on
$\beta,h,\delta,T_0,\ell,\varepsilon$, with the following property.
For each fixed $L\ge L_0$, one can choose a rate $c>0$ and prescribed
original-profile tolerances, all before choosing the original depth
$K$. For every sufficiently large fixed $K$, and all
$N\ge N_0(K,L,\varepsilon)$, these good original transcripts satisfy
\begin{equation}\label{curv:finite-dual-bound}
 \sup_{t\in[\delta,T_0]}
 \mathsf R_K\!\left(
   \mathscr U_{N,K,L}(\rho)>
      \sup_{x\in\mathcal X_t}\mathfrak D_{t,\ell}(x)
             +\varepsilon+C_\varepsilon\rho
           \,\middle|\,\mathcal C_K\right)
       \le \exp\{C(K+L)-cN\},\qquad 0\le\rho\le1.
\end{equation}
The rate $c$ and the constant $C$ are independent of $K$. For each such
fixed $K$, the prescribed original conditions hold with probability
$1-o(1)$.
\end{lem}

\begin{proof}
We give the coefficient reduction, the conditional concentration, and
the neighborhood perturbation separately. Write
$\mathbb E_N F=N^{-1}\sum_iF_i$ and $V_i=\sqrt N v_i$.

First replace the exact conditional mean by the maps in
\eqref{curv:old-map} and \eqref{curv:new-mean-map}, and remove their
small cross-projections. The fixed-tolerance estimates of
\S\ref{curv:second-regression} bound the resulting error uniformly
over $v$ by a prescribed number $\tau$. The same holds for replacing
$Qe$ by $e$ in the spike. Whiten the old and new input/output blocks
before parametrizing their linear combinations. The relative Gram
bounds imply that the coefficient vector of the resulting Gaussian
linear form has squared norm $1+O(\tau)$. Removing the projection of
$\xi$ onto the input spans costs at most
\[
 \frac{C}{\sqrt N}\bigl(\|P\xi\|+\|P_1\xi\|\bigr).
\]
The Gaussian norm bound controls this term at every prescribed
tolerance with probability $1-\exp\{C(K+L)-cN\}$.

The two selected output coordinates are the ones described in
\S\ref{curv:sf}. Their covariances with the Gaussian linear form are
the corresponding input moments. Replacing those inputs by $m$ and
$a^{L,*}$ costs at most the old consecutive-iterate error and the new
consecutive-iterate error, by Cauchy--Schwarz and the relative Gram
bounds. Choose $L_0$ so that the new error and
$d(t)-d_{L-1}(t)$ are uniformly small, and then require $K$ to be
sufficiently large for the old error. At the scalar level the remaining
Gaussian directions are independent of the two selected coordinates,
the teacher, and the observation noise. Their contribution is therefore
determined by its variance. This is precisely the reduction to $H_x$.
The approximate empirical constraints on $(x_0,x_1)$ can be projected
onto $x_0^2/q+x_1^2/d(t)\le1$ with an arbitrarily small error, uniformly
in $t$, since $d(t)$ is bounded away from zero on the compact interval.
The teacher-error moment is bounded by two.

It remains to justify this reduction uniformly over all coefficient
directions. Perform the pointwise square completion used to derive
\eqref{curv:explicit-dual}, inserting its bounded scalar minimizing
multiplier. Every coefficient that must be compared is an empirical
average of a bounded function times at most two Gaussian linear forms:
\[
 \mathbb E_N w_\ell,\qquad
 \mathbb E_N[w_\ell ff^{\mathsf T}],\qquad
 \mathbb E_N[w_\ell fH],\qquad
 \mathbb E_N[w_\ell H^2].
\]
Here the empirical $f$ is $(m,a^{L,*},X-n_L^*)^{\mathsf T}$.
For the comparison variables, the conditional means of these quantities
are the corresponding scalar moments, up to the prescribed original
profile errors and the late-iterate errors just described. Conditional
on the old transcript, the comparison rows are independent. The
weighted Chernoff bounds and Gaussian-product bounds of
\S\ref{curv:profiles} give an $e^{-cN}$ failure probability for each
test, with constants independent of $K$.

For the actual new transcript, use \eqref{curv:coupling-rate} and the
square-tail estimate \eqref{curv:weight-truncation}. Choose a fixed
truncation radius so that the tail contribution is less than a
prescribed fraction of $\varepsilon$, and then choose the coupling
tolerance. The tests are uniform on fixed-resolution nets in the
whitened coefficient sphere and in the compact set of scalar moments
and multipliers. Their total cardinality is at most
$\exp\{C(K+L)\}$. The Gram bounds and the same square-tail estimate
extend the bounds from these nets to every coefficient direction.
Thus the net dimension is $K+L+O(1)$; no net over the original
$N$-dimensional Hessian directions is required. All empirical dual
coefficients are within their prescribed scalar tolerances with the
probability in \eqref{curv:finite-dual-bound}.

Finally allow $\|w-n_L^*\|\le3\rho\sqrt N$.
The Onsager scalar and the TAP rank-one term change by at most
$C\beta^2\rho$ in a unit direction. For the entropy term, the
pointwise square completion replaces $w_\ell$ by
\[
   \frac{1-w_i^2}{1-\ell(1-w_i^2)}.
\]
This function is bounded and globally Lipschitz on $[-1,1]$. With the
already fixed truncation radius, \eqref{curv:weight-truncation} bounds
its effect on every dual coefficient by
$C_\varepsilon\rho$ plus the allocated tail error. This bound is
uniform over all such $w$, including coordinates close to $\pm1$.
Choose the finitely many profile, projection, and coupling tolerances
so that their total contribution is at most $\varepsilon$.
The resulting exponential rate depends on these tolerances and $L$,
but not on $K$. All requirements on the old transcript are the
relative Gram and weighted empirical conditions in
\S\ref{curv:old-transcript} and \S\ref{curv:profiles}; they hold
with probability $1-o(1)$ at every sufficiently large fixed $K$.
\end{proof}

\subsection{The scalar gap and the curvature neighborhood}
\label{curv:neighborhood}

We connect the matrix inequality of \S\ref{sec:longitudinal} to
the scalar comparison, retaining the teacher and longitudinal modes.
Define its Gram variables by
\[
 A=\mathbb E[sm^2],\quad B=\mathbb E[sma],\quad C=\mathbb E[sa^2],
 \quad u=\mathbb E[s^2],\quad b=\mathbb E s.
\]
Write $Q_2=\bigl(\begin{smallmatrix}A&B\\B&C\end{smallmatrix}\bigr)$ and
recall $f=(m,a,\epsilon)^{\mathsf T}$. The Bayes identity
$\mathbb E[X\mid U,n]=n$ gives
\begin{equation}\label{curv:weighted-gram}
 b=A+2B+C+u,\qquad
 \mathbb E[sff^{\mathsf T}]
   =\begin{pmatrix}A&B&0\\B&C&0\\0&0&u\end{pmatrix}.
\end{equation}
In particular the last coordinate is orthogonal to the first two under
the weighted inner product. Gaussian integration by parts, including
the score $X-m$ from the conditional law of $X$ when differentiating
$U$, yields
\begin{equation}\label{curv:ibp-first}
 \mathbb E\left[s
    \begin{pmatrix}m\\a\\\epsilon\end{pmatrix}
    \begin{pmatrix}Z_0&Z_1\end{pmatrix}\right]
 =
 \beta\begin{pmatrix}
  \sqrt q(b-3A-B)&-2\sqrt d(A+B)\\
  \sqrt q(-4B-2C)&\sqrt d(u-2B-2C)\\
  0&-\sqrt d\,u
 \end{pmatrix}.
\end{equation}
For completeness, with $R_{ij}=\mathbb E[sZ_iZ_j]$, the corresponding
second integration-by-parts identities are
\begin{align}
 R_{00}&=b-2\beta^2q(b-3A-3B-C),\notag\\
 R_{01}&=-2\beta^2\sqrt{qd}(b-3A-5B-2C),\notag\\
 R_{11}&=b+\beta^2d(4b-6u).
 \label{curv:ibp-second}
\end{align}
These formulas follow by differentiating $\tanh U$ and $\tanh V$,
where $V$ is the final field in \eqref{curv:scalar-coupling}, and then
using $\mathbb E[\epsilon\mid U,n]=0$ and
$\mathbb E[\epsilon^2\mid U,n]=s$. The score in the $Z_0$ derivative
is essential because the teacher is drawn conditionally on $U$.

For example, Gaussian integration
by parts gives
$\mathbb E[Z_0F]=\beta\sqrt q\,\mathbb E[\partial_UF+(X-m)F]$
for a test $F=F(U,X,Z_1,Z_2)$. Since $\partial_Un=s$ and
$\partial_Um=1-m^2$ with $X$ held fixed, this gives
\[
 \mathbb E[smZ_0]
 =\beta\sqrt q\,\mathbb E\bigl[-2nsm+s(1-m^2)+(X-m)sm\bigr]
 =\beta\sqrt q\,(b-3A-B).
\]
The last equality uses $\mathbb E[X\mid U,n]=n$ and $a=n-m$.
This verifies explicitly the teacher-score contribution in the first
entry of \eqref{curv:ibp-first}.

The two positive margins are $1-\beta^2u$ and
$1-\beta^2(u+2B+C)$. The first is at least $1-\kappa$ and the second
is at least $1-\alpha$, since
$u+2B+C=\mathbb E[(1-m^2)(1-n^2)]
                 \le\mathbb E(1-m^2)^2$.

We evaluate \eqref{curv:explicit-dual} at $\ell=0$. Put
$x_{01}=(x_0,x_1)^{\mathsf T}$ and
\[
 L_{\mathrm d}=(1-\beta^2u)I_2+\beta^2
       \begin{pmatrix}
       2A-B-C&2A+2B\\
       4B+2C&2B+2C
       \end{pmatrix}.
\]
Equations~\eqref{curv:weighted-gram}--\eqref{curv:ibp-first} give exactly
\begin{equation*}
 x-\beta\mathbb E[sfH]
       =\begin{pmatrix}L_{\mathrm d}x_{01}\\z+\beta^2u x_1\end{pmatrix}.
\end{equation*}
Let $M,K$ be the two symmetric matrices in
\eqref{long:eq:F}, with its parameter $\lambda=\beta^2$. Because
$\Xi$ is independent and the variance coefficients of $H$ sum to one,
\eqref{curv:ibp-second} gives
\begin{equation}\label{curv:quadratic-H}
 \beta^2\mathbb E[sH^2]-\beta^2b
       =\beta^4x_{01}^{\mathsf T}
          \left(K+u\begin{pmatrix}0&0\\0&1\end{pmatrix}\right)x_{01}.
\end{equation}
The spike terms expand as
\[
 \beta^2(x_1+z)^2+2\beta^2(x_0+x_1)^2
       =\beta^2x_{01}^{\mathsf T}Mx_{01}
                            +2\beta^2x_1z+\beta^2z^2.
\]
The last coordinate in the negative square of the dual is
$-(z+\beta^2u x_1)^2/u$. It cancels the mixed term
$2\beta^2x_1z$ and the extra $\beta^4u x_1^2$ in
\eqref{curv:quadratic-H}. Consequently
\begin{equation*}
 \mathfrak D_{t,0}(x)=
  -\frac{1-\beta^2u}{u}\,z^2+x_{01}^{\mathsf T}F x_{01},
 \qquad
 F=\beta^2M+\beta^4K-L_{\mathrm d}^{\mathsf T}Q_2^{-1}L_{\mathrm d}.
\end{equation*}
This is precisely the matrix in \eqref{long:eq:F}. The algebraic
negativity proved in \S\ref{sec:longitudinal} and the two margins
above give strict negativity of its nonzero moment part.

The coefficients of this quadratic form vary continuously with $\ell$.
Turning on a sufficiently small positive $\ell$ preserves their uniform
negative definiteness. The zero-moment constant becomes
\[
 -\beta^2b-\ell+\beta^2\mathbb E\frac{s}{1-\ell s}
 =-\int_0^\ell
       \left(1-\beta^2\mathbb E\frac{s^2}{(1-vs)^2}\right)dv<0.
\]
By continuity, the quantitative longitudinal margin, and compactness in
$t\in[\delta,T_0]$, there are $\ell>0$ and $g>0$ such that
\begin{equation*}
 \sup_{t\in[\delta,T_0]}\sup_{x\in\mathcal X_t}
                          \mathfrak D_{t,\ell}(x)\le-g.
\end{equation*}
In particular, $\Psi_t(V)\le-g$ for every admissible scalar direction.

Apply Lemma~\ref{curv:finite-dual-lemma} with error tolerance less than
$g/8$, and then choose $\rho>0$ so that its neighborhood term is less
than $g/8$. Choose the new depth $L$ sufficiently large for that lemma
and for the desired teacher-loss accuracy. It follows that
$\mathscr U_{N,K,L}(\rho)\le-g/2$, except with conditional probability
$\exp\{C(K+L)-cN\}$. The radius is chosen using the scalar gap and the
bounded-weight truncation constants; it does not depend on $K$.

For precision in passing to conditional expectations, first obtain this
failure bound jointly over the new transcript and the auxiliary Gaussian
$\xi$, conditional on a good old transcript. Conditional Markov at
threshold $e^{-cN/2}$ discards at most
$e^{C(K+L)}e^{-cN/2}$ of the new transcripts. On the remaining good new
transcripts, the positive part of the auxiliary supremum is bounded by
$C+C\|\xi\|/\sqrt N$, with a uniform second-moment bound. Its conditional
expectation is therefore at most $-g/3$ for all sufficiently large $N$.
Conditional Sudakov--Fernique and Gaussian concentration from
\S\ref{curv:sf} now imply, with an exponential rate independent of
the old depth, positive curvature throughout the neighborhood of
$n_L^*$. Finally \eqref{curv:completion-error}--\eqref{curv:product-projection}
place the observable center within $\rho\sqrt N$ of $n_L^*$. After
reducing constants, this proves
\begin{equation}\label{curv:reference-neighborhood}
 \mathsf R_K\left\{
  \inf_{\substack{w\in(-1,1)^N:\|w-\widehat n_L\|\le2\rho\sqrt N}}
       \lambda_{\min}\mathcal H_G(w)<2\gamma
                    \,\middle|\,\mathcal C_K\right\}
          \le e^{C(K+L)}e^{-cN}.
\end{equation}
All tolerances and $L$ are chosen before making $K$ large.

\subsection{Likelihood transfer, mean placement, and the order of constants}
\label{curv:likelihood-transfer}

We first justify the true teacher cutoff at a fixed original depth.
Let $\overline m=\mathbb E[X\mid J]$. The original finite-depth scalar
products established in \S\ref{sec:initializer-proof}, the AMP equivalence in
\S\ref{sec:amp}, and the fixed-depth equivalence of the canonical
and ordinary iterations in \cite[Section 4]{brennecke2023operator} give
\begin{equation}\label{curv:fixed-depth-mean-error}
 \frac{\|m^{(K)}-\overline m\|^2}{N}
                  \longrightarrow 2(q-\varrho_{K-1})
                  \quad\text{in probability},\qquad
                         \varrho_j\uparrow q.
\end{equation}
Here $\varrho_j$ is the original cavity-overlap sequence denoted by
$\rho_j$ in the initializer section; this different letter avoids
confusion with the neighborhood radius. The indexing is
$m^{(1)}=\sqrt q\,\mathbf1$, so $K-1$ is the number of denoising steps.
In particular the assertion is concentration at each fixed depth,
stronger than a double-limit expectation estimate.

The original overlap estimate gives
$\mathbb E\operatorname{Tr}(C_0^2)\le C N$. For the adaptive rank-$K$
projection $P$, condition on the full $G$ and use the covariance identity:
\begin{equation}\label{curv:true-projection}
 \frac1N\mathbb E\|P(X-\overline m)\|^2
    =\frac1N\mathbb E\operatorname{Tr}(PC_0)
    \le\frac K N\mathbb E\sqrt{\operatorname{Tr}(C_0^2)}
    \le\frac{CK}{\sqrt N}.
\end{equation}
Given any fixed $\zeta>0$, choose fixed $K$ with
$2(q-\varrho_{K-1})<\zeta^2/16$. Equations
\eqref{curv:fixed-depth-mean-error}--\eqref{curv:true-projection} imply
\begin{equation}\label{curv:true-cutoff}
                  \mathsf P(\|PX-m\|>\zeta\sqrt N)=o(1).
\end{equation}
Thus the teacher cutoff is established under the true law.

We now specify the constants in their required order. From the compact
scalar gap first choose $\ell,g,\rho,\gamma$ and the fixed truncation
and moment tolerances. Fix a desired mean error $\eta<\rho/2$. Choose a
new finite $L$ so that both the consecutive-iterate error and
$\sup_{t\in[\delta,T_0]}(d(t)-d_L(t))$ are sufficiently small, the
latter smaller than $\eta^2/8$. The finite-depth coupling, weighted
Chernoff estimates, completion stability, and conditional spectral
concentration then give a rate $c>0$ independent of $K$.

Next choose $\zeta>0$ sufficiently small and increase the fixed original
depth $K$ until the original consecutive-iterate, residual, TAP-value,
and mean errors are below the corresponding tolerances, and the bound
in \eqref{curv:rn-bound} is less than $c/2$. The original free-energy
convergence $N^{-1}\log Z_G\to f_{\mathrm{RS}}$, the bounded GOE
operator norm, \eqref{curv:old-tap-errors}, and
\eqref{curv:true-cutoff} show that
\begin{equation*}
       \mathsf P\left(\log\frac{d\mathsf P}{d\mathsf R_K}>cN/2\right)
                          \longrightarrow0.
\end{equation*}
The original overlap and pressure inputs follow from
\cite[Theorem~1.1]{kusuoka2026quantitative}; Gaussian concentration \cite[Section~5.4]{boucheron2013concentration}
upgrades expected-pressure convergence to convergence in probability.
The independent diagonal changes the normalized pressure by $o(1)$.

If $E$ is the curvature failure in \eqref{curv:reference-neighborhood}
or the teacher-loss failure in \eqref{curv:teacher-loss}, split off $\mathcal G_{N,K}^{\mathrm c}$ and the
likelihood cutoff. On the remaining event,
\begin{equation}\label{curv:transfer}
 \mathsf P(E)
 \le o(1)+e^{cN/2}
       \mathsf R_K(E\cap\mathcal G_{N,K})
 \le o(1)+e^{C(K+L)-cN/2}=o(1).
\end{equation}
The reference exponential estimate therefore absorbs the likelihood
cutoff.

For mean placement, the observable output $\widehat n_L$ is a function
of $J,Y_t$ and independent auxiliary seeds. It is coordinatewise bounded
by one. Posterior Bayes orthogonality and
\S\ref{sec:scalar-law} give
\begin{equation*}
 \mathbb E\left[
    \left\{\frac{\widehat n_L^{\mathsf T}(X-m_t)}N\right\}^2\right]
       \le\frac1N\mathbb E\|C_t\|_{\mathrm{op}}
       \le\frac{\sqrt{\mathbb E\operatorname{Tr}(C_t^2)}}N=o(1),
\end{equation*}
uniformly over fixed times in the compact interval. The transferred
teacher loss is at most $1-q-d_L(t)+\eta^2/8$ with probability $1-o(1)$.
Using $N^{-1}\|m_t\|^2\to r(t)$, expand the two squared distances to get
\[
 \frac{\|\widehat n_L-m_t\|^2}{N}
 =\frac{\|\widehat n_L-X\|^2}{N}-1
         +\frac{\|m_t\|^2}{N}
         +\frac{2\widehat n_L^{\mathsf T}(X-m_t)}N
 \le d(t)-d_L(t)+\eta^2/8+o_{\mathbb P}(1)<\eta^2
\]
with probability $1-o(1)$. This places the observable center within
$\eta\sqrt N$ of the true posterior mean.

The ball of radius $\rho\sqrt N$ around $m_t$ is then contained in the
ball of radius $2\rho\sqrt N$ around $\widehat n_L$. Combining
\eqref{curv:reference-neighborhood}, \eqref{curv:transfer}, and
\eqref{curv:diagonal} gives \eqref{curv:actual-neighborhood}. A finite
union covers $\mathcal T$. The resulting event in
\eqref{curv:actual-neighborhood} is measurable with respect to the hollow
disorder and actual observations alone; hence the diagonal and completion
seeds can be integrated out. This completes the proof of
Theorem~\ref{curv:fixed-anchor}.


\section{The terminal convex cap}
\label{sec:terminal-proof}

Write $W_{ij}=g_{ij}/\sqrt N$ for $i<j$, $W_{ji}=W_{ij}$ and $W_{ii}=0$,
and put $J=\beta W$. For $\tau\in\{-1,1\}^N$ and $\rho>0$, set
\[
 K_\tau(\rho)=\left\{m\in[-1,1]^N:
       \sum_{i=1}^N(1-\tau_i m_i)\le\rho N\right\}.
\]
This is a compact convex set. Use the TAP functional $ F_{\tap}(m,a)$ defined above; in this section it is also written as $\mathcal F_a(m)$.

\subsection{Uniform convexity on every cap}

The following sparse-block estimate gives a better cap radius than a
direct active-coordinate decomposition.

\begin{lem}[Sparse quadratic forms control a weighted matrix]
Suppose $J$ is symmetric and
\[
 \max_{|S|\le2k}\|J_{S,S}\|_{\op}\le b
\]
for an integer $k\ge1$. If $D=\operatorname{diag}(d_1,\ldots,d_N)$,
$0\le d_i\le1$, and $\sum_i d_i\le k$, then
\[
 \|D^{1/2}JD^{1/2}\|_{\op}\le4b.
\]
\end{lem}

\begin{proof}
Fix $x$ with $\|x\|_2=1$ and let $v=D^{1/2}x$. Then
$\|v\|_2\le1$ and
$\|v\|_1\le(\sum_i d_i)^{1/2}\|x\|_2\le\sqrt k$.
Order the coordinates by decreasing $|v_i|$, and partition them into
successive blocks $S_1,S_2,\ldots$, each of size $k$ except possibly the
last. Write $v_j=v\mathbf1_{S_j}$. For every $j\ge2$,
\[
 \|v_j\|_2\le\frac{\|v_{j-1}\|_1}{\sqrt k},
\]
because every coordinate in $S_j$ is at most the average absolute
coordinate in $S_{j-1}$. Hence
\[
 \sum_j\|v_j\|_2\le\|v\|_2+\frac{\|v\|_1}{\sqrt k}\le2.
\]
The union of any two blocks has at most $2k$ coordinates, so
$|v_j^{\mathsf T}Jv_\ell|\le b\|v_j\|_2\|v_\ell\|_2$.
Summing gives $|v^{\mathsf T}Jv|\le4b$. Taking the supremum over $x$
proves the claim.
\end{proof}

\begin{prop}[Deterministic convex-cap lemma]
Let $0<\rho<1$ and $k=\lceil2\rho N\rceil$. Suppose
\[
 \max_{|S|\le2k}\|J_{S,S}\|_{\op}\le\frac1{16},
 \qquad 4\beta^2\rho\le\frac14.
\]
Then, simultaneously for every sign vector $\tau$, every field $a$, and
every $m\in K_\tau(\rho)\cap(-1,1)^N$,
\[
 \nabla^2\mathcal F_a(m)
   \succeq\left(\frac34-4\beta^2\rho\right)
                   \operatorname{diag}\bigl((1-m_i^2)^{-1}\bigr)
   \succeq\frac12 I_N.
\]
In particular, $\mathcal F_a$ is $1/2$-strongly convex on the entire
closed convex cap $K_\tau(\rho)$.
\end{prop}

\begin{proof}
Let $d_i=1-m_i^2$ and $D=\operatorname{diag}(d_i)$. Since
\[
 d_i=(1-\tau_i m_i)(1+\tau_i m_i)
       \le2(1-\tau_i m_i),
 \qquad \sum_i d_i\le2\rho N\le k,
\]
the preceding lemma gives
$\|D^{1/2}JD^{1/2}\|_{\op}\le1/4$. Moreover,
\[
 \left\|\frac{2\beta^2}{N}
    D^{1/2}mm^{\mathsf T}D^{1/2}\right\|_{\op}
 =\frac{2\beta^2}{N}\sum_i d_i m_i^2
 \le\frac{2\beta^2}{N}\sum_i d_i\le4\beta^2\rho.
\]
The remaining Onsager Hessian term $\beta^2(1-q_m)D$ is positive
semidefinite. Thus
\[
 D^{1/2}\nabla^2\mathcal F_a(m)D^{1/2}
       \succeq\left(\frac34-4\beta^2\rho\right)I_N.
\]
Conjugating by $D^{-1/2}$ proves the Hessian claim. The strong convexity
inequality extends to boundary points by continuity: approximate every
point of the cap by its convex combination with
$(1-\rho/2)\tau$, which lies in the open cube and in the cap whenever
$0<\rho<2$. The application below has $\rho<1$.
\end{proof}

\begin{cor}[A simultaneous high-probability SK event]
Fix $\beta>0$. Choose $0<\delta\le1/2$ so small that
\[
 \delta\log(e/\delta)\le\frac1{8192\beta^2},
 \qquad \rho=\frac\delta8.
\]
For $N\delta\ge4$, the convex-cap conclusion holds with probability at
least $1-2\exp(-N/(8192\beta^2))$, simultaneously for every sign vector
$\tau$ and every field $a$. If desired, the event $\|W\|_{\op}\le3$ can
be imposed as well, at additional failure probability $2e^{-N/4}$.
The permitted $\rho$ can be chosen of order
$[\max(1,\beta^2)\log(e+\beta)]^{-1}$ up to a universal small constant.
\end{cor}

\begin{proof}
For each fixed nonempty $S$ with $|S|\le\delta N$, the Gaussian
matrix norm estimate \cite[Section~3.3]{bandeira2026mixing}, together with
Gaussian concentration \cite[Section~5.4]{boucheron2013concentration}, gives
\[
 \bbP\left(\|W_{S,S}\|_{\op}>2\sqrt{|S|/N}+u\right)
       \le2e^{-Nu^2/4}.
\]
The hollow-matrix version follows by averaging an added independent
centered Gaussian diagonal: Jensen's inequality bounds the expected
maximum eigenvalue of either sign by its GOE counterpart, and the
Gaussian Lipschitz constant has the same upper bound.

Our choice of $\delta$ implies $\delta\le1/(8192\beta^2)$ and hence
$2\sqrt\delta\le1/(32\beta)$. Taking $u=1/(32\beta)$, and then a union
bound over at most $\exp[N\delta\log(e/\delta)]$ subsets, gives
\[
 \bbP\left(\max_{|S|\le\delta N}\|J_{S,S}\|_{\op}>\frac1{16}\right)
 \le2\exp\left[-\frac N{4096\beta^2}
                +N\delta\log(e/\delta)\right]
 \le2e^{-N/(8192\beta^2)}.
\]
Since $k=\lceil\delta N/4\rceil$, the assumption $\delta N\ge4$ ensures
$2k\le\delta N$. Also,
$4\beta^2\rho=\beta^2\delta/2\le1/16384<1/4$.
The deterministic proposition applies. The full-matrix bound is the same
Gaussian inequality with $S=[N]$ and $u=1$.
\end{proof}

\subsection{The true posterior lies in a cap after a fixed time}

The fixed-sign containment estimate below is independent of SK and of the AT condition. Together with the preceding simultaneous convexity event, it proves Proposition~\ref{prop:terminal-cap}.

\begin{cor}[A single fixed cap works for the entire terminal phase]
For any prior on the cube, let $Y_t=tX+B_t$ and $m_t=\bbE[X\mid Y_t]$. Set $t_0=2\log(8/\rho)$ and freeze $\tau_0=\sign Y_{t_0}$. Then
\[
 \bbP\left(\exists t\ge t_0:m_t\notin K_{\tau_0}(\rho)\right)
      \le e^{-\rho N/4}.
\]
\end{cor}

\begin{proof}
Put $S_0=d_H(X,\tau_0)$. Conditional on the environment, this is binomial
with parameters $N$ and $p_0=\Phi(-\sqrt{t_0})\le\rho/8$.
Let $\mathcal F_t$ be the observation filtration, including the environment when one is present.
For $t\ge t_0$, the fixed sign vector $\tau_0$ is $\mathcal F_t$-measurable,
and
\[
 \sum_i(1-\tau_{0,i}m_{t,i})
   =2\bbE[S_0\mid\mathcal F_t].
\]
The right side is twice a continuous bounded martingale. For every finite $T\ge t_0$, conditional Jensen and Doob's maximal inequality \cite[Section~1.3.A]{karatzas1991brownian} give
\[
 \bbP\left(\sup_{t_0\le t\le T}\bbE[S_0\mid\mathcal F_t]>\rho N/2\right)
 \le e^{-\rho N/2}\bbE e^{S_0}
 \le\exp\{-\rho N/2+p_0(e-1)N\}
 \le e^{-\rho N/4}.
\]
Let $T\to\infty$. The proof also holds conditionally on an environment determining the prior.
\end{proof}


\section{Numerical implementation}\label{sec:numerical}

This appendix specifies a finite-precision version of the algorithm and verifies that its additional error in the final spin law is $o(1)$. All constants below may depend on the fixed pair $(\beta,h)$ and the chosen anchor record. The goal is a deterministic polynomial bound, rather than an optimized exponent.

\subsection{Input access and the fixed record}

We implement the algorithm on a symmetric dyadic approximation of $J$ with a prescribed polynomial input length. For an exact real Gaussian input, the corresponding access convention supplies an entry to $b$ binary digits at cost polynomial in $b$ and $\log N$. Reading these bits and all subsequent arithmetic are included in the computation. The same convention applies to the fixed parameters when their numerical values are supplied through precision access. A fixed valid record contains positive curvature and radius lower bounds, fixed AMP depths, the terminal time, and the finitely many anchors. Its existence is proved in \S\ref{sec:anchor-domains}. These quantities are fixed before $N$ varies. The strict inequalities allow rational lower bounds and a rational anchor mesh. Increase the terminal time slightly to a rational number larger than $2\log(8/\rho)$ if necessary; Proposition~\ref{prop:terminal-cap} remains valid.

This fixed-record formulation suffices for the quantifiers in Theorem~\ref{thm:main}. We do not need a procedure that detects a correct neighborhood or tests a hidden structural event. In particular, unsuccessful inputs still pass through bounded, prescribed routines.

Let $J'$ be a symmetric hollow approximation with $\max_{i,j}|J'_{ij}-J_{ij}|\le\delta_N$. For every spin configuration,
\[
 |H_{J'}(x)-H_J(x)|\le N^2\delta_N/2,
 \qquad \TV(\mu_{J'},\mu_J)\le e^{N^2\delta_N}-1.
\]
The same entrywise error changes the TAP objective uniformly on the cube by at most $N^2\delta_N/2$. Fixed-depth AMP is polynomially Lipschitz in the matrix entries on the input norm guard. Choosing $\delta_N=N^{-D}$ with a sufficiently large fixed $D$ makes these perturbations smaller than the objective tolerance and anchor slack. Implement the norm guard by rejecting when $\sum_{i,j}(J'_{ij})^2>(N^2-1)^2$, an exact rational comparison. If $N\delta_N<1$, passing implies $\|J\|_F<N^2$, and rejection implies $\|J\|_F>N^2-2$. Its exceptional original-disorder probability remains $O_\beta(N^{-3})$. The supplied input length and all precision requests have fixed polynomial caps; an unsuccessful request returns uniform spins.

To apply the disorder estimates, first keep the exact target $\Pi_J$ and its Hamiltonian $H_J$, while using $J'$ in every deterministic proposal routine. The preceding perturbation bounds preserve the anchor and objective guarantees with slightly smaller slack, so the proof in \S\ref{sec:completion} applies to this proposal under the original Gaussian disorder. In its log weights, subsequently replacing $H_J$ by $H_{J'}$ changes each score by at most $N^2\delta_N/2$. Since $z\mapsto\min(1,e^z)$ is $1$-Lipschitz, a common-uniform coupling changes the chain's output with probability at most $k_NN^2\delta_N=o(1)$. This argument uses the original disorder law throughout; no Gaussian-disorder assertion is needed for the rounded matrix.

\subsection{Deterministic scalar and drift routines}

The parameter $q$ is computed once by one-dimensional Gaussian quadrature and bisection for \eqref{eq:original-fixed-point}. To see the local inverse bound, put $F(s)=s-\bbE\tanh^2(h+\beta\sqrt s\,Z)$. Gaussian differentiation at the unique root gives, with $U=h+\beta\sqrt q\,Z$,
\[
 F'(q)=1-\beta^2\bbE[\sech^2U(1-3\tanh^2U)]
       \ge1-\beta^2\bbE\sech^4U>0.
\]
For a fixed $(\beta,h)$, boundedness away from zero off a neighborhood of the root and this strict local derivative bound give polynomial-time computation to any inverse-polynomial accuracy. Truncation of the Gaussian integral and deterministic quadrature may be chosen to the same accuracy. The observed overlap $r(t)$ is used to prove that a fixed record exists; it is not evaluated at every fine-grid point.

Before computing an anchor or a TAP approximation, round the supplied observation coordinatewise on a fixed lattice of mesh $N^{-D_0}$, where $D_0$ is a sufficiently large fixed exponent. Fix all iteration counts, elementary-function approximations, and projection tolerances in advance. Each call is then an exactly specified map of a finite input. The ideal continuous proposal in the proof uses this very map as its drift. Consequently its Gaussian transition density is exact, even though the drift has been numerically approximated and may be discontinuous at lattice boundaries.

Rounding the field by at most $\delta$ changes the TAP objective uniformly by at most $N\delta$. On the matrix guard the fixed-depth AMP map is polynomially Lipschitz in that field. Take $D_0$ large enough that both errors are below the prescribed objective accuracy and the fixed anchor slack. Final drift rounding uses the feasible inward adjustment described below, preserving both the cube bounds and the stored-domain constraints on every correct guarded domain.

The elementary functions $\log$, $\exp$, $\tanh$, and square roots are evaluated by deterministic approximations with specified absolute errors. On their guarded domains these require polynomially many bits and operations. Large arguments of $\tanh$ can be truncated with their explicit exponentially small error. In log-density calculations, use log weights and compare $\log V$ with the log acceptance ratio. There is no need to form a potentially enormous unnormalized density.

\subsection{Optimization and feasible projections}

Recall the prescribed accuracy and clipping parameters:
\[
 \varepsilon_N=N^{-8},\qquad
 u_N=\tfrac14\min(\rho,a,1)N^{-30},\qquad
 \mathcal B_N=[-1+u_N,1-u_N]^N.
\]
For a ball domain, coordinate clipping fixes its center and is nonexpansive, so it preserves feasibility. For a cap, contraction toward $(1-\rho/2)\tau$ by the factor $2u_N/\rho$ produces a point of its intersection with $\mathcal B_N$. The binary entropy modulus and the explicit interaction terms therefore give
\[
 \sup_{z\in\mathcal K}| F_{\tap}(z',f)- F_{\tap}(z,f)|
 \le C\{N^4u_N+Nu_N\log(e/u_N)\}
 =o(\varepsilon_N),
\]
where $z'$ is the indicated feasible modification and the guards are in force. This comparison applies to the original minimizer as well as to every feasible point. It proves that solving the clipped problem is sufficient without clipping the true posterior mean in the analysis.

On a correct clipped domain, the Hessian is between $\gamma I_N$ and $L_NI_N$, where $L_N=u_N^{-1}+\|J\|_F+3\beta^2$. The projected-gradient estimate of \cite[Theorem~3.10]{bubeck2015convex} gives contraction in squared distance. The guarded gradient bound then converts distance accuracy into objective accuracy. Thus the prescribed $n_{\mathrm{opt}}$ iterations in Algorithm~\ref{alg:tap-optimization} suffice. On a correct convex domain, the exact projected-gradient map is nonexpansive; if each numerical update has error at most $e_N$, accumulated distance error is at most $n_{\mathrm{opt}}e_N$. Choosing
\[
 e_N\le\frac{\varepsilon_N}
 {100\,n_{\mathrm{opt}}(1+L_N+B_N)}
\]
with an additional fixed polynomial factor when needed absorbs gradient, projection, and rounding errors into the objective budget. Every quantity in this precision prescription is polynomially bounded. Off the structural event, the same iteration count is used and a bounded feasible point is returned.

The projection formulas are given in \S\ref{sec:algorithm}. For the ball, the multiplier can be bracketed using
\[
 \|z(\lambda)-c_j\|\le\frac{\|v-c_j\|}{1+\lambda}.
\]
The guarded iterates and gradients give a polynomial upper endpoint for this bracket. For the cap, the scalar sum of clipped coordinates is continuous and monotone in the multiplier, and a polynomial bracket follows from the coordinate bounds. Bisection requires $O(\log N)$ iterations to any fixed inverse-polynomial precision. Choose the feasible side of the bracket and make a further small contraction toward the fixed feasible center before rounding. This leaves both constraints satisfied, with arbitrarily small inverse-polynomial projection error. The center of the cap is a fixed positive distance from its cutting hyperplane, and each clipped ball center has coordinate slack $u_N$, so the final inward correction also has a polynomial precision requirement.

The objective and projection routines always stop after their prescribed budgets. Their fallbacks are part of the drift definition, and the construction on a correct guarded domain proves that those fallbacks are not needed there.

\subsection{Time grids, Gaussian draws, and final spin coupling}

There are polynomially many random draws and coordinate tests over all $O(N\log N)$ proposals. The admissible-offset construction in \eqref{sampler:grid} gives a minimum ideal time step $d_N=N^{-20}$ and an offset density at most $2/\eta$. Approximate the offset and each grid time much more accurately than $d_N$. When forming $U+j\eta$, evaluate $\eta$ with an additional factor $M^{-1}$ in its error tolerance. The removed neighborhoods ensure that the ordering of the grid is preserved and that numerical intervals have length at least $d_N/2$, except on an event of arbitrarily small inverse-polynomial probability. Square roots of these interval lengths then have a polynomial Lipschitz bound.

For any iteration count written as the ceiling of a real expression, first compute a certified rational upper approximation with absolute error at most $1/4$ and then take its ceiling. This increases the displayed count by at most one. In particular, this convention for $M$, $k_N$, and $n_{\mathrm{opt}}$ preserves every lower bound needed in the proof and all polynomial cost estimates.

Truncate every scalar standard Gaussian at a threshold $C\sqrt{\log N}$, taking the fixed constant $C$ large enough for the polynomial number of draws. The total truncation probability is an arbitrarily small inverse polynomial. Gaussian and uniform variables can then be approximated with a fixed polynomial bit budget. Couple these approximations with their ideal continuous variables. Until a field-rounding, sign, or guard decision differs, both versions use identical deterministic drift maps. Their coordinate errors consequently accumulate through polynomially many bounded additions and interval-length perturbations; no Lipschitz estimate across a field-rounding discontinuity is needed.

Conditionally on the past, each new ideal observation coordinate is Gaussian with variance at least $d_N$. Its density is at most $(2\pi d_N)^{-1/2}$. On the guarded region there are only polynomially many lattice boundaries, and there are polynomially many sign and guard tests. A union bound using this conditional density estimate shows that the probability of any discrepant decision is arbitrarily small when the arithmetic precision is a sufficiently fine inverse polynomial. The offset has the corresponding bound using its density $2/\eta$ and its finitely many excluded endpoints. Both path versions invoke the identical fixed finite program on an identical rounded field, so their drift outputs agree until one of these decisions differs. Approximation of an exact projection by that program was already absorbed into the optimization error budget.

On agreement of these decisions, the difference between the two log weights in \eqref{sampler:log-weight} is bounded by a polynomial multiple of the numerical coordinate error. Choose the score precision to make this difference $o((N\log N)^{-1})$ per proposal. The acceptance function $z\mapsto\min(1,e^z)$ is $1$-Lipschitz. Coupling the Metropolis steps with common uniforms therefore makes the total probability of a different acceptance decision $o(1)$. Every retained state is one of the generated proposals, so the preceding Gaussian anti-concentration bounds suffice; no density bound for the accepted-state law is required.

The terminal signs coincide whenever all coupled decisions agree. Hence the total-variation distance between the ideal and finite-precision \emph{spin outputs} is $o(1)$. A rounded continuous path and its unrounded version need not be close in total variation, and no such comparison is used. Adding the input-score perturbation bound and this discrete-output coupling completes the numerical part of Theorem~\ref{thm:main}.


\bibliographystyle{alpha}
\bibliography{SKsampling}
\end{document}
